\documentclass{article}
%DIF LATEXDIFF DIFFERENCE FILE
%DIF DEL /mnt/nfsdata/nfsdata/linshuijin/DIRECTPP/revision_workspace/integration/derived/baseline_flat.tex   Thu Nov  5 19:29:11 2026
%DIF ADD /mnt/nfsdata/nfsdata/linshuijin/DIRECTPP/revision_workspace/integration/derived/current_flat.tex    Thu Nov  5 19:29:11 2026

\usepackage[utf8]{inputenc}
\usepackage{authblk}
\usepackage{setspace}
\usepackage[margin=1.25in]{geometry}
\usepackage{graphicx}
\graphicspath{{./figures/}}
\usepackage{subcaption}
\usepackage{amsmath}
\usepackage{amssymb}
\usepackage{booktabs}
%DIF 13a13-17
\usepackage{longtable,pdflscape} %DIF > 
\usepackage[T1]{fontenc} %DIF > 
\usepackage{lmodern} %DIF > 
\input{glyphtounicode} %DIF > 
\pdfgentounicode=1 %DIF > 
%DIF -------
\usepackage{array}
\usepackage{tabularx}
\usepackage{lineno}
\usepackage[
  style=nejm,
  citestyle=numeric-comp,
  sorting=none
]{biblatex}
%DIF 21c26
%DIF < \addbibresource{references.bib}
%DIF -------
\addbibresource{../revision_workspace/integration/derived/references_highlighted.bib} %DIF > 
%DIF -------
\usepackage[hidelinks]{hyperref}
\hypersetup{
%DIF 24c29
%DIF <   pdftitle={Precision Synthesis of Multi-Tracer PET via VLM-Modulated Rectified Flow for MCI Stratification},
%DIF -------
  pdftitle={Precision Synthesis of Multi-Tracer PET via VLM-Modulated Rectified Flow for MCI Stage Stratification}, %DIF > 
%DIF -------
  pdfauthor={Shuijin Lin; Shaozhen Yan; Tuo Liu; Haifeng Wang; Chunfeng Lian; Jie Lu; Jianhua Ma; Lin Xie},
  pdfsubject={Research Article},
  pdfkeywords={Controllable Cross-Modal Synthesis; Domain-Knowledge Encoding; PET Imaging; Alzheimer's Disease; Vision-Language Model; Rectified Flow}
}

%DIF 30c35
%DIF < \title{Precision Synthesis of Multi-Tracer PET via VLM-Modulated Rectified Flow for MCI Stratification}
%DIF -------
\title{Precision Synthesis of Multi-Tracer PET via VLM-Modulated Rectified Flow for MCI Stage Stratification} %DIF > 
%DIF -------

\author[1,4,$\dagger$]{Shuijin Lin}
\author[2,$\dagger$]{Shaozhen Yan}
\author[3,4,$\dagger$]{Tuo Liu}
\author[3,4]{Haifeng Wang}
\author[3,4]{Chunfeng Lian}
\author[2,*]{Jie Lu}
\author[1,4,*]{Jianhua Ma}
\author[1,4,*]{Lin Xie}

\affil[1]{Key Laboratory of Biomedical Information Engineering of Ministry of Education, School of Life Science and Technology, Xi'an Jiaotong University, No.~28 Xianning West Road, Xi'an, Shaanxi 710049, China.}
\affil[2]{Department of Radiology and Nuclear Medicine, Xuanwu Hospital, Capital Medical University, No.~45 Changchun Street, Xicheng District, Beijing 100053, China.}
\affil[3]{School of Mathematics and Statistics, Xi'an Jiaotong University, No.~28 Xianning West Road, Xi'an, Shaanxi 710049, China.}
\affil[4]{Research Center for Intelligent Medical Equipment and Devices (IMED), Xi'an Jiaotong University, No.~28 Xianning West Road, Xi'an, Shaanxi 710049, China.}
\affil[*]{Address correspondence to: Jie Lu, \href{mailto:imaginglu@hotmail.com}{imaginglu@hotmail.com}; Jianhua Ma, \href{mailto:jhma@xjtu.edu.cn}{jhma@xjtu.edu.cn}; Lin Xie, \href{mailto:linxie@xjtu.edu.cn}{linxie@xjtu.edu.cn}.}
\affil[$\dagger$]{These authors contributed equally to this work.}

\date{}

%DIF 50a55
\emergencystretch=2em %DIF > 
%DIF -------
\onehalfspacing
\linenumbers
%DIF PREAMBLE EXTENSION ADDED BY LATEXDIFF
%DIF CFONT PREAMBLE %DIF PREAMBLE
\RequirePackage{color}\definecolor{RED}{rgb}{1,0,0}\definecolor{BLUE}{rgb}{0,0,1} %DIF PREAMBLE
\DeclareOldFontCommand{\sf}{\normalfont\sffamily}{\mathsf} %DIF PREAMBLE
\providecommand{\DIFaddtex}[1]{{\protect\color{blue} \sf #1}} %DIF PREAMBLE
\providecommand{\DIFdeltex}[1]{{\protect\color{red} \ifmmode #1\else\scriptsize #1\fi}} %DIF PREAMBLE
%DIF SAFE PREAMBLE %DIF PREAMBLE
\providecommand{\DIFaddbegin}{} %DIF PREAMBLE
\providecommand{\DIFaddend}{} %DIF PREAMBLE
\providecommand{\DIFdelbegin}{} %DIF PREAMBLE
\providecommand{\DIFdelend}{} %DIF PREAMBLE
\providecommand{\DIFmodbegin}{} %DIF PREAMBLE
\providecommand{\DIFmodend}{} %DIF PREAMBLE
%DIF FLOATSAFE PREAMBLE %DIF PREAMBLE
\providecommand{\DIFaddFL}[1]{\DIFadd{#1}} %DIF PREAMBLE
\providecommand{\DIFdelFL}[1]{\DIFdel{#1}} %DIF PREAMBLE
\providecommand{\DIFaddbeginFL}{} %DIF PREAMBLE
\providecommand{\DIFaddendFL}{} %DIF PREAMBLE
\providecommand{\DIFdelbeginFL}{} %DIF PREAMBLE
\providecommand{\DIFdelendFL}{} %DIF PREAMBLE
%DIF HYPERREF PREAMBLE %DIF PREAMBLE
\providecommand{\DIFadd}[1]{\texorpdfstring{\DIFaddtex{#1}}{#1}} %DIF PREAMBLE
\providecommand{\DIFdel}[1]{\texorpdfstring{\DIFdeltex{#1}}{}} %DIF PREAMBLE
\newcommand{\DIFscaledelfig}{0.5}
%DIF HIGHLIGHTGRAPHICS PREAMBLE %DIF PREAMBLE
\RequirePackage{settobox} %DIF PREAMBLE
\RequirePackage{letltxmacro} %DIF PREAMBLE
\newsavebox{\DIFdelgraphicsbox} %DIF PREAMBLE
\newlength{\DIFdelgraphicswidth} %DIF PREAMBLE
\newlength{\DIFdelgraphicsheight} %DIF PREAMBLE
% store original definition of \includegraphics %DIF PREAMBLE
\LetLtxMacro{\DIFOincludegraphics}{\includegraphics} %DIF PREAMBLE
\newcommand{\DIFaddincludegraphics}[2][]{{\color{blue}\fbox{\DIFOincludegraphics[#1]{#2}}}} %DIF PREAMBLE
\newcommand{\DIFdelincludegraphics}[2][]{% %DIF PREAMBLE
\sbox{\DIFdelgraphicsbox}{\DIFOincludegraphics[#1]{#2}}% %DIF PREAMBLE
\settoboxwidth{\DIFdelgraphicswidth}{\DIFdelgraphicsbox} %DIF PREAMBLE
\settoboxtotalheight{\DIFdelgraphicsheight}{\DIFdelgraphicsbox} %DIF PREAMBLE
\scalebox{\DIFscaledelfig}{% %DIF PREAMBLE
\parbox[b]{\DIFdelgraphicswidth}{\usebox{\DIFdelgraphicsbox}\\[-\baselineskip] \rule{\DIFdelgraphicswidth}{0em}}\llap{\resizebox{\DIFdelgraphicswidth}{\DIFdelgraphicsheight}{% %DIF PREAMBLE
\setlength{\unitlength}{\DIFdelgraphicswidth}% %DIF PREAMBLE
\begin{picture}(1,1)% %DIF PREAMBLE
\thicklines\linethickness{2pt} %DIF PREAMBLE
{\color[rgb]{1,0,0}\put(0,0){\framebox(1,1){}}}% %DIF PREAMBLE
{\color[rgb]{1,0,0}\put(0,0){\line( 1,1){1}}}% %DIF PREAMBLE
{\color[rgb]{1,0,0}\put(0,1){\line(1,-1){1}}}% %DIF PREAMBLE
\end{picture}% %DIF PREAMBLE
}\hspace*{3pt}}} %DIF PREAMBLE
} %DIF PREAMBLE
\LetLtxMacro{\DIFOaddbegin}{\DIFaddbegin} %DIF PREAMBLE
\LetLtxMacro{\DIFOaddend}{\DIFaddend} %DIF PREAMBLE
\LetLtxMacro{\DIFOdelbegin}{\DIFdelbegin} %DIF PREAMBLE
\LetLtxMacro{\DIFOdelend}{\DIFdelend} %DIF PREAMBLE
\DeclareRobustCommand{\DIFaddbegin}{\DIFOaddbegin \let\includegraphics\DIFaddincludegraphics} %DIF PREAMBLE
\DeclareRobustCommand{\DIFaddend}{\DIFOaddend \let\includegraphics\DIFOincludegraphics} %DIF PREAMBLE
\DeclareRobustCommand{\DIFdelbegin}{\DIFOdelbegin \let\includegraphics\DIFdelincludegraphics} %DIF PREAMBLE
\DeclareRobustCommand{\DIFdelend}{\DIFOaddend \let\includegraphics\DIFOincludegraphics} %DIF PREAMBLE
\LetLtxMacro{\DIFOaddbeginFL}{\DIFaddbeginFL} %DIF PREAMBLE
\LetLtxMacro{\DIFOaddendFL}{\DIFaddendFL} %DIF PREAMBLE
\LetLtxMacro{\DIFOdelbeginFL}{\DIFdelbeginFL} %DIF PREAMBLE
\LetLtxMacro{\DIFOdelendFL}{\DIFdelendFL} %DIF PREAMBLE
\DeclareRobustCommand{\DIFaddbeginFL}{\DIFOaddbeginFL \let\includegraphics\DIFaddincludegraphics} %DIF PREAMBLE
\DeclareRobustCommand{\DIFaddendFL}{\DIFOaddendFL \let\includegraphics\DIFOincludegraphics} %DIF PREAMBLE
\DeclareRobustCommand{\DIFdelbeginFL}{\DIFOdelbeginFL \let\includegraphics\DIFdelincludegraphics} %DIF PREAMBLE
\DeclareRobustCommand{\DIFdelendFL}{\DIFOaddendFL \let\includegraphics\DIFOincludegraphics} %DIF PREAMBLE
%DIF LISTINGS PREAMBLE %DIF PREAMBLE
\RequirePackage{listings} %DIF PREAMBLE
\RequirePackage{color} %DIF PREAMBLE
\lstdefinelanguage{DIFcode}{ %DIF PREAMBLE
%DIF DIFCODE_CFONT %DIF PREAMBLE
  moredelim=[il][\color{red}\scriptsize]{\%DIF\ <\ }, %DIF PREAMBLE
  moredelim=[il][\color{blue}\sffamily]{\%DIF\ >\ } %DIF PREAMBLE
} %DIF PREAMBLE
\lstdefinestyle{DIFverbatimstyle}{ %DIF PREAMBLE
	language=DIFcode, %DIF PREAMBLE
	basicstyle=\ttfamily, %DIF PREAMBLE
	columns=fullflexible, %DIF PREAMBLE
	keepspaces=true %DIF PREAMBLE
} %DIF PREAMBLE
\lstnewenvironment{DIFverbatim}{\lstset{style=DIFverbatimstyle}}{} %DIF PREAMBLE
\lstnewenvironment{DIFverbatim*}{\lstset{style=DIFverbatimstyle,showspaces=true}}{} %DIF PREAMBLE
%DIF END PREAMBLE EXTENSION ADDED BY LATEXDIFF

\begin{document}

\maketitle

\begin{abstract}
\noindent\textbf{Objective:} \DIFdelbegin \DIFdel{Here, we introduce }\DIFdelend \DIFaddbegin \DIFadd{To evaluate }\DIFaddend DIReCT$++$, a \DIFdelbegin \DIFdel{Domain-Informed ReCTified flow model for synthesizing multi-tracer PET from MRI combined with fundamental clinical information}\DIFdelend \DIFaddbegin \DIFadd{domain-informed rectified flow framework for synthesizing FDG- and AV45-PET from MRI and clinical text inputs}\DIFaddend .

\noindent\textbf{Impact Statement:} This work \DIFdelbegin \DIFdel{advances a scalable, data-efficient tool for the early diagnosis and prognostic prediction of AD}\DIFdelend \DIFaddbegin \DIFadd{evaluates multi-tracer PET synthesis and cross-sectional MCI stage discrimination using a shared tracer-conditioned framework}\DIFaddend .

\noindent\textbf{Introduction:} \DIFdelbegin \DIFdel{The biological definition of Alzheimer's disease (AD) relies on multi-modal neuroimaging, yet the clinical utility of positron emission tomography (PET ) is limited by costand radiation exposure , hindering early screening at preclinical or prodromal stages. While generative models offer a promising alternative by synthesizing PET from magnetic resonance imaging (MRI), achieving subject-specific precision remains a primary challenge}\DIFdelend \DIFaddbegin \DIFadd{PET cost, radiation exposure and limited accessibility motivate MRI-based synthesis, but image similarity alone does not establish individual molecular validity}\DIFaddend .

\noindent\textbf{Methods:} \DIFdelbegin \DIFdel{Our approach integrates a }\DIFdelend \DIFaddbegin \DIFadd{A shared }\DIFaddend 3D rectified flow \DIFdelbegin \DIFdel{architecture to capture complex cross-modal and cross-tracer relationships with a domain-adapted vision-language model (BiomedCLIP) that provides text-guided, personalized generationusing clinical scores and imaging knowledge}\DIFdelend \DIFaddbegin \DIFadd{generator uses frozen BiomedCLIP-derived patient and tracer features for learned one-step inference. Separate analyses evaluate affine tracer adaptation, two-fold OOF generation, normalized regional uptake agreement and three-fold cross-sectional classification. ADNI includes 1,857 participants; historical OASIS-3 external evaluation is restricted to AV45 in CN/AD}\DIFaddend .

\noindent\textbf{Results:} \DIFdelbegin \DIFdel{Extensive evaluations on multi-center datasets demonstrate that DIReCT$++$ not only produces synthetic PET images ($^{18}$F-AV-45 and $^{18}$F-FDG) of superior fidelity and generalizability but also accurately recapitulates disease-specific patterns}\DIFdelend \DIFaddbegin \DIFadd{Masked OOF comparisons favored Full over a clinical regression U-Net and MRI-only Swin-UNETR in all 16 predefined comparisons within their respective Holm families, although some effects were small. Other conditioning comparisons showed metric-dependent tradeoffs. In AD/CN classification, MRI+Syn-both AUC was $93.82\pm3.21\%$ versus MRI-only $87.02\pm5.44\%$; MCI/CN did not improve, and EMCI/LMCI showed mixed endpoints. Clinical-only and MRI+Clinical AUC exceeded MRI+Syn-both. Historical regional intensities showed close agreement, without establishing clinical SUV/SUVR calibration or equivalence}\DIFaddend .

\noindent\textbf{Conclusion:} \DIFdelbegin \DIFdel{Crucially, combining these synthesized PET images with MRI enables precise personalized stratification of mild cognitive impairment (MCI), advancing a scalable, data-efficient tool for the early diagnosis and prognostic prediction of AD}\DIFdelend \DIFaddbegin \DIFadd{DIReCT$++$ supports further investigation of multi-tracer synthesis and cross-sectional classification. Extensive diagnosis-conditioned filling of clinical scores, source-identity limitations and mixed task effects preclude claims of independent diagnostic value, longitudinal prognosis or replacement of molecular confirmation}\DIFaddend .
\end{abstract}

\noindent\textbf{Keywords:} Controllable Cross-Modal Synthesis; Domain-Knowledge Encoding; PET Imaging; Alzheimer's Disease; Vision-Language Model; Rectified Flow


\section{Introduction}
Alzheimer's disease (AD), the most prevalent neurodegenerative disorder, constitutes a mounting global health crisis. Its pathological processes begin years before clinical dementia manifests, making the prodromal stage of mild cognitive impairment (MCI) a critical window for intervention~\cite{scheltens2021alzheimer}. However, the heterogeneous nature of MCI, where only a subset of individuals progress to AD, poses a significant prognostic challenge~\cite{haller2023neuroimaging}. \DIFdelbegin \DIFdel{This challenge is underscored by the modern biological definition of AD , which is based on in vivo evidence of core pathologies, specifically }\DIFdelend \DIFaddbegin \DIFadd{AD evaluation uses complementary molecular, functional and structural measurements: }\DIFaddend amyloid-$\beta$ (A$\beta$) deposition \DIFdelbegin \DIFdel{(e.g., on }\DIFdelend \DIFaddbegin \DIFadd{can be assessed with }\DIFaddend $^{18}$F-AV45 PET\DIFdelbegin \DIFdel{) and neuronal injury (e.g., }\DIFdelend \DIFaddbegin \DIFadd{, glucose-metabolic dysfunction with }\DIFaddend $^{18}$F-FDG PET\DIFdelbegin \DIFdel{hypometabolism), }\DIFdelend \DIFaddbegin \DIFadd{, and associated atrophy }\DIFaddend with structural MRI\DIFdelbegin \DIFdel{quantifying associated atrophy}\DIFdelend ~\cite{ossenkoppele2022amyloid,yan2021sex,haller2023neuroimaging}. \DIFaddbegin \DIFadd{These modalities are not interchangeable; in particular, FDG hypometabolism and MRI atrophy do not by themselves establish the biological diagnosis of AD~\cite{jack2024revised}. }\DIFaddend In clinical practice, however, the routine acquisition of PET is hampered by high cost, limited access, and radiation exposure, \DIFdelbegin \DIFdel{restricting this multi-modal biomarker profiling to research settings and hindering timely, widespread screening}\DIFdelend \DIFaddbegin \DIFadd{limiting access to complementary imaging in research and selected diagnostic workups}\DIFaddend ~\cite{nordberg2010use}.

Generative artificial intelligence (AI) offers \DIFdelbegin \DIFdel{a transformative pathway to bridge this clinical gap by synthesizing PET biomarkers }\DIFdelend \DIFaddbegin \DIFadd{methods for estimating PET images }\DIFaddend from routinely acquired MRI~\cite{dayarathna2024deep,lee2024synthesizing}. Initial studies, primarily using generative adversarial networks (GANs)~\cite{pan2018synthesizing,pan2020spatially,apostolopoulos2022applications} and, more recently, diffusion models~\cite{xie2024synthesizing,chen2025multi,theodorou2025mri2pet}, have established the feasibility of producing visually realistic PET images. Subsequent advances have sought to improve fidelity by incorporating domain knowledge, for instance, through guidance from large vision-language models (VLMs)~\cite{wang2025unisyn,liu2025direct}.  \DIFdelbegin \DIFdel{Despite these advances, a critical limitation remains: the }\DIFdelend \DIFaddbegin \DIFadd{However, visual realism alone does not establish }\DIFaddend subject-specific \DIFdelbegin \DIFdel{precision necessary for personalized diagnosis and prognostic stratification is still lacking, which is particularly crucial for }\DIFdelend \DIFaddbegin \DIFadd{accuracy for cross-sectional diagnosis or MCI stage discrimination, particularly in }\DIFaddend a heterogeneous condition \DIFdelbegin \DIFdel{like }\DIFdelend \DIFaddbegin \DIFadd{such as }\DIFaddend AD~\cite{ossenkoppele2022amyloid}. This \DIFdelbegin \DIFdel{challenges }\DIFdelend \DIFaddbegin \DIFadd{challenge }\DIFaddend stems from the inherent nature of MRI-to-PET synthesis as an ill-posed inverse problem; \DIFdelbegin \DIFdel{neurodegenration }\DIFdelend \DIFaddbegin \DIFadd{neurodegeneration }\DIFaddend visible on MRI is typically a downstream effect of prior molecular pathologies captured by PET~\cite{jagust2018imaging}. Consequently, individuals with similar structural MRI presentations at early disease stages can exhibit vastly different pathological burdens, leading to ambiguous mapping. \DIFdelbegin \DIFdel{Overcoming this ambiguity requires an AI synthesis framework that can comprehensively model complex cross-modal and cross-tracer associations to provide strong, }\DIFdelend \DIFaddbegin \DIFadd{We therefore investigate }\DIFaddend patient-specific \DIFdelbegin \DIFdel{guidance}\DIFdelend \DIFaddbegin \DIFadd{clinical and tracer information as complementary conditioning for MRI-to-PET synthesis}\DIFaddend .

\DIFdelbegin \DIFdel{To this end, we propose }\DIFdelend \DIFaddbegin \DIFadd{Building upon single-tracer DIReCT~\cite{liu2025direct}, }\DIFaddend DIReCT$++$ (Domain-Informed ReCTified flow) \DIFdelbegin \DIFdel{, a VLM-modulated generative model for the high-fidelity synthesis of patient-specific, multi-tracer PET images from baseline MRI and textual guidance. The DIReCT$++$ framework is built on two core innovations. First, it employs an efficient }\DIFdelend \DIFaddbegin \DIFadd{extends }\DIFaddend 3D rectified flow \DIFdelbegin \DIFdel{architecture to directly capture the complex, high-dimensional relationships between MRI and multi-tracer PET (specifically $^{18}$F-AV-45 and $^{18}$F-FDG), enabling concurrent and rapid generation . Second, to achieve subject-specific precision, the model is conditioned via a VLM (}\DIFdelend \DIFaddbegin \DIFadd{to a shared network conditioned on the target PET tracer, enabling FDG- or AV45-PET generation in separate forward passes. A second extension is lightweight affine adaptation of }\DIFaddend BiomedCLIP~\cite{zhang2025multimodal}\DIFdelbegin \DIFdel{) with application-dedicated adaptations. This domain-adapted VLM integrates fundamental patient knowledge (encoded as }\DIFdelend \DIFaddbegin \DIFadd{-derived tracer features using text--PET alignment and a cross-tracer soft similarity target. We distinguish this adaptation formulation from reconstructed generator configurations using raw frozen tracer vectors. The framework combines tracer knowledge with }\DIFaddend patient-specific clinical \DIFdelbegin \DIFdel{assessments and tracer-specific imaging domain knowledge) to guide the synthesis process, ensuring the generated PET images reflect individual pathological states rather than population averages. By reconciling an efficient generative process with precise, text-based patient conditioning, our approach directly addresses the ill-posed nature of cross-modal synthesis in heterogeneous disease populations}\DIFdelend \DIFaddbegin \DIFadd{text and learned one-step inference. Its components are evaluated with configuration-specific comparisons rather than assumed to resolve the biological ambiguity of MRI-to-PET translation}\DIFaddend .

We \DIFdelbegin \DIFdel{rigorously evaluate DIReCT$++$ on large-scale, multi-center datasets, including the Alzheimer's Disease Neuroimaging Initiative (ADNI)}\DIFdelend \DIFaddbegin \DIFadd{evaluate image synthesis on ADNI}\DIFaddend ~\cite{weiner2010alzheimer} and \DIFdelbegin \DIFdel{Open Access Series of Imaging Studies (OASIS)}\DIFdelend \DIFaddbegin \DIFadd{retain historical AV45 external results from OASIS-3}\DIFaddend ~\cite{lamontagne2019oasis}. \DIFdelbegin \DIFdel{Comprehensive experiments demonstrate that our model generates multi-tracer PET images with superior fidelity and generalizability compared to state-of-the-art methods. Crucially, the synthesized images accurately recapitulate known disease-specific patternsof amyloid deposition and hypometabolism. When combined with the original MRI scans to train deep learning classifiers~\cite{lian2018hierarchical,huang2017densely}, the synthesized PET data enable accurate personalized diagnosis and , most importantly, the stratification of MCI patients based on their risk of progression to AD. Our work establishes a translational paradigm for precision biomarker synthesis, offering a scalable and data-efficient tool for early intervention in dementia}\DIFdelend \DIFaddbegin \DIFadd{Regional analyses examine normalized uptake agreement and AD--CN effect patterns. Separate DenseNet-based classification~\cite{huang2017densely} evaluates AD/CN, MCI/CN and EMCI/LMCI under fixed three-fold splits, alongside Clinical-only and MRI+Clinical baselines. The study addresses cross-sectional performance and methodological tradeoffs; it does not test future AD conversion, establish clinical equivalence or validate synthetic PET as molecular confirmation}\DIFaddend .


\section{Results}
\subsection{The \texorpdfstring{DIReCT$++$}{DIReCT++} Framework for Precision PET Synthesis}
DIReCT$++$ is a VLM-modulated rectified flow framework that synthesizes \DIFdelbegin \DIFdel{high-fidelity $^{18}$F-AV-45 and $^{18}$F-FDG PET images from a single routine MRI scan, guided by fundamental text prompts }\DIFdelend \DIFaddbegin \DIFadd{FDG- or AV45-PET estimates from MRI with patient and tracer text guidance }\DIFaddend (Fig.~\ref{fig:1}A). \DIFdelbegin \DIFdel{The core of our method is a }\DIFdelend \DIFaddbegin \DIFadd{Its }\DIFaddend conditional 3D \DIFdelbegin \DIFdel{rectified flow (RF) model, implemented as a multi-task }\DIFdelend U-Net \DIFdelbegin \DIFdel{that predicts velocity fields }\DIFdelend \DIFaddbegin \DIFadd{predicts a tracer-specific velocity field for learned one-step inference }\DIFaddend (Fig.~\ref{fig:1}C). \DIFdelbegin \DIFdel{This architecture establishes a direct mapping from MRI to multi-tracer PET, enabling high-fidelity image generation in a single sampling step}\DIFdelend \DIFaddbegin \DIFadd{The shared network handles both tracer targets; this does not imply simultaneous dual-tracer output or clinical molecular validity}\DIFaddend .

\DIFdelbegin \DIFdel{To achieve subject-wise precision, text guidance incorporating both common knowledge (e.g., general descriptions of PET tracers) and personal }\DIFdelend \DIFaddbegin \DIFadd{Text guidance incorporates tracer descriptions and }\DIFaddend patient information (\DIFdelbegin \DIFdel{e.g., age, gender, }\DIFdelend \DIFaddbegin \DIFadd{age, sex, weight and }\DIFaddend clinical scores) \DIFdelbegin \DIFdel{is encoded by a pre-trained VLM and integrated }\DIFdelend \DIFaddbegin \DIFadd{encoded by frozen BiomedCLIP encoders and injected }\DIFaddend into the RF model \DIFdelbegin \DIFdel{via }\DIFdelend \DIFaddbegin \DIFadd{through }\DIFaddend cross-attention. \DIFdelbegin \DIFdel{We fine-tune the VLM from BiomedCLIP using a straightforward prompt learning strategy }\DIFdelend \DIFaddbegin \DIFadd{The affine fine-tuning formulation optimizes lightweight adapters of BiomedCLIP-derived tracer guidance, rather than BiomedCLIP encoder weights, using a soft cross-tracer similarity penalty }\DIFaddend (Fig.~\ref{fig:1}B). \DIFdelbegin \DIFdel{This strategy maximizes the intra-tracer alignment between text and image embeddings while preserving the inherent inter-tracer differences, reflecting the distinct yet complementary molecular activities captured by $^{18}$F-AV-45 and $^{18}$F-FDG PET }\DIFdelend \DIFaddbegin \DIFadd{In a separate adapter-level analysis, varying $\tau$ from 0.70 to 0.98 changed the mean cross-tracer embedding similarity from 0.7032 to 0.9807. At $\tau=0.82$, the similarity was 0.8214, compared with 0.9897 without the cross-tracer penalty (Table~\ref{tab:supp_tau_adapter}). These are representation-level results, not a comparison of PET generation quality across $\tau$. The reconstructed Full generation experiments used the original fixed tracer embeddings; their performance should not be attributed to affine adaptation}\DIFaddend .

\DIFdelbegin \DIFdel{After training, DIReCT$++$ performs one-step inference: given an input MRIand corresponding text, it simultaneously generates the corresponding $^{18}$F-AV-45 and $^{18}$F-FDG PET images }\DIFdelend \DIFaddbegin \DIFadd{A separate E1 fixed-weight sensitivity analysis replaced the clinical-score portion of the patient prompt while retaining age, sex and weight, MRI, tracer and model parameters (Table~\ref{tab:e1_score_removal}). Removing scores increased ROI MAE by approximately 0.000036 for FDG and 0.000079 for AV45, with statistical support under the endpoint-specific two-tracer Holm families but very small absolute changes. SSIM increased while PSNR decreased, and the average absolute prediction change was approximately 0.001019/0.001105 for FDG/AV45. These observations show conditional output sensitivity; they are not a retrained No-Scores ablation or proof of independent diagnostic benefit.
}

\DIFadd{At inference, the shared network predicts a velocity field for the selected tracer from MRI and the corresponding two-token text context. A single forward pass produces one tracer-specific PET estimate; FDG and AV45 are obtained in separate tracer-conditioned passes }\DIFaddend (Fig.~\ref{fig:1}D). The resulting \DIFdelbegin \DIFdel{synthetic images maintain both subject-wise }\DIFdelend \DIFaddbegin \DIFadd{images are evaluated through image-quality measures, regional uptake analysis }\DIFaddend and \DIFdelbegin \DIFdel{tracer-wise realism, making them directly applicable to downstream analytical tasks such as statistical analysis , early diagnosis, and stratification }\DIFdelend \DIFaddbegin \DIFadd{cross-sectional classification }\DIFaddend (Fig.~\ref{fig:1}E).

\begin{figure}[htbp]
    \centering
    \includegraphics[width=\textwidth]{figure_1}
    \DIFdelbeginFL %DIFDELCMD < \caption{\textbf{Overview of the DIReCT$++$ framework for multimodal PET synthesis and downstream analysis.}
%DIFDELCMD < (A) Schematic of the rectified flow (RF) process that synthesizes dual-tracer $^{18}$F-AV-45 and $^{18}$F-FDG PET images from MRI. (B) Fine-tuning of the pre-trained BiomedCLIP model via prompt learning to encode tracer-specific and subject-level text guidance. (C) Training pipeline of the conditional 3D RF model implemented as a multi-task U-Net that predicts velocity fields for both tracers. (D) One-step inference process that simultaneously generates synthetic AV45- and FDG-PET from MRI and textual input. (E) Downstream analytical tasks, including statistical analysis, early diagnosis, and stratification, demonstrating the clinical applicability of the synthesized PET images.}
%DIFDELCMD <     %%%
\DIFdelendFL \DIFaddbeginFL \caption{{\color{blue}\textbf{Overview of the DIReCT$++$ framework for multimodal PET synthesis and downstream analysis.}
(A) Tracer-conditioned rectified flow from MRI to FDG- or AV45-PET. (B) Lightweight affine fine-tuning of tracer guidance derived from frozen BiomedCLIP encoders. The symbols $\alpha$ and $\beta$ denote FDG adapter parameters $(\alpha_1,\alpha_2)$ and AV45 adapter parameters $(\beta_1,\beta_2)$, respectively; these adaptations are separate from the raw tracer embeddings used by the reconstructed Full generation experiments. (C) A shared conditional 3D U-Net predicts velocity fields for the requested tracer. (D) Learned one-step PET estimation by residual addition of the predicted velocity to MRI; each forward pass produces one tracer. (E) Image-quality, regional uptake and cross-sectional classification analyses.}}
    \DIFaddendFL \label{fig:1}
\end{figure}

\DIFdelbegin \par\noindent\DIFdel{\textbf{Superior Fidelity and Generalizability of Synthesized PET}}\par
%DIFDELCMD < %%%
\DIFdel{We rigorously evaluated the multi-tracer PET synthesis performance of DIReCT$++$ against a comprehensive suite of representative methods, including generative adversarial networks (CycleGAN~\cite{zhu2017unpaired}), transformer-based architectures (Swin-UNet~\cite{Cao2021swin}), and diffusion/flow-based models (SegGuidedDif ~\cite{konz2024anatomically}, 3D DDIM~\cite{Graf2023}, 3D RF ~\cite{Xing2022}), alongside our baseline model, DIReCT~\cite{liu2025direct}, which performs mono-tracer generation without prompt adaptation. Image quality was quantified }\DIFdelend \DIFaddbegin \subsection{\texorpdfstring{\DIFadd{PET Synthesis Quality and External Evaluation}}{PET Synthesis Quality and External Evaluation}}
\DIFadd{PET synthesis was evaluated }\DIFaddend using structural similarity (SSIM), peak signal-to-noise ratio (PSNR), mean squared error (MSE) \DIFdelbegin \DIFdel{, }\DIFdelend and mean absolute error (MAE)\DIFdelbegin \DIFdel{against real PET ground truth. Models were trained on the ADNI dataset and evaluated under two settings: internally via 3-fold cross-validation on ADNI, and externally on the independent OASIS cohort (for which only$^{18}$F-AV-45 was available)}\DIFdelend \DIFaddbegin \DIFadd{. Figure~\ref{fig:2} retains the historical comparisons with CycleGAN~\cite{zhu2017unpaired}, Swin-UNet~\cite{Cao2021swin}, segmentation-guided diffusion~\cite{konz2024anatomically}, 3D DDIM~\cite{song2021ddim,pinaya2023generative}, 3D rectified flow~\cite{Xing2022} and single-tracer DIReCT~\cite{liu2025direct}. Absolute historical aggregate values are given in Table~\ref{tab:historical_generation}; subject-level statistics were not recoverable for every historical method and are not reconstructed from these aggregates.
PET synthesis used two-fold ADNI evaluation; historical external evaluation used OASIS-3 AV45 only. The reconstructed OOF comparisons and their statistical protocols are reported separately from the historical Figure~\ref{fig:2} summaries}\DIFaddend .

\DIFdelbegin \DIFdel{On internal evaluation , DIReCT$++$ demonstrated consistent superiority, outperforming all competing methods by substantial margins for both $^{18}$F-AV-45 }\DIFdelend \DIFaddbegin \DIFadd{The reconstructed OOF evaluation does not support a uniform advantage of Full over all alternatives (Tables~\ref{tab:whole_volume_generation},\ref{tab:paired_generation}). Full outperformed the clinically conditioned regression U-Net across the four whole-volume metrics for both tracers, while the magnitude of the AV45 MSE effect was small ($|d_z|=0.075$). Relative to TabularMLP, Full had higher SSIM but lower PSNR and higher MSE for both tracers; MAE did not differ after the specified Holm correction. TracerOnly had better FDG results on all four metrics and higher AV45 SSIM than Full. Single-tracer DIReCT had higher FDG SSIM, whereas Full had better FDG PSNR/MSE and AV45 SSIM. These independently trained models differ in initialization, conditioning or implementation, so the comparisons do not isolate a single causal component or prove multi-tracer synergy.
}

\DIFadd{In the masked evaluation (Tables~\ref{tab:r3_11_masked_summary},\ref{tab:r3_11_masked_paired}), Full had higher mean SSIM/PSNR and lower mean MSE/MAE than Clinical U-Net }\DIFaddend and \DIFdelbegin \DIFdel{$^{18}$F-FDG synthesis (Fig.~\ref{fig:2}A). Quantitative gains were significant (e.g. }\DIFdelend \DIFaddbegin \DIFadd{MRI-only Swin-UNETR for both tracers; all 16 predefined comparisons met $p_{\mathrm{Holm}}<0.05$ within their two separate eight-test families. The masked cohort contained 1}\DIFaddend ,\DIFdelbegin \DIFdel{$\Delta$ $+2.21$ to $+11.46$ dB in PSNR for $^{18}$F-FDG; $\Delta$ $+0.98$  to $+12.66$ dB for $^{18}$F-AV-45), with similar robust improvements observed across all metrics. This quantitative superiority was corroborated by qualitative assessment (Fig.~\ref{fig:2}B). Syntheses from DIReCT $++$ realistically captured characteristic patterns of metabolic activity and amyloid deposition across the disease spectrum,exhibiting markedly higher fidelity than competitor outputs. }\DIFdelend \DIFaddbegin \DIFadd{597 FDG and 1,364 AV45 observations after the existing MRI/mask QC exclusion. For Clinical U-Net AV45 MSE/MAE, the standardized effects were small ($|d_z|=0.071/0.076$), and only 56.4--57.8\% of subjects favored Full. Thus statistical support does not imply uniform individual or clinical benefit. Full/Clinical U-Net used epoch 165, while MRI-only Swin-UNETR used epoch 200; information and training endpoints differ.
}\DIFaddend 

\DIFdelbegin \DIFdel{External validation on the OASIS cohort confirmed that the advantages of DIReCT$++$ generalize beyond the training data (Fig.~\ref{fig:2}A).
Our model sustained its leading performance, achieving a mean PSNR of $26.79$ dB(compared to $28.09$ dB on ADNI ) and outperforming competitors by $1.34$ to $14.41$ dB. Visual results on OASIS cases further demonstrated excellent generalization, with DIReCT$++$ reproducing realistic tracer uptake patterns where other methods showed deviations (}\DIFdelend \DIFaddbegin \DIFadd{Separate whole-volume and masked summaries are retained because background exclusion changes the metric definitions and values. Supplemental FDG and AV45 examples use the median, 10th-percentile and worst QC-passed subjects selected by saved whole-volume SSIM, and show MRI, real PET, synthetic PET and absolute-error maps (Figs.~\ref{fig:cases_fdg},\ref{fig:cases_av45}).
}

\DIFadd{The historical OASIS-3 evaluation reports AV45 PSNR of 26.79 dB, compared with 28.09 dB for historical internal ADNI AV45 evaluation (Fig.~\ref{fig:2}A; Table~\ref{tab:historical_generation}). The cross-dataset examples in }\DIFaddend Fig.~\ref{fig:2}C \DIFdelbegin \DIFdel{). This robust performance across diverse datasets affirms the state-of-the-art status of DIReCT$++$, highlighting its superior fidelity, generalization capacity, and effective cross-modal and cross-tracer integration}\DIFdelend \DIFaddbegin \DIFadd{illustrate CN/AD AV45 synthesis. These aggregate and illustrative results do not provide subject-level confidence intervals, FDG external evidence or external MCI-stage validation, and should not be combined with the reconstructed OOF statistics as a single experiment}\DIFaddend .

\begin{figure}[htbp]
\centering
\DIFdelbeginFL %DIFDELCMD < \includegraphics[width=0.65\textwidth]{figure_2}
%DIFDELCMD <     \caption{\textbf{Comparison of PET reconstruction quality across methods and datasets.}
%DIFDELCMD < (A) Radar plots show normalized quantitative results (SSIM, MSE, MAE, PSNR) for FDG-PET and AV45-PET reconstruction on ADNI, and cross-dataset generalization on OASIS. Notably, each metric is normalized by the best value among all methods. 
%DIFDELCMD < (B) Representative PET synthesis results for cognitively normal (CN, top), mild cognitive impairment (MCI, middle), and Alzheimer’s disease (AD, bottom) subjects. Columns: MRI input, competing methods, DIReCT (i.e., ours single), ours (DIReCT$++$), and real PET. Each sub-row depicts FDG (Up) and AV45 (Down) reconstruction. 
%DIFDELCMD < (C) Generalization to OASIS using models trained on ADNI. Rows: CN and AD subjects. Our method yields consistently higher quantitative scores and superior visual fidelity across datasets.}
%DIFDELCMD <     %%%
\DIFdelendFL \DIFaddbeginFL \includegraphics[width=\textwidth,height=.62\textheight,keepaspectratio]{figure_2}
\caption{{\color{blue}\textbf{Historical PET synthesis comparisons and illustrative examples.} (A) Radar plots summarize the historical aggregate SSIM, MSE, MAE and PSNR values for ADNI FDG/AV45 and OASIS-3 AV45. Metric values are displayed relative to the best value among methods; absolute aggregate values and scaling conventions are listed in Table~\ref{tab:historical_generation}. These aggregates are separate from the reconstructed OOF and masked analyses. (B) Illustrative ADNI CN, MCI and AD examples with MRI, competing synthesis outputs and real PET; paired subrows show FDG and AV45. (C) Historical CN/AD OASIS-3 examples are limited to AV45. Examples were randomly selected within diagnostic groups as reported by the authors; the selection seed could not be verified from the available records. This differs from the supplemental quantile-based selection and does not establish population-level superiority.}}
\DIFaddendFL \label{fig:2}
\end{figure}


%DIF >  BEGIN supplementary/tables/historical_generation
%DIF >  Source-backed publication fragment; source records retained independently.
%DIF >  Source: reports/baseline_metrics_provenance_20261020/paper_figure2_values.csv
\DIFaddbegin \begingroup
\footnotesize
\setlength{\tabcolsep}{4pt}
\begin{longtable}{llllll}
\caption{{\color{blue}Absolute historical Figure 2 summary values}}\label{tab:historical_generation}\\
\toprule
cohort & method & SSIM\_percent & PSNR & MSE\_percent & MAE\_percent \\
\midrule
\endfirsthead
\toprule
cohort & method & SSIM\_percent & PSNR & MSE\_percent & MAE\_percent \\
\midrule
\endhead
\bottomrule
\endfoot
ADNI\_FDG & Cyclegan & 67.82 & 17.26 & 2.35 & 6.47 \\
ADNI\_FDG & Swin-Unet & 75.29 & 19.82 & 1.27 & 5.76 \\
ADNI\_FDG & DDIM(3D) & 82.62 & 23.09 & 0.98 & 3.98 \\
ADNI\_FDG & SegGuidedDif & 84.85 & 23.76 & 0.57 & 3.73 \\
ADNI\_FDG & Rectified Flow & 85.69 & 24.96 & 0.53 & 3.52 \\
ADNI\_FDG & Ours(single) & 87.48 & 26.51 & 0.43 & 3.08 \\
ADNI\_FDG & Ours & 89.94 & 28.72 & 0.33 & 2.96 \\
ADNI\_AV45 & Cyclegan & 59.56 & 15.43 & 3.24 & 8.39 \\
ADNI\_AV45 & Swin-Unet & 69.16 & 17.54 & 2.67 & 7.95 \\
ADNI\_AV45 & DDIM(3D) & 73.37 & 21.33 & 1.73 & 7.01 \\
ADNI\_AV45 & SegGuidedDif & 80.29 & 21.99 & 1.12 & 6.92 \\
ADNI\_AV45 & Rectified Flow & 85.43 & 24.74 & 0.69 & 4.31 \\
ADNI\_AV45 & Ours(single) & 87.37 & 27.11 & 0.51 & 3.16 \\
ADNI\_AV45 & Ours & 88.21 & 28.09 & 0.39 & 2.73 \\
OASIS & Cyclegan & 40.27 & 12.38 & 4.86 & 10.78 \\
OASIS & Swin-Unet & 51.34 & 14.96 & 3.96 & 10.16 \\
OASIS & DDIM(3D) & 68.59 & 17.05 & 3.48 & 8.24 \\
OASIS & SegGuidedDif & 71.85 & 19.63 & 2.51 & 7.59 \\
OASIS & Rectified Flow & 80.44 & 21.33 & 1.87 & 4.29 \\
OASIS & Ours(single) & 82.66 & 25.48 & 1.12 & 3.55 \\
OASIS & Ours & 83.97 & 26.79 & 0.94 & 3.01 \\
\end{longtable}
\noindent\parbox{\linewidth}{\footnotesize \DIFadd{21method/cohort rows x4metrics =84historical values. SSIM/MSE/MAE retain source percent scales, PSNR dB. OASIS is AV45 only. No recovered common per-subject final inputs: SD, CI and paired p unavailable. These are not the fixed W1 OOF experiment and are not paired with it. Original radar figure unchanged.}}
\endgroup

%DIF >  END supplementary/tables/historical_generation


\DIFaddend \subsection{Recapitulation of Disease-Specific Biomarker Patterns}
To assess \DIFdelbegin \DIFdel{the clinical validity of the synthesized PET images, we parcellated both synthetic and real PET scans into anatomical regions using \texttt{mri\_synthseg} (FreeSurfer v7.4.1)~\cite{fischl2012freesurfer,dalca2022synthseg}, a deep learning-based segmentation tool that provides high-resolution cortical and subcortical labeling consistent with the }\DIFdelend \DIFaddbegin \DIFadd{regional uptake patterns, the corresponding T1-weighted MRI scans were parcellated using SynthSeg~\cite{synthseg}, and the resulting }\DIFaddend Desikan--Killiany and \textit{aseg} \DIFdelbegin \DIFdel{atlases}\DIFdelend \DIFaddbegin \DIFadd{ROI masks were applied to spatially aligned real and synthetic PET images}\DIFaddend . This approach yielded a total of 98 distinct regions of interest (ROIs), including cortical, subcortical, cerebellar, and brainstem structures. Regional \DIFdelbegin \DIFdel{standardized uptake value  (SUV}\DIFdelend \DIFaddbegin \DIFadd{normalized uptake values (arbitrary units}\DIFaddend ) were computed \DIFdelbegin \DIFdel{by extracting the mean tracer uptake }\DIFdelend \DIFaddbegin \DIFadd{as the mean voxel intensity }\DIFaddend within each ROI \DIFdelbegin \DIFdel{and normalizing to the cerebellar gray matter (for \textsuperscript{18}F-AV-45 PET) or the pons (for \textsuperscript{18}F-FDG PET )}\DIFdelend \DIFaddbegin \DIFadd{from PET images normalized using per-image min--max intensity scaling}\DIFaddend . Fig.~\ref{fig:3} illustrates this analysis for key ROIs in two representative ADNI subjects, \DIFdelbegin \DIFdel{demonstrating strong regional consistency between synthetic and real images, as well as clear distinctions between syntheses from different diagnostic groups}\DIFdelend \DIFaddbegin \DIFadd{illustrating regional patterns in these selected cases}\DIFaddend .

\begin{figure}[htbp]
    \centering
    \includegraphics[width=\textwidth]{figure_3}
    \DIFdelbeginFL %DIFDELCMD < \caption{
%DIFDELCMD < \textbf{Detailed regional visualization of synthetic and real PET images.}
%DIFDELCMD < The figure displays magnified regional views of FDG-PET (left two columns) and AV45-PET (right two columns) for representative AD and CN subjects. 
%DIFDELCMD < For each modality and group, the first sub-column shows real PET, and the second sub-column shows synthetic PET. 
%DIFDELCMD < Key anatomical regions are highlighted in the bottom-right corner of each sub-column, with their names labeled at the top-center. 
%DIFDELCMD < Each row corresponds to a distinct brain region, demonstrating strong regional consistency between synthetic and real PET, and clear differences between diagnostic groups.
%DIFDELCMD < }
%DIFDELCMD <     %%%
\DIFdelendFL \DIFaddbeginFL \caption{{\color{blue}
\textbf{Detailed regional visualization of synthetic and real PET images.}
The figure displays magnified regional views of FDG-PET (left two columns) and AV45-PET (right two columns) for representative AD and CN subjects. 
For each modality and group, the first sub-column shows real PET, and the second sub-column shows synthetic PET. 
Key anatomical regions are highlighted in the bottom-right corner of each sub-column, with their names labeled at the top-center. 
Each row corresponds to a distinct brain region. Displayed uptake values are arbitrary units after per-image min--max normalization, without anatomical reference-region normalization. Cases were randomly selected within their diagnostic groups as reported by the authors; the selection seed could not be verified from the available records. These selected cases illustrate regional appearance and do not establish population-level agreement or clinical SUV/SUVR calibration.
}}
    \DIFaddendFL \label{fig:3}
\end{figure}

For a comprehensive quantitative evaluation, we performed two \DIFdelbegin \DIFdel{group-wise statistical comparisons }\DIFdelend \DIFaddbegin \DIFadd{complementary regional statistical analyses }\DIFaddend on the ADNI cohort. First, we compared synthetic and real images to \DIFdelbegin \DIFdel{validate the regional SUV precision }\DIFdelend \DIFaddbegin \DIFadd{assess the normalized regional uptake agreement }\DIFaddend of DIReCT$++$. Second, we compared \DIFdelbegin \DIFdel{synthetic images across diagnostic groups (cognitively normal, i.e. , CN, MCI, AD ) to verify their ability to capture disease-relevant patterns}\DIFdelend \DIFaddbegin \DIFadd{the AD and CN groups separately in real and synthetic PET to assess whether disease-related regional differences were retained. Figure~\ref{fig:4}A,B uses the pooled available subjects without diagnostic-group stratification; panels C--F compare AD and CN only}\DIFaddend .

Paired-sample $t$-tests between synthetic and real images revealed no statistically significant differences in \DIFdelbegin \DIFdel{mean SUV across all }\DIFdelend \DIFaddbegin \DIFadd{normalized regional uptake across the }\DIFaddend principal brain regions \DIFdelbegin \DIFdel{for the CN, MCI, and AD groups ($p>0.05$; }\DIFdelend \DIFaddbegin \DIFadd{shown in }\DIFaddend Fig.~\ref{fig:4}A, B \DIFdelbegin \DIFdel{; figs. ~\ref{fig:supp_av45_1}, \ref{fig:supp_av45_2}, \ref{fig:supp_fdg_1}, and \ref{fig:supp_fdg_2}). This high fidelity was consistent for both }\DIFdelend \DIFaddbegin \DIFadd{for the overall cohorts of both tracers ($p>0.05$). Across all 98 ROIs, no overall comparisons remained significant after Benjamini--Hochberg false discovery rate (FDR) correction (all $q>0.05$). Among the diagnosis-stratified paired comparisons, only the right pericalcarine cortex (ROI 2021) in the $^{18}$F-FDG MCI subgroup showed a significant mean paired difference after FDR correction ($q=0.0481$). The median concordance correlation coefficients (CCC) across the 98 ROIs were 0.9745 for }\DIFaddend $^{18}$F-AV-45 \DIFdelbegin \DIFdel{(}\DIFdelend \DIFaddbegin \DIFadd{and 0.9450 for $^{18}$F-FDG, with corresponding median mean absolute errors (MAE) of 0.0160 and 0.0159 (normalized intensity units). For the 16 representative tracer--ROI combinations in }\DIFaddend Fig.~\ref{fig:4}A\DIFdelbegin \DIFdel{) and }\DIFdelend \DIFaddbegin \DIFadd{, B, descriptive regression slopes ranged from 0.9794 to 1.0127 (median, 1.0007) for $^{18}$F-AV-45 and from 0.9844 to 1.0101 (median, 1.0001) for }\DIFaddend $^{18}$F-FDG\DIFdelbegin \DIFdel{(Fig. ~\ref{fig:4}B), indicating that DIReCT$++$ effectively captures the regional characteristics of the ground truth}\DIFdelend \DIFaddbegin \DIFadd{, without a consistent direction of deviation from unity. These results indicate close agreement in normalized regional uptake between synthetic and real PET, with detailed statistics for the representative ROIs provided in Supplementary Table~\ref{tab:regional_agreement}. Non-significant mean differences do not establish formal equivalence or clinical interchangeability. These historical regional-CSV results are separate from the reconstructed Full generation comparisons}\DIFaddend .

Critically, independent-sample $t$-tests demonstrated that the significant differences between diagnostic groups observed in real PET data were \DIFdelbegin \DIFdel{precisely preserved }\DIFdelend \DIFaddbegin \DIFadd{also observed }\DIFaddend in the synthetic images (Fig.~\ref{fig:4}C, D; figs.~\ref{fig:supp_disease_av45} and \ref{fig:supp_disease_fdg}). For instance, DIReCT$++$ recapitulated the statistically significant differences \DIFdelbegin \DIFdel{($p<0.001$) }\DIFdelend between AD and CN groups in the \DIFaddbegin \DIFadd{left }\DIFaddend precuneus for $^{18}$F-AV-45 (\DIFaddbegin \DIFadd{real: $q=7.73\times10^{-66}$; synthetic: $q=1.25\times10^{-63}$; }\DIFaddend Fig.~\ref{fig:4}C) and in the \DIFaddbegin \DIFadd{left }\DIFaddend hippocampus for $^{18}$F-FDG (\DIFaddbegin \DIFadd{real: $q=4.27\times10^{-15}$; synthetic: $q=1.27\times10^{-13}$; }\DIFaddend Fig.~\ref{fig:4}D). Similarly, regional differences that were non-significant in real PET data \DIFdelbegin \DIFdel{were also consistently }\DIFdelend \DIFaddbegin \DIFadd{remained }\DIFaddend non-significant in the synthetic images \DIFdelbegin \DIFdel{($p>0.05$, }\DIFdelend \DIFaddbegin \DIFadd{after FDR correction (all $q\geq0.05$ for the regions shown in }\DIFaddend Fig.~\ref{fig:4}E, F). \DIFdelbegin \DIFdel{These findings confirm that DIReCT$++$ reliably preserves the essential biomarker information necessary for clinical differentiation}\DIFdelend \DIFaddbegin \DIFadd{The AD--CN effect-size directions were consistent between real and synthetic PET for all ten regions shown in Fig.~\ref{fig:4}C, D, although the absolute Cohen's $d$ decreased in nine regions (maximum reduction, 0.0771). These findings support the preservation of group-level disease-related regional uptake patterns, while indicating variations in effect-size magnitude. Complete effect-size estimates and FDR-adjusted results for the displayed regions are provided in Supplementary Table~\ref{tab:regional_effect_sizes}. The AD--CN comparison illustrates representative differences between two diagnostic states rather than all diagnostic-group contrasts. MCI was evaluated separately through MCI-versus-CN and EMCI-versus-LMCI classification (Fig.~\ref{fig:5}B,C); the regional paired analysis also includes the original MCI categories}\DIFaddend .

\begin{figure}[htbp]
\centering
\DIFdelbeginFL %DIFDELCMD < \includegraphics[width=0.75\textwidth]{figure_4}
%DIFDELCMD <     \caption{
%DIFDELCMD < \textbf{Regional comparison of mean uptake between real and MRI-synthetic PET images.}
%DIFDELCMD < (A,B) Violin plots show regional mean uptake across cortical and subcortical regions for real and synthetic PET, comparing AV45 (A) and FDG (B). Statistical significance was assessed for each region. (C–F) Disease-related patterns. Violin plots display mean uptake in AD (red) and CN (green) subjects for selected regions. Each region uses paired half-violins: left (solid) = real PET, right (dashed) = synthetic PET; colors match across modalities. Solid lines indicate real PET significance, dashed lines indicate synthetic PET significance. (C,D) show representative disease-related regions for AV45 and FDG, (E,F) show non–disease-related regions. Cblm-WM denotes Cerebellum-White-Matter. Additional analyses are in the Supplementary Materials.
%DIFDELCMD < }
%DIFDELCMD < 

%DIFDELCMD <     %%%
\DIFdelendFL \DIFaddbeginFL \includegraphics[width=\textwidth,height=.62\textheight,keepaspectratio]{figure_4}
\caption{{\color{blue}
\textbf{Regional comparison of mean normalized uptake between real and MRI-synthetic PET images.}
(A,B) Violin plots show regional mean normalized uptake (arbitrary units) across cortical and subcortical regions for real and synthetic PET, comparing AV45 (A) and FDG (B). All regional uptake values were derived from per-image min--max-normalized PET intensities without anatomical reference-region normalization.
Panels A,B summarize the pooled available subjects without diagnostic stratification. Non-significant differences are not evidence of equivalence. (C--F) AD--CN regional comparisons. Violin plots display mean normalized uptake in AD (red) and CN (green) subjects for selected regions. Each region uses paired half-violins: upper (solid) = real PET, lower (hatched) = synthetic PET; colors match across modalities. Solid lines indicate real PET significance, dashed lines indicate synthetic PET significance. The statistical annotations in panels (C--F) follow the BH-FDR analysis: \texttt{****} corresponds to $q<10^{-4}$ and \texttt{ns} to $q\geq0.05$. (C,D) show representative regions with significant AD--CN differences for AV45 and FDG, whereas (E,F) show selected regions without significant AD--CN differences after FDR correction. Cblm-WM denotes Cerebellum-White-Matter. Additional regional agreement and effect-size analyses are provided in the Supplementary Materials.
}}
\DIFaddendFL \label{fig:4}
\end{figure}

\DIFdelbegin \par\noindent\DIFdel{\textbf{Precise Prognostic Stratification of Mild Cognitive Impairment}}\par
%DIFDELCMD < %%%
\DIFdel{To evaluate the clinical utility of the synthetic images for personalized diagnosis and prognostic stratification, we conducted classification tasks on the ADNI cohort using a 3-fold cross-validation scheme. We trained models with a DenseNet backbone~\cite{huang2017densely} to differentiate between key diagnostic pairs: (1)AD vs. CN , (2)MCI vs. CN, and }\DIFdelend \DIFaddbegin \subsection{\texorpdfstring{\DIFadd{Diagnostic Classification and MCI Stage Discrimination}}{Diagnostic Classification and MCI Stage Discrimination}}
\DIFadd{Classification used independent subject-level three-fold evaluation for AD versus CN ($n=470$), EMCI+LMCI versus CN }\DIFaddend (\DIFdelbegin \DIFdel{3)early MCI (EMCI)vs. late MCI (LMCI).
For each task, models were trained and evaluated using various input modalities, including mono-modal (real PET, synthetic PET, or MRI alone) and multi-modal (}\DIFdelend \DIFaddbegin \DIFadd{$n=741$), and EMCI versus LMCI ($n=436$). The MCI/CN task includes EMCI and LMCI, while MCI-U is excluded. The completed analysis comprises 63 imaging and 18 clinical fits (27 task--input configurations, three folds each). Results are reported as the equally weighted fold mean $\pm$ sample s.d., alongside conditional bootstrap 95\% CIs (Table~\ref{tab:classification_summary}). This is cross-sectional MCI stage discrimination, not future conversion or longitudinal prognostic prediction.
}

\DIFadd{For AD/CN, }\DIFaddend MRI+\DIFdelbegin \DIFdel{real PET; MRI }\DIFdelend \DIFaddbegin \DIFadd{Syn-FDG}\DIFaddend +\DIFdelbegin \DIFdel{synthetic PET)configurations. Performance was quantified by accuracy (ACC), sensitivity (SEN), specificity (SPE), and the area under the receiver operating characteristic curve (AUC ).
}%DIFDELCMD < 

%DIFDELCMD < %%%
\DIFdel{The synthesized PET images from DIReCT$++$ yielded substantial performance gains that are close to real images across all tasks (}\DIFdelend \DIFaddbegin \DIFadd{Syn-AV45 achieved ACC $88.08\pm2.06\%$ and AUC $93.82\pm3.21\%$, compared with MRI-only ACC $79.56\pm6.20\%$ and AUC $87.02\pm5.44\%$ (Fig.~\ref{fig:5}A). The paired AUC difference was $+6.80$ percentage points (95\% CI $[+3.39,+10.22]$). For MCI/CN, MRI+Syn-both had ACC $65.45\pm1.64\%$ and AUC $69.46\pm1.34\%$, versus MRI-only $67.48\pm0.23\%$ and $71.59\pm1.80\%$; the paired AUC difference was $-2.13$ points (95\% CI $[-6.30,+2.10]$; }\DIFaddend Fig.~\ref{fig:5}\DIFdelbegin \DIFdel{). In the fundamental task of AD recognition (Fig.
~\ref{fig:5}A), the combination of MRIwith synthetic $^{18}$F-AV-45 and $^{18}$F-FDG PET achieved high performance (ACC : $93.47\%$, AUC : $95.65\%$), significantly outperforming MRI alone (ACC: $87.69\%$, AUC: $91.88\%$) and performing comparably to the combination of MRI with real multi-tracer PET (ACC: $94.27\%$, AUC: $96.10\%$).
}\DIFdelend \DIFaddbegin \DIFadd{B). Thus the synthetic combination did not improve all tasks.
}\DIFaddend 

\DIFdelbegin \DIFdel{The value of synthetic PET was even more pronounced in the challenging task of MCI diagnosis }\DIFdelend \DIFaddbegin \DIFadd{For EMCI/LMCI, MRI+Syn-both had ACC $65.82\pm2.06\%$ and AUC $66.33\pm0.54\%$, compared with MRI-only $64.90\pm3.09\%$ and $62.81\pm2.50\%$. Its balanced-ACC difference favored the synthetic combination by $6.91$ points (95\% CI $[2.02,11.98]$), while the ACC and AUC difference CIs crossed zero }\DIFaddend (Fig.~\ref{fig:5}\DIFdelbegin \DIFdel{B). While overall performance decreased, the improvements afforded by synthetic multi-tracer PET were more significant (e.g.
}\DIFdelend \DIFaddbegin \DIFadd{C). The Syn-minus-Real dual-PET AUC intervals crossed zero for all three tasks, which does not establish equivalence or interchangeability (Table~\ref{tab:classification_comparisons}).
}

\DIFadd{Clinical-only AUC was $100.00\pm0.00\%$ for AD/CN, $98.93\pm0.70\%$ for MCI/CN, and $99.94\pm0.08\%$ for EMCI/LMCI. MRI+Clinical AUC was $99.17\pm0.96\%$}\DIFaddend , \DIFdelbegin \DIFdel{SEN: $78.17\%$ for }\DIFdelend \DIFaddbegin \DIFadd{$93.07\pm6.10\%$, and $89.60\pm9.08\%$, respectively. }\DIFaddend MRI+\DIFdelbegin \DIFdel{synthetic PET vs. $71.43\%$ for MRIalone), underscoring their utility for early disease detection}\DIFdelend \DIFaddbegin \DIFadd{Syn-both AUC was below both clinical baselines in each task, with the specified paired intervals below zero. These observations do not support independent incremental diagnostic value over the clinical variables, which include diagnosis-related scores and diagnosis-conditioned fills. The zero observed AD/CN error for Clinical-only and its degenerate empirical bootstrap CI do not establish perfect population generalization}\DIFaddend .

\DIFdelbegin \DIFdel{Crucially, the combination of baseline MRI and synthetic multi-tracer PET also enabled accurate prognostic stratification, achieving $81.96\%$ accuracy in differentiating EMCI from LMCI, compared to $77.65\%$ using MRI alone }\DIFdelend \DIFaddbegin \begin{figure}[htbp]
\centering
\includegraphics[width=\textwidth,height=0.68\textheight,keepaspectratio]{figure_5}
\caption{{\color{blue}\textbf{Cross-sectional classification under the revised three-fold protocol.} Rows A--C correspond to AD/CN, EMCI+LMCI/CN and EMCI/LMCI. Within each row, the five metric columns compare nine input conditions: MRI; MRI+Real FDG; MRI+Syn FDG; MRI+Real AV45; MRI+Syn AV45; MRI+Real-both; MRI+Syn-both; Clinical-only; and MRI+Clinical. Points are equally weighted three-fold means; horizontal error bars show sample s.d. in percentage points, not 95\% CIs. ACC, AUC, SEN, SPE and balanced ACC are reported with their CIs in Table~\ref{tab:classification_summary}. The results do not demonstrate longitudinal prognosis, PET equivalence or benefit independent of diagnosis-related clinical inputs.}}
\label{fig:5}
\end{figure}

\DIFadd{Supplemental ROC curves use fold-wise averaging and pointwise confidence bands }\DIFaddend (Fig.~\DIFdelbegin \DIFdel{\ref{fig:5}C). This result highlights the potential of DIReCT$++$ to stratify patients at the prodromal stage based on their risk of progression}\DIFdelend \DIFaddbegin \DIFadd{\ref{fig:classification_roc}). Paired DeLong tests are reported by outer fold with their raw and family-specific Holm-adjusted $p$ values (Table~\ref{tab:classification_delong}); they are not combined into a three-fold overall DeLong test. All comparisons are conditional on the fixed trained models and splits}\DIFaddend .

\DIFdelbegin %DIFDELCMD < \begin{figure}[htbp]
%DIFDELCMD <     \centering
%DIFDELCMD <     \includegraphics[width=\textwidth]{figure_5}
%DIFDELCMD <     \caption{\textbf{Diagnostic classification performance using MRI-synthetic PET images across different clinical tasks.} (A–C) Bar plots illustrate the classification performance—accuracy (ACC), sensitivity (SEN), specificity (SPE), and area under the ROC curve (AUC)—of eleven input-modality combinations evaluated for three diagnostic tasks: (A) Alzheimer’s disease (AD) vs. cognitively normal (CN), (B) mild cognitive impairment (MCI) vs. CN and (C) early MCI (EMCI) vs. late MCI (LMCI). The input modalities include: MRI only, real FDG-PET, synthetic FDG-PET, real AV45-PET, synthetic AV45-PET, MRI + real FDG-PET, MRI + synthetic FDG-PET, MRI + real AV45-PET, MRI + synthetic AV45-PET, MRI + real FDG + real AV45, and MRI + synthetic FDG + synthetic AV45. For visual clarity, bars corresponding to synthetic PET inputs are shown with the same colour scheme as their real PET counterparts but with dashed (hatched) fills. The results demonstrate that MRI-synthetic PET images achieve comparable diagnostic accuracy to real PET scans, highlighting their clinical utility for downstream classification tasks.}
%DIFDELCMD <     \label{fig:5}
%DIFDELCMD < \end{figure}
%DIFDELCMD < 

%DIFDELCMD < %%%
\DIFdel{Taken together,these findings confirm the translational potential of DIReCT$++$, as its synthetic outputs serve as effective proxies for real PET biomarkers in critical downstream diagnostic and prognostic tasks}\DIFdelend \DIFaddbegin \DIFadd{Age ($<75$ versus $\geq75$ years) and sex were analyzed separately using existing classification predictions (Tables~\ref{tab:r2-16-ad-cn},\ref{tab:r2-16-mci-cn},\ref{tab:r2-16-emci-lmci}). For AD/CN, MRI+Syn-both minus MRI balanced-ACC differences ranged from 7.45 to 10.97 percentage points across the four subgroups. EMCI/LMCI point differences were positive in all four subgroups, including 9.64 points for age $\geq75$ (95\% CI $[1.23,18.05]$); MCI/CN did not show consistent gains. The subgroup intervals are marginal and unadjusted. These stratified descriptions do not remove demographic confounding or establish age/sex interactions}\DIFaddend .



\section{Discussion}
\DIFdelbegin \DIFdel{Our study introduces }\DIFdelend DIReCT$++$ \DIFdelbegin \DIFdel{, a VLM-modulated rectified flow framework for the precise synthesis of }\DIFdelend \DIFaddbegin \DIFadd{combines shared tracer-conditioned 3D rectified flow with biomedical text guidance for FDG- and AV45-PET synthesis from MRI. Its contribution extends the single-tracer DIReCT framework to }\DIFaddend multi-tracer \DIFdelbegin \DIFdel{PET from routine MRIunder fundamental text guidance. We demonstrate that this approach not only generates images of superior fidelity and generalizability but, more importantly, yields synthetic biomarkers that accurately recapitulate disease-specific patterns and enable precise prognostic stratification of MCI. By effectively making multi-modal biomarker profiling accessible from a single MRI, DIReCT$++$ represents a scalable, AI-empowered strategy that addresses a critical bottleneck in the early diagnosis and management of Alzheimer's disease}\DIFdelend \DIFaddbegin \DIFadd{modeling and a separately evaluated affine-adapter formulation. We assessed image quality, normalized regional uptake and cross-sectional classification, reporting historical and reconstructed analyses separately}\DIFaddend .

The \DIFdelbegin \DIFdel{superior performance of DIReCT$++$ stems from its core methodological innovations, which are designed to overcome the fundamental challenges of cross-modal medical image synthesis. First, the multi-task rectified flow (RF) architecture provides an efficient and stable framework for learning the complex, high-dimensional mapping from MRI to multi-tracer PET, enabling rapid, single-step generation of high-fidelity images. The critical innovation, however, lies in the modulation of this flow by a domain-adapted VLM. This integration directly addresses the ill-posed inverse problem of predicting molecular pathology from macrostructural anatomy. By conditioning the synthesis on patient-specific text guidance that encodes both individual clinical scores and general tracer knowledge, the model resolves the inherent ambiguity where similar MRI presentations correspond to distinct pathological states. Our results confirm this mechanism: DIReCT$++$ not only achieves superior pixel-level fidelity (Fig.~\ref{fig:2}), but also accurately recapitulates disease-specific patterns of amyloid deposition and hypometabolism (Fig.~\ref{fig:4}). Collectively, the VLM guidance and multi-task flow setting act as a powerful regularizer, steering the generation toward biologically plausible and subject-specific outcomes, thereby advancing beyond population averages to enable personalized biomarker profiling}\DIFdelend \DIFaddbegin \DIFadd{generation results show both advantages and tradeoffs. Full had better masked image-quality means than the clinical regression U-Net and MRI-only Swin-UNETR, while some AV45 error effects were small and did not imply universal individual benefit. Whole-volume comparisons favored TabularMLP for PSNR/MSE and TracerOnly for several endpoints. The adapter analysis supports control of tracer-embedding geometry, but it does not establish a causal improvement in PET generation. Fixed-weight conditioning interventions indicate that patient features can influence outputs; independently initialized Full/TracerOnly or text/tabular models do not isolate the clinical or VLM component as the sole cause of their performance differences}\DIFaddend .

\DIFdelbegin \DIFdel{Beyond its technical merits, the primary significance of DIReCT$++$ lies in its potential to democratize access to critical biomarker profiling, thereby addressing a major bottleneck in the early and precise management of dementia. By generating accurate proxies for both amyloid and metabolic PET from a single routine MRIscan, our framework effectively translates the gold-standard biological definition of AD into a scalable tool for clinical practice. The clinical viability of this approach is strongly evidenced by our results on personalized MCI diagnosis and stratification (Fig.~\ref{fig:5}B, C), where the combination of MRI and synthetic multi-tracer PETachieved performance comparable to the resource-intensive ideal of MRI with real PET, significantly outperforming models using MRI or real PET alone}\DIFdelend \DIFaddbegin \DIFadd{The potential value of this framework is to improve access to biomarker-related imaging information for multimodal imaging research. The classification findings are task dependent: MRI+Syn-both improved AD/CN relative to MRI, did not improve MCI/CN, and showed mixed EMCI/LMCI endpoints. Clinical-only and MRI+Clinical AUC exceeded MRI+Syn-both in the completed comparisons. Because clinical inputs include diagnosis-related assessments and diagnosis-conditioned fills, these findings cannot be used to establish independent diagnostic value of synthetic PET. Similar observed Real/Syn performance does not demonstrate equivalence or molecular validity}\DIFaddend .

\DIFdelbegin \DIFdel{This capability has profound implications. It could enable widespread screening of individuals at the preclinical or prodromal stages in primary care settings, identifying those with underlying AD pathology who are most likely to benefit from disease-modifying therapies. Furthermore, the ability to serially generate synthetic PET biomarkers from longitudinal MRI without additional radiation exposure opens new avenues for cost-effective monitoring of disease progression and treatment response . DIReCT$++$ thus moves us closer to a future where multi-modal diagnostic precision is accessible as part of the standard clinical workflow, unconstrained by current resource limitations}\DIFdelend \DIFaddbegin \DIFadd{Synthetic PET may support exploratory MRI-based analyses when PET is unavailable, but clinical screening or treatment selection would require independent molecular biomarker confirmation and prospective validation. No longitudinal conversion, progressive-versus-stable MCI, survival or treatment-response endpoints were analyzed. EMCI/LMCI classification is therefore MCI stage discrimination; it does not predict an individual's progression to AD. Whether serially synthesized PET can reflect progression or treatment response remains a question for longitudinal studies}\DIFaddend .

\DIFdelbegin \DIFdel{When contextualized within the field of }\DIFdelend \DIFaddbegin \DIFadd{Within }\DIFaddend medical image synthesis, \DIFdelbegin \DIFdel{the advancements of DIReCT$++$ represent a qualitative leap beyond prior approaches. While existing generative adversarial networks and diffusion modelshave demonstrated the feasibility of MRI-to-PET translation, they have primarily focused on achieving visual realism.
}\DIFdelend DIReCT$++$ \DIFdelbegin \DIFdel{moves beyond this by achieving precise generation of subject-specific biomarkers that are clinically actionable. This is a direct consequence of its VLM-modulated architecture, which incorporates domain knowledge as a critical regularizer, a feature largely absent in preceding purely, data-driven approaches. The core principle of using a domain-adapted foundation model to guide an efficient generative process establishes a new paradigm that is extendible well beyond the synthesis of $^{18}$F-AV-45 and $^{18}$F-FDG PET. By prioritizing clinically grounded precision over mere voxel-level fidelity, our work points the way toward a new generation of AI models that are truly co-designed with clinical application in mind. }\DIFdelend \DIFaddbegin \DIFadd{evaluates disease-related regional uptake patterns and downstream classification in addition to image-level metrics. The learned one-step approximation permits a single tracer-conditioned forward pass; measured forward latency does not constitute an end-to-end workflow or an equal-protocol speedup over competing models.
}\DIFaddend 

%DIF <  Limitations
%DIF <  1. guidance from more modalities, and more sophiscated prompt learning strategies to align heterogenous data
%DIF <  2. synthesis of Tau-PET, and application to preclinical identificaton and longitudinal prediction
%DIF <  3. technically, include downstream guidance
%DIF <  4. evaluation on more heterogenous data
\DIFdelbegin \DIFdel{Several limitations of this study present clear avenues for future work.
First, the current text guidance relies on fundamental patient information and imaging knowledge. Incorporating additional clinically accessible data, such as blood biomarkers~\cite{hansson2023blood}, could further enhance }\DIFdelend \DIFaddbegin \DIFadd{Several limitations require explicit consideration. Clinical-record reconstruction showed extensive diagnosis-conditioned filling of missing scores, which can introduce group information into conditioning and classification. Source weight entries also contain unresolved validity issues. The original random filling states and complete actual acquisition mappings were not recovered. The independent three-fold classifier evaluation verifies its own partitions, but the historical synthetic-image inputs could not be bound directly to the checkpoints in the reconstructed generation comparisons. These limitations constrain claims of reproducibility and independent clinical discrimination.
}

\DIFadd{Regional measurements use per-image min--max-normalized intensities in arbitrary units, without established anatomical-reference SUV/SUVR calibration. No clinically justified equivalence margin, formal TOST result or valid amyloid-positivity agreement analysis was established. Close agreement, non-significant paired differences or preservation of AD--CN significance is not evidence of clinical equivalence or reliable atypical A$\beta$ status classification. Calibrated quantitative PET and prespecified, clinically meaningful equivalence margins are needed for such assessments.
}

\DIFadd{The deterministic one-step formulation provides a single output per input without calibrated }\DIFaddend subject-level \DIFdelbegin \DIFdel{precision, necessitating the development of targeted cross-modal pretraining or prompt-learning strategies. Second, while DIReCT$++$ synthesizes $^{18}$F-AV-45 and $^{18}$F-FDG, extending the framework to include tau-PET imaging would provide a more comprehensive biomarker profile for sensitive screening at the preclinical stage.
Third, to maximize clinical utility, future iterations could integrate downstream task guidance (e. g. , diagnosis or prognosis labels) directly into the generative model, enabling controllable synthesis~\cite{wang2025flexibly,chang2025controllable} tailored for specific clinical decisions. Finally, the evaluation was conducted on well-characterized but limited research cohorts (ADNI, OASIS)}\DIFdelend \DIFaddbegin \DIFadd{predictive uncertainty. An MSE-trained deterministic model may represent an average conditional mapping and cannot guarantee recovery of individual molecular pathology from MRI and text. Ensemble or stochastic uncertainty estimation remains future work. The available comparisons do not isolate every model component or establish the quality--latency tradeoff between one-step and multistep sampling. The fixed-model bootstrap intervals do not quantify training-seed or split variability.
}

\DIFadd{External evaluation was limited to historical AV45 synthesis in OASIS-3 CN/AD, with age differences and incomplete run-specific prompt records; FDG and MCI external generalizability remain unestablished}\DIFaddend . Prospective validation \DIFdelbegin \DIFdel{on more diverse , real-world clinical populations(including data from thick-slice or low-field MRI ~\cite{sorby2024portable}) is essential to fully establish the generalizability and robustness of the framework}\DIFdelend \DIFaddbegin \DIFadd{in more diverse clinical populations, including alternative MRI acquisitions~\cite{sorby2024portable}, is needed. Additional clinical modalities such as blood biomarkers~\cite{hansson2023blood} may provide complementary information, but their incremental value must be evaluated independently. The 2024 Alzheimer's Association criteria emphasize biological diagnosis using validated core biomarkers, including amyloid PET and sufficiently accurate blood-based assays such as plasma p-tau217~\cite{jack2024revised}. Synthetic PET has not been validated as a substitute for molecular confirmation, and its added value relative to those biomarker assessments remains to be determined}\DIFaddend .

In conclusion, \DIFdelbegin \DIFdel{we have presented DIReCT$++$, a framework that synergizes rectified flow with domain-adapted vision-language modeling to achieve precise synthesis of multi-tracer PET from routine MRI. By generating clinically actionable biomarkers that enable accurate diagnosis and prognostic stratification of MCI, our work effectively bridges the gap between the biological definition of AD and its clinical application. }\DIFdelend DIReCT$++$ \DIFdelbegin \DIFdel{represents a significant step toward scalable, accessible, and precise tools for early intervention in neurodegenerative disorders}\DIFdelend \DIFaddbegin \DIFadd{provides a shared tracer-conditioned framework for PET synthesis with biomedical text inputs and learned one-step inference. The reported image-quality and regional results support further methodological investigation, while mixed classification effects, clinical-input circularity and unresolved source limitations constrain clinical interpretation. Longitudinal prognosis, clinical equivalence and replacement of molecular biomarker confirmation require separate validation}\DIFaddend .



\section{Materials and Methods}
\subsection{Experimental Design}

The clinical management of Alzheimer's disease (AD) is limited by the
practical challenges associated with routine multi-tracer positron emission
tomography (PET), including acquisition cost, radiation exposure, and
accessibility. This study was designed to investigate whether molecular PET
images associated with Alzheimer's disease pathology could be synthesized
from a single structural MRI scan while preserving biologically and
clinically relevant information.

We developed DIReCT$++$, an MRI-conditioned generative framework for
subject-specific synthesis of amyloid and FDG PET images. The framework was
designed to address the inherently ill-posed nature of MRI-to-PET translation
by incorporating patient-specific clinical information and tracer-related
knowledge into the generation process.

The study was evaluated using a multi-level validation strategy. Paired MRI
and PET data from the Alzheimer's Disease Neuroimaging Initiative (ADNI) were
used for model development and internal evaluation, while an independent
Open Access Series of Imaging Studies 3 (OASIS-3) cohort was used for
external validation. The evaluation framework included three predefined
components: image-level fidelity assessment between synthetic and real PET,
regional consistency analysis of PET-derived biomarkers, and downstream
clinical classification and mild cognitive impairment stratification.


\subsection{\texorpdfstring{DIReCT$++$}{DIReCT++}}
DIReCT$++$ \DIFdelbegin \DIFdel{is a VLM-modulated rectified flow (RF) framework for synthesizing $^{18}$F-AV-45 and $^{18}$F-FDG PET images from a }\DIFdelend \DIFaddbegin \DIFadd{combines a frozen BiomedCLIP text encoder with a conditional 3D rectified flow model to synthesize FDG- or AV45-PET from }\DIFaddend T1-weighted MRI \DIFdelbegin \DIFdel{scan, conditioned on text prompts derived from patient information and tracerknowledge }\DIFdelend \DIFaddbegin \DIFadd{and tracer/patient text guidance }\DIFaddend (Fig.~\ref{fig:1}). The \DIFdelbegin \DIFdel{framework integrates a pre-trained VLM, fine-tuned from BiomedCLIP via prompt learning, with a conditional 3D rectified flow model. The VLM encodes text guidance, which is incorporated into the flow model via cross-attention to achieve subject-wise precision. The prompt learning strategy is designed to maximize intra-tracer text-image alignment while preserving inter-tracer distinctions, reflecting the complementary nature of the biomarkers}\DIFdelend \DIFaddbegin \DIFadd{affine-adapter formulation aligns tracer-specific text features with PET features while regulating cross-tracer representational similarity. Adapter parameters are optimized separately from the frozen encoder and generator. We distinguish this formulation and its representation-level analyses from generator configurations that use raw frozen tracer embeddings}\DIFaddend . The following subsections \DIFdelbegin \DIFdel{detail the architecture and training}\DIFdelend \DIFaddbegin \DIFadd{specify the adaptation and conditioning procedures}\DIFaddend .

\subsubsection{Domain-Adapted VLM}
To incorporate domain-specific knowledge into the generation process, we develop an adaptive guidance mechanism that leverages a frozen, pre-trained BiomedCLIP model. Specifically, we learn two \textit{affine} transformations, $Adapt_f$ and $Adapt_a$, to align the textual embeddings of FDG and AV-45 PET modalities with their respective image feature spaces. Let $c_f, c_a$ denote the BiomedCLIP textual embeddings for FDG and AV-45, respectively. \DIFaddbegin \DIFadd{The frozen BiomedCLIP text and image encoders, including their projection layers, produce 512-dimensional features. For the affine-adapter procedure, each preprocessed 3D PET volume is represented by 160 two-dimensional slices along its third array dimension. Each slice is intensity-normalized, resized to $224\times224$ pixels and converted to a three-channel image before encoding; the slice features are averaged into a subject-level PET representation. The same extraction procedure is used for both tracers in the revised adapter analysis.
}\DIFaddend For simplicity, we adopt lightweight affine transformations parameterized by a scalar scale and a vector bias:
\[
Adapt_f(c_f) = \alpha_1 \, c_f + \alpha_2, \quad
Adapt_a(c_a) = \beta_1 \, c_a + \beta_2,
\]
where \DIFdelbegin \DIFdel{$\alpha_1, \beta_1 $ are learnable scalar weights and $\alpha_2, \beta_2 $ are learnable }\DIFdelend \DIFaddbegin \DIFadd{$\alpha_1,\beta_1$ are scalar scales and $\alpha_2,\beta_2$ are 512-dimensional }\DIFaddend bias vectors. \DIFdelbegin \DIFdel{We choose scalar scaling rather than full linear projection to minimize trainable parameters and avoid overfitting,given the limited size of medical imaging datasets. The optimal parameters are obtained by solving a constrained optimization problem:
}\DIFdelend \DIFaddbegin \DIFadd{Each adapter has 513 trainable parameters (1,026 in total). The parameters are optimized using the objective
}\DIFaddend \begin{equation}
\DIFdelbegin %DIFDELCMD < \begin{aligned}
%DIFDELCMD < & \underset{Adapt_f,\, Adapt_a}{\arg \min} \quad 
%DIFDELCMD <    \mathcal{L}_{\text{align}} = \big(1 - \operatorname{sim}(Adapt_f(c_f), \mathbb{E}[f_f])\big) + \big(1 - \operatorname{sim}(Adapt_a(c_a), \mathbb{E}[f_a])\big) \\
%DIFDELCMD < & \text{subject to} \quad 
%DIFDELCMD <    \operatorname{sim}\big(Adapt_f(c_f), Adapt_a(c_a)\big) \geq \tau,
%DIFDELCMD < \end{aligned}%%%
\DIFdelend \DIFaddbegin \begin{aligned}
\mathcal{L}_{\mathrm{align}}={}&-\frac{1}{N_f}\sum_{i=1}^{N_f}\operatorname{sim}(Adapt_f(c_f),f_{f,i})\\
&-\frac{1}{N_a}\sum_{j=1}^{N_a}\operatorname{sim}(Adapt_a(c_a),f_{a,j})\\
&+\lambda_{\mathrm{cross}}\left[\operatorname{sim}(Adapt_f(c_f),Adapt_a(c_a))-\tau\right]^2,
\end{aligned}\DIFaddend 
\label{eq:align_opt}
\end{equation}
where \DIFdelbegin \DIFdel{$\operatorname{sim}(\cdot, \cdot)$ denotes }\DIFdelend \DIFaddbegin \DIFadd{$\operatorname{sim}$ is }\DIFaddend cosine similarity, \DIFaddbegin \DIFadd{$f_{f,i}$ and $f_{a,j}$ are normalized subject-level PET features and $N_f,N_a$ count the available training-fold features. The settings are $\lambda_{\mathrm{cross}}=100$ }\DIFaddend and \DIFdelbegin \DIFdel{$\tau \in [0,1]$ is a predefined threshold . The constraint ensures the cross-modality guidance vectors maintain a certain level of semantic consistency. }%DIFDELCMD < 

%DIFDELCMD < %%%
\DIFdel{By solving~}%DIFDELCMD < \eqref{eq:align_opt}%%%
\DIFdel{, we obtain the modality-specific guidance vectors $Adapt_f(c_f)$ and $Adapt_a(c_a)$, which are subsequently used to steer the rectified flow translation process in a domain-informed manner}\DIFdelend \DIFaddbegin \DIFadd{$\tau=0.82$. This is a soft squared penalty targeting a reference similarity, not a hard lower-bound constraint or a threshold optimized for generation quality. The target is numerically consistent with the original tracer-text cosine similarity of 0.8205 before adaptation. The two adapters are optimized jointly using Adam (learning rate $10^{-3}$, 5,000 updates); the encoders remain frozen. Adapted vectors remain fixed when used in the affine-adapter conditioning configuration. The reconstructed Full generator uses raw tracer embeddings, while the separate adapter sensitivity analysis tests the representation geometry under Eq.~}\eqref{eq:align_opt}\DIFaddend .


\subsubsection{Rectified Flow}
\DIFdelbegin \DIFdel{The core of }\DIFdelend DIReCT$++$ \DIFdelbegin \DIFdel{is a continuous generative dynamics framework grounded in the principle of rectified flow , designed to construct a near-optimal, domain-informed trajectory that minimizes sampling complexity while preserving the structural and semantic fidelity of medical image synthesis}\DIFdelend \DIFaddbegin \DIFadd{models MRI-to-PET translation using a conditional rectified flow velocity field}\DIFaddend . Specifically, given two heterogeneous medical distributions---$\mathbf{x}_0 \sim \mathcal{M}_{\text{MRI}}$ and $\mathbf{x}_1 \sim \mathcal{P}_{\text{PET}}$ (multi-tracer)---we model the transformation between them via a learnable velocity field $V_\theta$\DIFdelbegin \DIFdel{denotes the }\DIFdelend \DIFaddbegin \DIFadd{, where $\theta$ denotes its }\DIFaddend network parameters and \DIFdelbegin \DIFdel{$c$ encodes modality-specific contextual cues (e.g., anatomical priors or clinical metadata)}\DIFdelend \DIFaddbegin \DIFadd{$\mathbf{c}$ is the conditioning context containing tracer and clinical text features}\DIFaddend . This velocity field, parameterized by a U-Net architecture tailored for high-dimensional medical imagery, governs the continuous transport of samples along a deterministic path governed by an ordinary differential equation (ODE):
\begin{equation}
\frac{d\mathbf{x}_t}{dt} = V_\theta\big(\mathbf{x}_t,\, \mathbf{c},\, t\big), \quad t \in [0, 1].
\label{eq:rf_ode}
\end{equation}

\DIFdelbegin \DIFdel{The modality-specific guidance embeddings $\text{Adapt}_f(c_f)$ and $\text{Adapt}_a(c_a)$ are pre-computed by optimizing Eq.~}%DIFDELCMD < \eqref{eq:align_opt} %%%
\DIFdel{and remain fixed throughout training and inference. These pre-computed vectors are subsequently used to steer the rectified flow translation process in a domain-informed manner. Specifically, for each subject, the composite conditioning context $\mathbf{c} $ is formed by stacking the appropriate modality guidance vector with a subject-specific clinical embedding along the token dimension:
}\DIFdelend \DIFaddbegin \DIFadd{For each participant, the conditioning context stacks the patient feature before the fixed target-tracer feature:
}\DIFaddend \begin{equation}
\mathbf{c}=\DIFdelbegin %DIFDELCMD < \begin{bmatrix}
%DIFDELCMD < \text{Adapt}_m(c_m) \\
%DIFDELCMD < c_{\text{demo}}
%DIFDELCMD < \end{bmatrix}%%%
\DIFdelend \DIFaddbegin \begin{bmatrix}c_{\mathrm{demo}}\\g_m\end{bmatrix}\DIFaddend ,
\label{eq:composite_conditioning}
\end{equation}
where \DIFdelbegin \DIFdel{$m \in \{f, a\}$ denotes the target PET tracer ($f$: FDG ; $a$: }\DIFdelend \DIFaddbegin \DIFadd{$m\in\{f,a\}$ indicates FDG or }\DIFaddend AV45\DIFdelbegin \DIFdel{). Here, $\text{Adapt}_m(c_m)$ is the fixed guidance embedding obtained from the alignment optimization, and $c_{\text{demo}} $ is a frozen BiomedCLIP text embedding derived from a standardized clinical prompt encoding the individual’s demographic and cognitive profile (e.g., }\DIFdelend \DIFaddbegin \DIFadd{. The raw tracer text feature is $c_m$ and its affine-adapted form is $Adapt_m(c_m)$; $g_m$ denotes the feature actually supplied by the specified configuration. In the reconstructed Full generator, $g_m=c_m$. The patient feature $c_{\mathrm{demo}}$ is a 512-dimensional frozen BiomedCLIP embedding of a prompt encoding }\DIFaddend age, sex, weight, \DIFdelbegin \DIFdel{and neuropsychological scores )}\DIFdelend \DIFaddbegin \DIFadd{Global CDR, MMSE, GDS, FAQ and NPI-Q, with descriptions of the assessments and their medical interpretations. Logical Memory II and an explicit diagnostic label are not prompt fields, whereas CDR and other scores are diagnosis-related, and diagnosis-conditioned filling can encode group information indirectly. The prompt template and tracer texts are provided in Supplementary Methods. Patient prompts are padded and truncated to a maximum of 256 tokens in the reconstructed implementation}\DIFaddend .

\DIFdelbegin \DIFdel{This }\DIFdelend \DIFaddbegin \DIFadd{The }\DIFaddend two-token context \DIFdelbegin \DIFdel{representation is injected into the velocity U-Net via }\DIFdelend \DIFaddbegin \DIFadd{is injected by }\DIFaddend cross-attention \DIFdelbegin \DIFdel{layers at multiple resolutions, enabling spatial features at each scale to dynamically attend to both the target imaging modality and the patient’s unique clinical attributes. Critically, no additional trainable parameters are introduced for clinical metadata encoding—all semantic information is derived from the pre-trained BiomedCLIP model—thereby enhancing generalizability and mitigating overfitting in data-limited medical imaging scenarios}\DIFdelend \DIFaddbegin \DIFadd{only at the deepest of the four U-Net resolution levels (attention levels: false, false, false, true); the conditioning feature dimension is 512. No additional trainable clinical text encoder is introduced: patient features are obtained from the frozen BiomedCLIP model}\DIFaddend .

This ODE governs the transition from the initial state $\mathbf{x}_0 \in \mathcal{M}_{\text{MRI}}$ to the final state $\mathbf{x}_1 \in \mathcal{P}_{\text{PET}}$, ensuring that the velocity field $V_\theta(\mathbf{x}_t, \mathbf{c}, t)$ drives the flow along the desired path. To determine the velocity field $V_\theta$, we solve a straightforward least squares regression problem:
\begin{equation}
\min_{\theta} \int_0^1 \mathbb{E} \left\| (\mathbf{x}_1 - \mathbf{x}_0) - V_\theta(\mathbf{x}_t, \mathbf{c}, t) \right\|^2 dt, \quad \text{with} \quad \mathbf{x}_t = t\mathbf{x}_1 + (1-t)\mathbf{x}_0,
\label{eq:rf_regression}
\end{equation}
where $\mathbf{x}_t$ represents the linear interpolation of $\mathbf{x}_0$ and $\mathbf{x}_1$. In order to train the velocity U-Net, we formulate the loss function based on Eq.~\eqref{eq:rf_regression} as follows:
\begin{equation}
\mathcal{L}(\mathbf{x}_0, \mathbf{x}_1; \theta) = \mathbb{E}_{t,\, \mathbf{x}_0 \sim \mathcal{M},\, \mathbf{x}_1 \sim \mathcal{P}} \left\| (\mathbf{x}_1 - \mathbf{x}_0) - V_\theta(\mathbf{x}_t, \mathbf{c}, t) \right\|^2,
\label{eq:rf_loss}
\end{equation}
where $t \sim \mathcal{U}(0, 1)$ is sampled from the uniform distribution.
\subsubsection{Multi-Task Learning and One-Step Generation}
\DIFaddbegin \DIFadd{Extending single-tracer DIReCT~\cite{liu2025direct}, DIReCT$++$ uses a shared generator conditioned on the target tracer. }\DIFaddend To accommodate subjects for whom only one PET tracer is available, we define a unified training objective $\mathcal{L}(\theta)$ that conditionally supervises the model based on the presence of ground-truth PET data. Specifically, the total loss combines modality-specific reconstruction terms, each gated by a binary availability indicator:
\begin{equation}
\mathcal{L}(\theta) = \lambda_F \, \delta_F \, \mathcal{L}_{\text{FDG}}(\mathbf{x}_{\text{MRI}}, \mathbf{y}_f; \theta) + \lambda_A \, \delta_A \, \mathcal{L}_{\text{AV45}}(\mathbf{x}_{\text{MRI}}, \mathbf{y}_a; \theta),
\label{eq:multitask_loss}
\end{equation}
where $\mathbf{x}_{\text{MRI}}$ denotes the input magnetic resonance image, $\mathbf{y}_f$ and $\mathbf{y}_a$ are the corresponding ground-truth FDG and AV45 PET images (when available), and $\delta_F, \delta_A \in \{0,1\}$ indicate the presence ($1$) or absence ($0$) of each PET modality for a given subject. The \DIFdelbegin \DIFdel{hyperparameters $\lambda_F$ and $\lambda_A$ balance the relative contribution of each task during training}\DIFdelend \DIFaddbegin \DIFadd{loss weights were $\lambda_F=1.0$ and $\lambda_A=1.0$}\DIFaddend . This formulation ensures that supervision is applied solely to available modalities, enabling \DIFdelbegin \DIFdel{robust }\DIFdelend training in heterogeneous clinical cohorts with incomplete multi-tracer imaging.
\DIFdelbegin \DIFdel{A key advantage of rectified flow is its ability to support exact one-step sampling. In DIReCT$++$, the target PET image is generated directly from the input MRI in a single forward pass, bypassing iterative refinement. %DIF < This is made possible by the near-straight optimal transport trajectories learned by the underlying ordinary differential equation (ODE), which closely approximate the minimal-path coupling between MRI and PET domains and eliminate the need for temporal discretization. As a result, computational overhead is substantially reduced while preserving high-fidelity synthesis.
}%DIFDELCMD < 

%DIFDELCMD < %%%
\DIFdel{To further enhance the accuracy of }\DIFdelend \DIFaddbegin \DIFadd{Rectified flow supports a learned }\DIFaddend one-step \DIFdelbegin \DIFdel{generation, we incorporate a distillation step following the initial ODE-based training phase. This procedure trains the velocity network to directly predict the end-point mapping from $\mathbf{x}_0$ }\DIFdelend \DIFaddbegin \DIFadd{approximation; a single network evaluation is not exact integration of the continuous ODE. In the reconstructed implementation, epochs before 140 sample integer timesteps from 0 }\DIFaddend to \DIFdelbegin \DIFdel{$\mathbf{x}_1$, effectively collapsing the continuous flow into a single, refined inference step. The resulting }\DIFdelend \DIFaddbegin \DIFadd{999 and fit the velocity target $\mathbf{x}_1-\mathbf{x}_0$ at the interpolated input. From epoch 140 onward, $t$ is fixed at zero, so the same supervised objective fits the endpoint residual from MRI. This stage is referred to as }\DIFaddend one-step \DIFdelbegin \DIFdel{sampler is both faster and more precise, making it suitable for time-sensitive clinical applications. Mathematically, the generation process is expressed as:
}\DIFdelend \DIFaddbegin \DIFadd{distillation in the framework, but the recovered implementation does not introduce a separate teacher trajectory or a teacher-step schedule. A one-step-versus-multistep quality/latency comparison is not established by the available measurements.
}\DIFaddend \begin{equation}
\DIFdelbegin \DIFdel{\mathbf{x}}\DIFdelend \DIFaddbegin \widehat{\mathbf{x}}\DIFaddend _1=\DIFaddbegin \operatorname{clip}\DIFadd{_{[0,1]}\!}\left(\DIFaddend \mathbf{x}_0+\DIFdelbegin %DIFDELCMD < \tilde{V}%%%
\DIFdelend \DIFaddbegin \widetilde{V}\DIFaddend _\theta(\mathbf{x}_0,\mathbf{c},0)\DIFaddbegin \right)\DIFaddend ,
\label{eq:onestep_generation}
\end{equation}
where \DIFdelbegin \DIFdel{$\mathbf{x}_0 $ denotes the subject-specific MRIinput, $\mathbf{x}_1 $ the corresponding synthesized PETimages}\DIFdelend \DIFaddbegin \DIFadd{$\mathbf{x}_0$ is MRI, $\widehat{\mathbf{x}}_1$ is the selected tracer's synthetic PET}\DIFaddend , $\mathbf{c}$ \DIFdelbegin \DIFdel{the structured conditioning context encoding both anatomical priors and clinical metadata (defined }\DIFdelend \DIFaddbegin \DIFadd{is the context }\DIFaddend in Eq.~\eqref{eq:composite_conditioning}\DIFdelbegin \DIFdel{), and $\tilde{V}_\theta$ the distilled velocity field implemented by the }\DIFdelend \DIFaddbegin \DIFadd{, and $\widetilde{V}_\theta$ is the endpoint-fitted }\DIFaddend velocity U-Net. \DIFaddbegin \DIFadd{FDG and AV45 are generated in separate tracer-conditioned forward passes.
}\DIFaddend \subsection{Materials and Data Preprocessing}
This study leveraged neuroimaging data from 1,857 participants enrolled in the Alzheimer’s Disease Neuroimaging Initiative (ADNI), spanning cohorts from ADNI-1 through ADNI-4. All participants underwent structural T1-weighted MRI, while subsets also had \textsuperscript{18}F-FDG-PET ($n = 1,598$) and \textsuperscript{18}F-AV45-PET ($n = 1,365$) scans available, enabling multi-tracer PET synthesis. \DIFaddbegin \DIFadd{Only one prepared imaging record per participant was retained; repeated longitudinal visits were not treated as independent participants. The recovered clinical-record selection rule took the first CSV record matching each participant, which should not be interpreted as a nearest-date MRI--PET pairing rule. The actual acquisition identifiers and maximum MRI--PET interval could not be established for the complete imaging cohort; the recovered clinical-source Image IDs do not resolve that file-level correspondence.
}\DIFaddend 

Participants were stratified into five clinically defined diagnostic categories to reflect the full continuum of Alzheimer’s disease pathology: Cognitively Normal (CN; $n = 577$), Alzheimer’s disease dementia (AD; $n = 317$), and three subtypes of Mild Cognitive Impairment (MCI). To account for the heterogeneity of MCI, we further subclassified individuals based on delayed recall performance on the WMS-R Logical Memory II Story A task: Early MCI (EMCI; $n = 291$) denoted a milder impairment stage, whereas Late MCI (LMCI; $n = 152$) \DIFdelbegin \DIFdel{indicated proximity to dementia conversion}\DIFdelend \DIFaddbegin \DIFadd{denoted a more severe impairment stage}\DIFaddend . The remaining MCI participants \DIFdelbegin \DIFdel{without clear progression status }\DIFdelend \DIFaddbegin \DIFadd{not classified as EMCI or LMCI }\DIFaddend were grouped as MCI-unclassified (MCI-U; $n = 520$). \DIFdelbegin \DIFdel{This granular stratification facilitates a biologically grounded analysis of neuroanatomical and metabolic alterations across the AD spectrum, thereby enhancing the clinical validity of synthesized PET images}\DIFdelend \DIFaddbegin \DIFadd{These categories define cross-sectional diagnostic and MCI-stage comparisons, without establishing longitudinal prognosis}\DIFaddend .

\DIFdelbegin \DIFdel{A consistent gradient in disease severity was observed across groups. Mean }\DIFdelend \DIFaddbegin \DIFadd{In the prepared clinical inputs, mean }\DIFaddend Mini-Mental State Examination (MMSE) scores \DIFdelbegin \DIFdel{declined progressively }\DIFdelend \DIFaddbegin \DIFadd{decreased }\DIFaddend from 29.40 in CN to 20.84 in AD\DIFdelbegin \DIFdel{($p < 0.001$, trend across groups). Correspondingly, functional and neuropsychiatric burden increased with diagnostic severity, as evidenced by higher scores on the Functional Activities Questionnaire (FAQ), Geriatric Depression Scale (GDSCALE), and Neuropsychiatric Inventory–Questionnaire (}\DIFdelend \DIFaddbegin \DIFadd{; other clinical scores also varied across diagnostic groups (Table~\ref{tab:supp_adni}). Clinical information was selected by participant identity, taking the first record in the source CSV order, rather than the nearest or earliest imaging visit. Missing MMSE, GDSCALE, Global CDR, FAQ and }\DIFaddend NPI-Q \DIFdelbegin \DIFdel{) in AD relative to CN and MCI. Significant intergroup differences were also observed in demographic variables, including age and sex distribution (see Table~\ref{tab:supp_adni}).
}\DIFdelend \DIFaddbegin \DIFadd{scores were filled using diagnostic-group-conditioned rules. Of 9,285 score cells, 1,751 were observed and retained, while 7,534 were missing in those source records and had values consistent with the recovered filling rules. The resulting score distributions include simulated inputs and should not be interpreted as independent evidence of between-group clinical differences. Group-wise missingness is reported in Table~\ref{tab:clinical_missingness}; no new imputation was performed for the revised analyses. The prepared weight data retain unresolved source-value issues, including CN entries ranging from 0 to 227 kg.
}\DIFaddend 

\DIFdelbegin \DIFdel{To further validate the generalizability of our model beyond ADNI, we performed an independent external evaluation using data from }\DIFdelend \DIFaddbegin \DIFadd{External evaluation used }\DIFaddend the Open Access Series of Imaging Studies (OASIS-3), a longitudinal \DIFdelbegin \DIFdel{, multi-modal }\DIFdelend neuroimaging and clinical cohort focused on normal aging and Alzheimer\DIFdelbegin \DIFdel{’}\DIFdelend \DIFaddbegin \DIFadd{'}\DIFaddend s disease (\DIFdelbegin \DIFdel{available at: }%DIFDELCMD < \url{https://www.medrxiv.org/content/10.1101/2019.12.13.19014902v1}%%%
\DIFdel{). To ensure methodological consistency with our primary analysis, we restricted our OASIS-3 analysis to individuals with \textsuperscript{18}F-AV45 PET scans processed through the standardized }\DIFdelend \DIFaddbegin \url{https://www.oasis-brains.org}\DIFadd{). The retained external cohort comprised 419 participants with AV45-PET processed using the }\DIFaddend Pet Unified Pipeline (PUP)\DIFdelbegin \DIFdel{, which provides quantitatively reliable, cross-site standardized SUV metrics. This yielded a high-quality cohort of 419 participants . In contrast, \textsuperscript{18}F-FDG-PET data from a }\DIFdelend \DIFaddbegin \DIFadd{: 334 CN and 85 AD, with no MCI participants (Table~\ref{tab:supp_oasis}). FDG-PET from the }\DIFaddend smaller subset (\DIFdelbegin \DIFdel{$n = 117$) were excluded due to lack of PUP processing , preserving homogeneity in image quantification. The demographic characteristics of this OASIS-3 validation cohortare detailed in Table~\ref{tab:supp_oasis}}\DIFdelend \DIFaddbegin \DIFadd{$n=117$) was excluded because it lacked the PUP processing used for this external cohort. Subsequent model inputs were min--max normalized; the external evaluation does not establish clinical SUV/SUVR calibration. The CN and AD groups differed in age distribution. These data therefore evaluate external AV45 synthesis in this CN/AD cohort, without establishing external generalizability for FDG synthesis or MCI-stage discrimination. Complete run-specific records of the external clinical prompt fields and missing-field handling were not retained, limiting reproducibility of that historical evaluation}\DIFaddend .

To ensure precise spatial correspondence between MRI and dual-modality PET images, we implemented a standardized preprocessing pipeline. T1-weighted MRI scans were first skull-stripped using FreeSurfer~\cite{Freesurfer} and subjected to histogram matching for contrast normalization. Each PET image was then rigidly registered to its corresponding MRI. Both modalities were jointly affinely normalized to the MNI152 standard space using SPM~\DIFdelbegin \DIFdel{\cite{spm}}\DIFdelend \DIFaddbegin \DIFadd{\cite{spm,friston1995spatial}}\DIFaddend , with identical transformation matrices applied to preserve cross-modality alignment. Following spatial normalization, all volumes were cropped to a common field of view of $160 \times 192 \times 160$ voxels ($1\,\text{mm}^3$ isotropic resolution) and intensity-normalized to the $[0, 1]$ range via min–max scaling. For quantitative regional analysis, \DIFdelbegin \DIFdel{brain parenchyma was }\DIFdelend \DIFaddbegin \DIFadd{T1-weighted MRI scans were }\DIFaddend parcellated into 98 anatomical \DIFdelbegin \DIFdel{structures }\DIFdelend \DIFaddbegin \DIFadd{regions }\DIFaddend using SynthSeg~\cite{synthseg}\DIFaddbegin \DIFadd{, and the resulting ROI masks were applied to the spatially aligned PET images. ADNI provided preprocessed PET images upstream of this pipeline. Spatial registration and intensity normalization do not establish comprehensive scanner-specific harmonization; acquisition and upstream processing details could not be verified for every actual model input}\DIFaddend .

\DIFdelbegin \DIFdel{Model evaluation was conducted under a }\DIFdelend \DIFaddbegin \DIFadd{PET synthesis used subject-level }\DIFaddend two-fold cross-validation \DIFdelbegin \DIFdel{scheme: the cohort was randomly partitioned into two equal-sized subsets, with alternating roles as training}\DIFdelend \DIFaddbegin \DIFadd{with fixed, approximately equal partitions (929 and 928 participants; split seed 42). In each fold, training and held-out participants were disjoint. Held-out predictions from the two folds were combined into an out-of-fold (OOF) dataset for generation analyses, with 1,598 valid FDG and 1,365 valid AV45 observations before the separately specified MRI}\DIFaddend /\DIFdelbegin \DIFdel{validation and test sets to ensure unbiased performance estimation}\DIFdelend \DIFaddbegin \DIFadd{mask QC exclusion used in masked evaluation. Downstream classification used a separate task-specific stratified three-fold protocol, described below. The classification analysis retained the existing author-provided synthetic PET dataset under the intended two-fold OOF design; its image-to-checkpoint correspondence to the reconstructed Full generation experiments could not be independently established. The classification results are therefore analyzed separately from the reconstructed generator comparisons. Fixed splits, seeds and implementation details are summarized in Table~\ref{tab:reproducibility}}\DIFaddend .

\subsection{Experimental Setup and Evaluation}
\subsubsection{Competing Methods}
To \DIFdelbegin \DIFdel{rigorously benchmark the performance of }\DIFdelend \DIFaddbegin \DIFadd{benchmark }\DIFaddend DIReCT$++$, \DIFdelbegin \DIFdel{we selected six representative baseline methods spanning diverse generative paradigms, including unsupervised translation, diffusion-based synthesis, and segmentation-guided generation. All models were trained and evaluated on identical data splits, with hyperparameters optimized to their reported best configurations to ensure a fair and reproducible comparison}\DIFdelend \DIFaddbegin \DIFadd{six historical baseline methods were adapted to PET synthesis as described below. Reconstructed comparisons additionally used Clinical Regression 3D U-Net, MRI-only Swin-UNETR, TracerOnly, TabularMLP and separately trained Single models. Their input conditions, training endpoints and implementation differences are retained; the comparisons are not treated as uniformly optimized, equal-budget or single-component ablations}\DIFaddend .

\begin{enumerate}
    \item \textbf{CycleGAN}~\cite{zhu2017unpaired}: A widely adopted unsupervised image-to-image translation framework that enforces cycle-consistency to learn bidirectional mappings between unpaired domains. We include it as a representative of non-probabilistic, GAN-based cross-modality synthesis.

    \item \textbf{Swin-UNet}~\cite{Cao2021swin}: A hybrid architecture combining Swin Transformer blocks with a U-Net encoder–decoder structure, excelling in high-resolution medical image generation and fine-grained anatomical modeling.

    \item \textbf{SegGuidedDiff}~\cite{konz2024anatomically}: A \DIFdelbegin \DIFdel{state-of-the-art diffusion model that leverages segmentation maps as anatomical priors to guide PET synthesis . It serves as a high-fidelity reference for structure-preserving generative performance}\DIFdelend \DIFaddbegin \DIFadd{segmentation-guided diffusion model adapted for 2D slice-based PET synthesis using anatomical segmentation maps as guidance}\DIFaddend .

    \item \textbf{3D DDIM}~\DIFdelbegin \DIFdel{\cite{Graf2023}}\DIFdelend \DIFaddbegin \DIFadd{\cite{song2021ddim,pinaya2023generative}}\DIFaddend : A deterministic fast-sampling variant of 3D diffusion models, \DIFdelbegin \DIFdel{included to assess DIReCT$++$’s computational efficiency and sampling quality relative to accelerated diffusion approaches}\DIFdelend \DIFaddbegin \DIFadd{implemented using a 3D diffusion-model U-Net and a DDIM scheduler for the historical synthesis comparison}\DIFaddend .

    \item \textbf{3D Rectified Flow (3D RF)}~\cite{Xing2022}: A flow-based generative model that shares DIReCT$++$’s underlying ODE framework but lacks domain-informed conditioning.
\DIFdelbegin \DIFdel{This ablation model isolates the contribution of our BiomedCLIP-enhanced guidance module.
}\DIFdelend 

    \item \textbf{DIReCT (Single-Modality Variant)}~\cite{liu2025direct}: \DIFdelbegin \DIFdel{A simplified version of our framework that synthesizes only }\DIFdelend \DIFaddbegin \DIFadd{Separate tracer-specific models synthesize }\DIFaddend one PET tracer \DIFdelbegin \DIFdel{type (either \textsuperscript{18}F-FDG or \textsuperscript{18}F-AV45) }\DIFdelend per forward pass\DIFdelbegin \DIFdel{. In contrast to }\DIFdelend \DIFaddbegin \DIFadd{, while }\DIFaddend DIReCT$++$ \DIFdelbegin \DIFdel{—which jointly generates both tracers through cross-modality interaction—this variant enables direct assessment of the benefits conferred by multi-tracer synergistic generation in terms of anatomical consistency and quantitative accuracy}\DIFdelend \DIFaddbegin \DIFadd{uses a shared network conditioned on the target tracer. The comparison includes other implementation differences and does not isolate joint-training synergy}\DIFaddend .
\end{enumerate}

\DIFaddbegin \DIFadd{The 3D Swin-UNETR was adapted to PET intensity regression with a continuous single-channel output and MSE loss, using MRI alone. Clinical Regression 3D U-Net used MRI and the same eight clinical variables with target-tracer information, providing a clinically matched framework comparison. TracerOnly excluded patient-specific conditioning while retaining target-tracer guidance; TabularMLP replaced the clinical text representation with a structured-variable encoder. Separate Single models were trained per tracer. The principal comparison endpoint was epoch 165, except the prespecified Swin-UNETR endpoint of epoch 200. These independently initialized training runs were not designed to estimate variability over training seeds or to isolate every module as a single-factor causal ablation.
}

\DIFaddend \subsection{Statistical Analysis}
\DIFaddbegin \subsubsection{\texorpdfstring{\DIFadd{Image-Quality Evaluation}}{Image-Quality Evaluation}}
\DIFadd{Whole-volume metrics include background; masked metrics use the frozen SynthSeg nonzero-label union, which includes ventricles and CSF and is not a pure brain-parenchymal mask. MSE/MAE average squared/absolute errors within that mask, and PSNR uses data range 1. Masked 3D SSIM uses a Gaussian $7\times7\times7$ kernel ($\sigma=1.5$) and averages only centers whose complete window lies inside the mask; outside-mask voxels were not set to zero and included in a whole-image SSIM average. The masked comparison uses identical target and mask arrays across models. The existing MRI/mask QC exclusion removes one subject per tracer, yielding FDG $n=1,597$ and AV45 $n=1,364$; this exclusion is specific to the masked analysis and case selection and is not imposed on the classification cohort.
}

\DIFadd{Metrics are reported as across-subject mean $\pm$ sample s.d. Paired two-sided $t$-tests and 10,000 paired bootstrap resamples stratified by held-out fold (seed 42) assess the frozen generation predictions. For each model comparison, Holm correction covers eight tests (two tracers $\times$ four metrics); bootstrap 95\% percentile CIs are unadjusted. Whole-volume and masked results use their respective subject sets and target protocols. These intervals do not estimate retraining or split variability. Supplemental cases are selected within each tracer's QC-passed cohort by saved whole-volume SSIM: nearest subject to the median or 10th percentile (linear quantiles), or the minimum for the worst case; ties are resolved by stable participant-ID ordering, without visual selection.
}\DIFaddend \subsubsection{Regional Group-Wise Comparison}

\DIFdelbegin \DIFdel{To evaluate the clinical validity of synthesized PET (Syn-PET) images, we conducted a region-of-interest (ROI) -based statistical analysis focused on canonical Alzheimer’sdisease (AD) -vulnerable brain regions. T1-weighted MRI scans were first parcellated into }\DIFdelend \DIFaddbegin \DIFadd{Regional analysis used }\DIFaddend 98 \DIFdelbegin \DIFdel{anatomical ROIs using SynthSeg ~\cite{synthseg}. Among these, we prioritized a set of AD-relevant regions : for metabolic assessment with \textsuperscript{18}F-FDG PET,}\DIFdelend \DIFaddbegin \DIFadd{SynthSeg ROIs defined on T1-weighted MRI. Representative regions shown in Fig.~\ref{fig:4}C,D were selected for disease-related comparisons. FDG regions included }\DIFaddend the bilateral hippocampus, \DIFdelbegin \DIFdel{temporoparietal junction, and posterior cingulate cortex ~\cite{mueller2005alzheimers}; and for amyloid burden evaluation with \textsuperscript{18}F-AV45 PET, the precuneus, frontal lobes, }\DIFdelend \DIFaddbegin \DIFadd{isthmus cingulate cortex }\DIFaddend and \DIFdelbegin \DIFdel{temporal lobes}\DIFdelend \DIFaddbegin \DIFadd{supramarginal gyrus (atlas labels \texttt{Left-Hippocampus}, \texttt{Right-Hippocampus}, and the \texttt{ctx-lh-}/}\allowbreak\DIFadd{\texttt{ctx-rh-} labels with suffixes \texttt{isthmuscingulate} and \texttt{supramarginal})~\cite{mosconi2005reduced,he2015meta,strom2022cortical}. AV45 regions included bilateral precuneus and inferior, middle, superior and transverse temporal gyri (the corresponding \texttt{ctx-lh-}/}\allowbreak\DIFadd{\texttt{ctx-rh-} labels)}\DIFaddend ~\cite{wong2010in}. \DIFaddbegin \DIFadd{Atlas-defined names are used rather than the ambiguous temporoparietal-junction label.
}\DIFaddend 

The validation comprised two complementary analyses:

\begin{enumerate}
    \item \textbf{Intra-group consistency}: \DIFdelbegin \DIFdel{For each diagnostic group (AD and CN), we }\DIFdelend \DIFaddbegin \DIFadd{We }\DIFaddend performed paired two-sided $t$-tests to compare \DIFdelbegin \DIFdel{the mean standardized uptake values (SUVs) }\DIFdelend \DIFaddbegin \DIFadd{normalized regional uptake }\DIFaddend between real and Syn-PET images within each \DIFdelbegin \DIFdel{key ROI}\DIFdelend \DIFaddbegin \DIFadd{ROI, both across the overall cohorts and within individual diagnostic groups}\DIFaddend . The null hypothesis was that the mean \DIFdelbegin \DIFdel{difference in SUV between real and synthesized scans is zero. A non-significant result ($p > 0.05$, corrected for multiple comparisons)would indicate high fidelity of Syn-PET to ground-truth PET at the group level}\DIFdelend \DIFaddbegin \DIFadd{paired difference in normalized regional uptake was zero. Mean paired differences with 95\% confidence intervals, mean absolute error (MAE), Lin's concordance correlation coefficient (CCC), the two-way absolute-agreement single-measure intraclass correlation coefficient (ICC(A,1)), and Bland--Altman bias and 95\% limits of agreement were also calculated to assess regional agreement. For the 16 representative tracer--ROI combinations shown in Fig.~\ref{fig:4}A, B, we additionally performed subject-level ordinary least-squares regression of synthetic on real normalized regional uptake in the overall cohorts, reporting slopes, intercepts, and their unadjusted 95\% confidence intervals}\DIFaddend .

    \item \textbf{Inter-group discriminability}: To assess whether Syn-PET preserves diagnostic signal, we conducted independent two-sided \DIFaddbegin \DIFadd{Student's }\DIFaddend $t$-tests \DIFdelbegin \DIFdel{comparing mean SUVs }\DIFdelend \DIFaddbegin \DIFadd{(assuming equal variances) comparing mean normalized regional uptake }\DIFaddend between the AD and CN groups in the same ROIs—first using real PET and then repeating the analysis with Syn-PET. \DIFdelbegin \DIFdel{Preservation of statistically significant group differences(e.g., hypometabolism in FDG, elevated amyloid in AV45) in Syn-PET data would demonstrate its capacity to recapitulate clinically meaningful biomarkers}\DIFdelend \DIFaddbegin \DIFadd{Cohen's $d$ was calculated separately for real and synthetic PET as the AD--CN mean difference divided by the pooled within-group sample standard deviation, to characterize the direction and magnitude of the group differences}\DIFaddend .
\end{enumerate}

All \DIFdelbegin \DIFdel{statistical tests }\DIFdelend \DIFaddbegin \DIFadd{$t$-tests }\DIFaddend were two-tailed, and $p$-values were adjusted for multiple ROI comparisons using the \DIFaddbegin \DIFadd{Benjamini--Hochberg }\DIFaddend false discovery rate (FDR) procedure~\cite{benjamini1995controlling}. \DIFdelbegin \DIFdel{This dual-validation framework ensures that Syn-PET not only approximates real PET quantitatively but also retains its discriminative power for early diagnostic stratification. }%DIFDELCMD < 

\par\noindent\DIFdel{\textbf{Diagnosis and Prognostic Stratification}}\par
%DIFDELCMD < %%%
\DIFdelend \DIFaddbegin \DIFadd{For the paired real--synthetic comparisons, FDR correction was applied across 98 ROIs separately for each tracer and analysis group (overall, AD, CN, EMCI, LMCI, and MCI), yielding 12 families (2 tracers $\times$ 6 analysis groups) and 1,176 tests in total. For the AD--CN comparisons, FDR correction was applied across 98 ROIs separately for each tracer and PET type (real or synthetic), yielding four families and 392 tests in total. Statistical significance was defined as an FDR-adjusted $q<0.05$.
}\DIFaddend 

\DIFdelbegin \DIFdel{To rigorously evaluate the clinical utility of synthesized PET (Syn-PET)images, we assessed their capacity to support diagnostic and prognostic classification across the AD continuum. Specifically, we trained identical DenseNet classifiers }\DIFdelend \DIFaddbegin \subsubsection{\texorpdfstring{\DIFadd{Diagnostic Classification and MCI Stage Discrimination}}{Diagnostic Classification and MCI Stage Discrimination}}
\DIFadd{Image-based classifiers used a DenseNet121 backbone}\DIFaddend ~\cite{huang2017densely} to \DIFdelbegin \DIFdel{discriminate between three clinically salient diagnostic pairs: (1)  ADversus CNindividuals, a canonical diagnostic boundary; (2)  MCI versus CN, representing early disease detection; and (3) EMCIversus late MCI (LMCI), a prognostic distinction reflecting differential progression risk. }%DIFDELCMD < 

%DIFDELCMD < %%%
\DIFdel{For each classification task, models were independently trained and evaluated using the following input modalities: (i) real\textsuperscript{18}F-FDG and/or \textsuperscript{18}F-AV45 PET; (ii) Syn-PET counterparts; (iii) structural MRIalone; and (iv)multi-modal combinations, including MRI\,+\,Syn-FDG, MRI\,}\DIFdelend \DIFaddbegin \DIFadd{evaluate three cross-sectional tasks: AD/CN; EMCI}\DIFaddend +\DIFdelbegin \DIFdel{\,Syn-AV45, }\DIFdelend \DIFaddbegin \DIFadd{LMCI/CN; }\DIFaddend and \DIFdelbegin \DIFdel{MRI\,}\DIFdelend \DIFaddbegin \DIFadd{EMCI/LMCI. Positive classes were AD, EMCI}\DIFaddend +\DIFdelbegin \DIFdel{\,Syn-FDG\,}\DIFdelend \DIFaddbegin \DIFadd{LMCI and LMCI, respectively. The task-specific cohorts contained 470, 741 and 436 participants, excluding MCI-U. Seven imaging conditions used MRI alone or MRI combined with real/synthetic FDG, AV45 or both tracers. Two further conditions were Clinical-only (an 18-parameter linear classifier) and MRI}\DIFaddend +\DIFdelbegin \DIFdel{\,Syn-AV45. All models were optimized on identical training protocols and validated via 3-fold cross-validation with stratified sampling to preserve class distribution across folds. Performance was quantified using accuracy (ACC), sensitivity (SEN}\DIFdelend \DIFaddbegin \DIFadd{Clinical. Both clinical baselines used the same eight clinical variables as the conditioning inputs; their information content is diagnosis-related and affected by historical group-conditioned filling.
}

\DIFadd{All conditions used the same task-specific subject-level stratified three-fold outer splits (seed 42). Within each outer training set, a fixed stratified 15\% validation subset selected the checkpoint by validation ACC; the outer test was evaluated once after reloading that checkpoint. Imaging models used Adam (learning rate $10^{-4}$}\DIFaddend ), \DIFdelbegin \DIFdel{and specificity (SPE )}\DIFdelend \DIFaddbegin \DIFadd{batch size 5, up to 200 epochs and validation every two epochs. Clinical and combined-model settings are retained separately in the reproducibility records rather than assumed identical to every imaging model. Thresholded metrics use 0.5; AUC uses continuous positive-class scores. We report ACC, AUC, SEN, SPE and balanced ACC as equally weighted means across the three folds with sample s.d}\DIFaddend .

\DIFdelbegin \DIFdel{This comparative framework enables a direct assessment of whether Syn-PET preserves the discriminative biomarker information encoded in real PET, and whether multi-tracer Syn-PET enhances early stratification beyond single-modality or MRI-only approaches. Crucially, equivalence in classification performance between real and Syn-PET would indicate that the synthesized images retain sufficient biological and pathological fidelityfor use in diagnostic pipelines.
}\DIFdelend \DIFaddbegin \DIFadd{Bootstrap CIs use 5,000 paired resamples within each fold-by-class stratum (seed 20261102); each condition shares participant multiplicities within a resample. Metrics are recalculated by fold and averaged equally, rather than pooling differently trained model scores. Paired differences use the same resamples. The 95\% percentile intervals are marginal and unadjusted for simultaneous coverage; they condition on the trained models, class proportions and fixed splits and do not estimate retraining variability or eliminate dependence from overlapping training sets. DeLong tests compare matched predictions separately within each outer fold. The 18 original imaging comparisons and 36 additional clinical comparisons form separate Holm families; raw $p$ values are retained. No overall three-fold DeLong $p$ value is reported. ROC curves average fold-specific TPR over an FPR grid with step 0.005; pointwise 95\% bands use 2,000 fold-by-class resamples. The legend AUC is the mean empirical fold AUC, not pooled AUC or an integral of the finite grid.
}

\DIFadd{Similar observed Real/Syn performance or a non-significant difference does not establish biological fidelity, equivalence or clinical interchangeability. Clinical-only and MRI+Clinical comparisons assess whether the reported synthetic-image classifier adds information beyond the provided clinical fields; they do not resolve label circularity introduced by diagnostic-group-conditioned simulated inputs.
}

\DIFadd{Demographic sensitivity analyses separately stratified recorded age at 75 years and sex (F/M), without crossing the strata or changing thresholds/checkpoints. Each task, fold, subgroup and true-class stratum was resampled 5,000 times (seed 20261103); seven imaging inputs shared resample multiplicities, and fold metrics were averaged equally. These unadjusted marginal subgroup CIs were not used to infer interactions. Recorded ages could not be verified as MRI-visit ages because the available demographic export lacked visit dates. Age and sex were not mutually adjusted, and clinical-score or residual within-group confounding was not eliminated.
}\DIFaddend \subsection{Implementation Details}

\DIFdelbegin \DIFdel{All models were implemented in PyTorch (v2.6.0)and trained on NVIDIA RTX A6000 GPUs (48 GB memory per card). The Adam optimizer was used with an initial learning rate of $2.5 \times 10^{-5}$. Training was conducted for }\DIFdelend \DIFaddbegin \DIFadd{The reconstructed generator is a four-level 3D U-Net with channels $(32,64,64,128)$, one residual block per level, cross-attention only at the deepest level, and conditioning dimension 512. It contains 9,925,921 trainable parameters, excluding the frozen BiomedCLIP encoders. Generator training used Adam (learning rate $2.5\times10^{-5}$), batch size 1 and AMP. The principal generation endpoint is epoch }\DIFaddend 165\DIFdelbegin \DIFdel{epochs, with knowledge distillation applied from epoch 140 onward. Automatic Mixed Precision (AMP) was employed to accelerate trainingand reduce memory consumption. Each model required approximately }\DIFdelend \DIFaddbegin \DIFadd{; documented epoch-200 outputs and the MRI-only Swin-UNETR endpoint are distinguished from that endpoint. Recorded training/split settings by experiment family are summarized in Table~\ref{tab:reproducibility}; hardware and software differed between historical, reconstructed and timing runs. The historical estimate of }\DIFaddend 138 \DIFdelbegin \DIFdel{GPU hours for full convergence . A complete list of software dependencies and their versions is provided in Table~\ref{tab:supp_software} for full reproducibility}\DIFdelend \DIFaddbegin \DIFadd{GPU-hours was not independently verified and is not used to claim full convergence or a speedup over other methods.
}

\DIFadd{Inference timing used the Full fold-1 epoch-165 generator on an RTX 4090 (24 GB), PyTorch 2.5.1+cu121, CUDA FP32 and GPU-resident $1\times1\times160\times192\times160$ MRI with cached text context. After ten warm-up iterations, 50 synchronized repetitions per tracer gave FDG $251.015\pm0.187$ ms and AV45 $251.411\pm0.256$ ms (mean $\pm$ sample s.d. across repeated measurements of one fixed input). Peak PyTorch-allocated memory was 8.29 GiB and reserved memory 13.27 GiB, excluding other processes. These measurements exclude text encoding, preprocessing, transfers and file I/O; they describe a single-tracer forward pass, not end-to-end clinical workflow latency. Validated full-network FLOP counts and identical-protocol baseline timing were unavailable, so no FLOP estimate or baseline speedup ratio is claimed}\DIFaddend .

\subsection{Ethical Statements}

This study used publicly available, de-identified data from the Alzheimer's Disease Neuroimaging Initiative (ADNI) and the Open Access Series of Imaging Studies 3 (OASIS-3). The original studies were approved by the institutional review boards of the participating institutions, and written informed consent was obtained from all participants or their legally authorized representatives. No additional participant recruitment or identifiable data collection was performed in this study.



\section*{Acknowledgments}

\subsection*{General}

The authors thank the participants and investigators of the Alzheimer's Disease Neuroimaging Initiative (ADNI) and the Open Access Series of Imaging Studies (OASIS-3) for providing the data used in this study.

\subsection*{Author Contributions}

S.L., S.Y., and T.L. contributed equally to this work. S.L. and T.L. conceived
the methodology, implemented the core algorithms, and conducted the main
experiments. S.Y. contributed to data curation, preprocessing, and experimental
implementation. H.W. contributed to model evaluation, validation experiments,
and statistical analysis. C.L. contributed to study design, project supervision,
and interpretation of experimental results. J.L. provided clinical guidance and
contributed to the interpretation of neuroimaging findings. J.M. and L.X.
provided methodological guidance, technical supervision, and critical revision
of the manuscript. S.L. and T.L. drafted the manuscript. All authors reviewed,
edited, and approved the final manuscript.

\subsection*{Funding}
The authors acknowledge the funding of the National Natural Science Foundation of China (No.~T2522028 and 12326616), Natural Science Basic Research Program of Shaanxi (No.~2024JC-TBZC-09), and Shaanxi Provincial Key Industrial Innovation Chain Project (No.~2024SF-ZDCYL-02-10).

\subsection*{Competing Interests}
The authors declare that they have no competing interests.

\subsection*{Data Availability}
The data used in this study were obtained from two \DIFdelbegin \DIFdel{publicly available }\DIFdelend \DIFaddbegin \DIFadd{research }\DIFaddend datasets: the Alzheimer's Disease Neuroimaging Initiative (ADNI) and the Open Access Series of Imaging Studies (OASIS-3). ADNI data are \DIFdelbegin \DIFdel{publicly accessible }\DIFdelend \DIFaddbegin \DIFadd{available to approved researchers }\DIFaddend through the LONI Image and Data Archive (IDA) at \url{https://ida.loni.usc.edu}\DIFdelbegin \DIFdel{after registration and agreement to the data use terms}\DIFdelend \DIFaddbegin \DIFadd{, subject to the ADNI Data Use Agreement}\DIFaddend . OASIS-3 data are available \DIFdelbegin \DIFdel{at }%DIFDELCMD < \url{https://www.oasis-brains.org} %%%
\DIFdel{under a Creative Commons Attribution 4.0 International License. The specific imaging and clinical data used in this work include structural MRI}\DIFdelend \DIFaddbegin \DIFadd{through the official OASIS data access process at }\url{https://sites.wustl.edu/oasisbrains/request-access/}\DIFadd{, subject to the corresponding data use terms. ADNI inputs included T1-weighted MRI, $^{18}$F-FDG PET, $^{18}$F-AV45 PET, and demographic and clinical assessments; external OASIS-3 evaluation used MRI and AV45-PET. Fixed subject-level split records and source manifests for the revised analyses were retained. Sharing of subject- }\DIFaddend and \DIFdelbegin \DIFdel{AV-45 PETscans, along with associated demographic and cognitive assessments}\DIFdelend \DIFaddbegin \DIFadd{image-level identifiers is subject to the applicable data access and redistribution requirements. Clinical-source identifiers must not be interpreted as verified acquisition identifiers of the actual model inputs}\DIFaddend .

\subsection*{Code Availability}
The training and inference scripts, and trained models have been publicly released. Please refer to \url{https://github.com/ladderlab-xjtu/DIReCT-PLUS} for more details.

\section*{Supplementary Materials}
\begingroup
\setcounter{figure}{0}
\renewcommand{\thefigure}{S\arabic{figure}}
\DIFaddbegin \renewcommand{\theHfigure}{S\arabic{figure}}
\DIFaddend \setcounter{table}{0}
\renewcommand{\thetable}{S\arabic{table}}
\DIFaddbegin \renewcommand{\theHtable}{S\arabic{table}}
\DIFaddend 

% BEGIN supplementary/supplementary
\subsection*{Supplementary Results}

\begin{figure}[htbp]
\centering
\includegraphics[width=\textwidth]{supplementary_figure_1}
\DIFdelbeginFL %DIFDELCMD < \caption{\textbf{Comparison of real and synthetic AV45 uptake across the remaining brain regions (Part 1 of 2).}}
%DIFDELCMD <     %%%
\DIFdelendFL \DIFaddbeginFL \caption{{\color{blue}\textbf{Comparison of real and synthetic AV45 uptake across the remaining brain regions (Part 1 of 2).} Values represent mean regional intensities from per-image min--max-normalized PET images (arbitrary units), rather than SUVR.}}
\DIFaddendFL \label{fig:supp_av45_1}
\end{figure}

\begin{figure}[htbp]
\centering
\includegraphics[width=\textwidth]{supplementary_figure_2}
\DIFdelbeginFL %DIFDELCMD < \caption{\textbf{Comparison of real and synthetic AV45 uptake across the remaining brain regions (Part 2 of 2).}}
%DIFDELCMD <     %%%
\DIFdelendFL \DIFaddbeginFL \caption{{\color{blue}\textbf{Comparison of real and synthetic AV45 uptake across the remaining brain regions (Part 2 of 2).} Values represent mean regional intensities from per-image min--max-normalized PET images (arbitrary units), rather than SUVR.}}
\DIFaddendFL \label{fig:supp_av45_2}
\end{figure}

\begin{figure}[htbp]
\centering
\includegraphics[width=\textwidth]{supplementary_figure_3}
\DIFdelbeginFL %DIFDELCMD < \caption{\textbf{Comparison of real and synthetic FDG uptake across the remaining brain regions (Part 1 of 2).}}
%DIFDELCMD <     %%%
\DIFdelendFL \DIFaddbeginFL \caption{{\color{blue}\textbf{Comparison of real and synthetic FDG uptake across the remaining brain regions (Part 1 of 2).} Values represent mean regional intensities from per-image min--max-normalized PET images (arbitrary units), rather than SUVR.}}
\DIFaddendFL \label{fig:supp_fdg_1}
\end{figure}

\begin{figure}[htbp]
\centering
\includegraphics[width=\textwidth]{supplementary_figure_4}
\DIFdelbeginFL %DIFDELCMD < \caption{\textbf{Comparison of real and synthetic FDG uptake across the remaining brain regions (Part 2 of 2).}}
%DIFDELCMD <     %%%
\DIFdelendFL \DIFaddbeginFL \caption{{\color{blue}\textbf{Comparison of real and synthetic FDG uptake across the remaining brain regions (Part 2 of 2).} Values represent mean regional intensities from per-image min--max-normalized PET images (arbitrary units), rather than SUVR.}}
\DIFaddendFL \label{fig:supp_fdg_2}
\end{figure}

\begin{figure}[htbp]
\centering
\includegraphics[width=\textwidth]{supplementary_figure_5}
\DIFdelbeginFL %DIFDELCMD < \caption{\textbf{Validation of disease-related AV45 uptake patterns across the remaining brain regions.} Violin plots display mean uptake in Alzheimer’s disease (AD, red) and cognitively normal (CN, green) subjects for selected brain regions. Each region contains paired half-violins: the left half (solid fill) represents real PET data, and the right half (dashed fill) represents MRI-synthetic PET data, with consistent color coding between modalities. Solid lines above the violins indicate statistical significance for real PET, while dashed lines indicate significance for synthetic PET.}
%DIFDELCMD <     %%%
\DIFdelendFL \DIFaddbeginFL \caption{{\color{blue}\textbf{Validation of disease-related AV45 uptake patterns across the remaining brain regions.} Violin plots display mean normalized regional uptake (arbitrary units) in Alzheimer’s disease (AD, red) and cognitively normal (CN, green) subjects for selected brain regions. Each region contains paired half-violins: the left half (solid fill) represents real PET data, and the right half (dashed fill) represents MRI-synthetic PET data, with consistent color coding between modalities. Solid lines above the violins indicate statistical significance for real PET, while dashed lines indicate significance for synthetic PET. Uptake values were derived from per-image min--max-normalized PET images without reference-region normalization.}}
\DIFaddendFL \label{fig:supp_disease_av45}
\end{figure}

\begin{figure}[htbp]
\centering
\includegraphics[width=0.75\textwidth]{supplementary_figure_6}
\DIFdelbeginFL %DIFDELCMD < \caption{\textbf{Validation of disease-related FDG uptake patterns across the remaining brain regions.} Violin plots display mean uptake in Alzheimer’s disease (AD, red) and cognitively normal (CN, green) subjects for selected brain regions. Each region contains paired half-violins: the left half (solid fill) represents real PET data, and the right half (dashed fill) represents MRI-synthetic PET data, with consistent color coding between modalities. Solid lines above the violins indicate statistical significance for real PET, while dashed lines indicate significance for synthetic PET.}
%DIFDELCMD <     %%%
\DIFdelendFL \DIFaddbeginFL \caption{{\color{blue}\textbf{Validation of disease-related FDG uptake patterns across the remaining brain regions.} Violin plots display mean normalized regional uptake (arbitrary units) in Alzheimer’s disease (AD, red) and cognitively normal (CN, green) subjects for selected brain regions. Each region contains paired half-violins: the left half (solid fill) represents real PET data, and the right half (dashed fill) represents MRI-synthetic PET data, with consistent color coding between modalities. Solid lines above the violins indicate statistical significance for real PET, while dashed lines indicate significance for synthetic PET. Uptake values were derived from per-image min--max-normalized PET images without reference-region normalization.}}
\DIFaddendFL \label{fig:supp_disease_fdg}
\end{figure}

\begin{table}[htbp]
\DIFdelbeginFL %DIFDELCMD < \caption{\textbf{Demographic characteristics of ADNI.} Data are mean $\pm$ s.d. unless otherwise indicated. $P$ values from one-way ANOVA (continuous) or chi-square test (sex). CDR: Clinical Dementia Rating; MMSE: Mini-Mental State Examination.}
%DIFDELCMD < %%%
\DIFdelendFL \DIFaddbeginFL \caption{{\color{blue}\textbf{Demographic characteristics and prepared clinical inputs of ADNI.} Data are mean $\pm$ sample s.d., except counts. Demographic $P$ values use one-way ANOVA (age, weight) or a chi-square test (sex). Clinical summaries combine observed and diagnostic-group-conditioned filled scores; their group-comparison $P$ values are omitted because the distributions are not independent clinical evidence. Source missingness is reported in Table~\ref{tab:clinical_missingness}. CDR: Clinical Dementia Rating; MMSE: Mini-Mental State Examination.}}
\DIFaddendFL \label{tab:supp_adni}
\centering
\DIFdelbeginFL %DIFDELCMD < \begin{tabularx}{\textwidth}{l *{5}{>{\centering\arraybackslash}X} >{\centering\arraybackslash}X}
%DIFDELCMD < \toprule
%DIFDELCMD <  & \textbf{CN} & \textbf{EMCI} & \textbf{LMCI} & \textbf{MCI} & \textbf{AD} & \textbf{P} \\
%DIFDELCMD < \midrule
%DIFDELCMD < $n$ & 577 & 291 & 152 & 520 & 317 & \\
%DIFDELCMD < Age (years) & 72.55 $\pm$ 6.65 & 71.62 $\pm$ 7.51 & 72.44 $\pm$ 7.69 & 73.20 $\pm$ 7.51 & 74.90 $\pm$ 7.99 & $<$0.001 \\
%DIFDELCMD < Weight (kg) & 73.20 $\pm$ 24.13 & 80.77 $\pm$ 16.67 & 77.15 $\pm$ 16.64 & 72.90 $\pm$ 22.33 & 73.21 $\pm$ 18.64 & 0.0023 \\
%DIFDELCMD < Sex (M/F) & 252/325 & 165/126 & 80/72 & 310/210 & 186/131 & $<$0.001 \\
%DIFDELCMD < MMSE & 29.40 $\pm$ 0.70 & 27.50 $\pm$ 0.50 & 24.47 $\pm$ 0.50 & 26.83 $\pm$ 1.24 & 20.84 $\pm$ 1.97 & $<$0.001 \\
%DIFDELCMD < GDSCALE & 0.24 $\pm$ 0.73 & 1.01 $\pm$ 0.14 & 3.02 $\pm$ 0.28 & 1.81 $\pm$ 0.99 & 4.14 $\pm$ 1.62 & $<$0.001 \\
%DIFDELCMD < Global CDR & 0.00 $\pm$ 0.02 & 0.15 $\pm$ 0.23 & 0.50 $\pm$ 0.01 & 0.42 $\pm$ 0.19 & 0.95 $\pm$ 0.15 & $<$0.001 \\
%DIFDELCMD < FAQ & 0.52 $\pm$ 0.53 & 3.11 $\pm$ 2.42 & 8.83 $\pm$ 1.54 & 5.33 $\pm$ 1.51 & 20.20 $\pm$ 2.92 & $<$0.001 \\
%DIFDELCMD < NPI-Q & 0.01 $\pm$ 0.19 & 1.55 $\pm$ 2.17 & 4.00 $\pm$ 0.03 & 2.92 $\pm$ 0.50 & 4.98 $\pm$ 0.65 & $<$0.001 \\
%DIFDELCMD < \bottomrule
%DIFDELCMD < \end{tabularx}
%DIFDELCMD < %%%
\DIFdelendFL \DIFaddbeginFL \begin{tabularx}{\textwidth}{l *{5}{>{\centering\arraybackslash}X} >{\centering\arraybackslash}X}
\toprule
 & \textbf{CN} & \textbf{EMCI} & \textbf{LMCI} & \textbf{MCI} & \textbf{AD} & \textbf{P} \\
\midrule
$n$ & 577 & 291 & 152 & 520 & 317 & \\
Age (years) & 72.55 $\pm$ 6.65 & 71.62 $\pm$ 7.51 & 72.44 $\pm$ 7.69 & 73.20 $\pm$ 7.51 & 74.90 $\pm$ 7.99 & $<$0.001 \\
Weight (kg) & 73.20 $\pm$ 24.13 & 80.77 $\pm$ 16.67 & 77.15 $\pm$ 16.64 & 72.90 $\pm$ 22.33 & 73.21 $\pm$ 18.64 & 0.0023 \\
Sex (M/F) & 252/325 & 165/126 & 80/72 & 310/210 & 186/131 & $<$0.001 \\
MMSE & 29.40 $\pm$ 0.70 & 27.50 $\pm$ 0.50 & 24.47 $\pm$ 0.50 & 26.83 $\pm$ 1.24 & 20.84 $\pm$ 1.97 & -- \\
GDSCALE & 0.24 $\pm$ 0.73 & 1.00 $\pm$ 0.00 & 3.00 $\pm$ 0.00 & 1.81 $\pm$ 0.99 & 4.14 $\pm$ 1.62 & -- \\
Global CDR & 0.00 $\pm$ 0.02 & 0.15 $\pm$ 0.23 & 0.50 $\pm$ 0.00 & 0.42 $\pm$ 0.19 & 0.95 $\pm$ 0.15 & -- \\
FAQ & 0.52 $\pm$ 0.53 & 3.11 $\pm$ 2.42 & 8.83 $\pm$ 1.54 & 5.33 $\pm$ 1.51 & 20.20 $\pm$ 2.92 & -- \\
NPI-Q & 0.01 $\pm$ 0.19 & 1.55 $\pm$ 2.17 & 4.00 $\pm$ 0.00 & 2.92 $\pm$ 0.50 & 4.98 $\pm$ 0.65 & -- \\
\bottomrule
\end{tabularx}
\DIFaddendFL 

\vspace{2mm}
\footnotesize
\textit{Note:} CN = Cognitively Normal; EMCI = Early MCI; LMCI = Late MCI; MCI = \DIFdelbeginFL \DIFdelFL{Mild Cognitive Impairment}\DIFdelendFL \DIFaddbeginFL \DIFaddFL{the remaining MCI participants not assigned to EMCI or LMCI (MCI-U in the text)}\DIFaddendFL ; AD = Alzheimer’s dementia. \DIFaddbeginFL \DIFaddFL{Zero standard deviations reflect constant prepared scores, including diagnosis-conditioned fills; they do not establish absence of clinical variability. The original CN weight entries span 0--227 kg and are retained without exclusions or replacements; their clinical validity and the meaning of zero entries remain unresolved.
}\DIFaddendFL \end{table}

\DIFdelbegin \DIFdel{To obtain the required text-guided information for our analysis, we performedsecondary screening within the }\DIFdelend \DIFaddbegin \DIFadd{The historical }\DIFaddend OASIS-3 \DIFdelbegin \DIFdel{cohort (Open Access Series of Imaging Studies) , a longitudinal neuroimaging, clinical, and cognitive dataset for normalaging and Alzheimer’s disease(available at: }%DIFDELCMD < \url{https://www.medrxiv.org/content/10.1101/2019.12.13.19014902v1.full.pdf}%%%
\DIFdel{) . We focused on individuals with amyloid-beta PET imaging data processed through the standardized Pet Unified Pipeline (PUP). Specifically,we selected }\DIFdelend \DIFaddbegin \DIFadd{external cohort comprised }\DIFaddend 419 \DIFdelbegin \DIFdel{participants who had undergone AV45 PET scans and were PUP-processed,ensuring high-quality, quantitatively reliable neuroimaging data. These individuals formed the basis of our subsequent multimodal analysis. In contrast, FDG-PET data from a smaller subset ($n = 117$) were not processed with the PUP pipeline and were therefore excludedto maintain methodological consistency.
The demographic characteristics of this selected cohort, including age and gender distribution, are detailed in Table~\ref{tab:supp_oasis}. }\DIFdelend \DIFaddbegin \DIFadd{PUP-processed AV45 participants (334 CN, 85 AD); the 117 FDG participants without the corresponding processing were not included. The age distributions differed between diagnostic groups (Table~\ref{tab:supp_oasis}). Reported clinical summaries do not establish which complete prompt fields were used in every external run or how unavailable fields were handled. The external results are restricted to AV45 synthesis in CN/AD and do not provide FDG or MCI external evidence.
}\DIFaddend 

\begin{table}[htbp]
\caption{\textbf{Demographic characteristics of OASIS-3.}}
\label{tab:supp_oasis}
\centering
\begin{tabularx}{\textwidth}{l *{3}{>{\centering\arraybackslash}X}}
\toprule
 & \textbf{CN} & \textbf{AD} & \textbf{P} \\
\midrule
$n$ & 334 & 85 & \\
Age (years) & 66.33 $\pm$ 8.38 & 72.76 $\pm$ 6.96 & $<$0.001 \\
Sex (M/F) & 143/191 & 37/48 & $>$0.05 \\
MMSE score & 29.24 $\pm$ 0.90 & 23.21 $\pm$ 3.51 & $<$0.001 \\
Global CDR & 0.01 $\pm$ 0.08 & 0.73 $\pm$ 0.59 & $<$0.001 \\
\bottomrule
\end{tabularx}

\vspace{2mm}
\footnotesize
\textit{Note:} CN = \DIFdelbeginFL \DIFdelFL{Cognitively Normal; AD }\DIFdelendFL \DIFaddbeginFL \DIFaddFL{Cognitively Normal; AD }\DIFaddendFL = Alzheimer’s dementia.
\end{table}

\begin{table}[htbp]
\DIFdelbeginFL %DIFDELCMD < \caption{\textbf{Core software dependencies and versions for DIReCT$++$.} Only packages directly relevant to model implementation, training, and evaluation are listed. Full environment details are available in the code repository.}
%DIFDELCMD < %%%
\DIFdelendFL \DIFaddbeginFL \caption{{\color{blue}\textbf{Historical core software dependencies and reported versions for DIReCT$++$.} This retained environment list describes the historical configuration; reconstructed training and the timing run used separately recorded environments. In particular, the RTX 4090 timing run used PyTorch 2.5.1+cu121, as stated in Methods. The table is not a single environment specification for every experiment.}}
\DIFaddendFL \label{tab:supp_software}
\centering
\begin{tabularx}{\textwidth}{lX lX}
\toprule
\textbf{Package} & \textbf{Version} & \textbf{Package} & \textbf{Version} \\
\midrule
Python & 3.10 & PyTorch & 2.6.0+cu118 \\
torchvision & 0.21.0+cu118 & torchaudio & 2.6.0+cu118 \\
CUDA & 11.8 & cuDNN & 9.1.0 \\
einops & 0.8.1 & transformers & 4.40.2 \\
tokenizers & 0.19.1 & safetensors & 0.5.3 \\
xformers & 0.0.29.post3 & faiss & 1.9.0 \\
monai & 1.5.1 & numpy & 2.2.4 \\
scipy & 1.15.2 & pandas & 2.2.3 \\
scikit-learn & 1.6.1 & matplotlib & 3.10.1 \\
seaborn & 0.13.2 & opencv-python & 4.12.0.88 \\
nibabel & 5.3.2 & tensorboard & 2.19.0 \\
fsspec & 2024.6.1 & huggingface-hub & 0.30.1 \\
protobuf & 6.30.2 & triton & 3.2.0 \\
\bottomrule
\end{tabularx}
\end{table}
\DIFaddbegin 


%DIF >  BEGIN supplementary/tables/clinical_missingness
\begin{table}[htbp]
\caption{{\color{blue}\textbf{Missing clinical scores in the source records before historical filling.} Each cell gives missing count / group size (percentage); missingness is defined relative to the first source CSV record selected for each participant, not all ADNI visits.}}
\label{tab:clinical_missingness}
\centering\small
\begin{tabularx}{\textwidth}{l*{5}{>{\centering\arraybackslash}X}}
\toprule
Score & CN & EMCI & LMCI & MCI-U & AD \\
\midrule
MMSE & 427/577 (74.0\%) & 291/291 (100.0\%) & 152/152 (100.0\%) & 289/520 (55.6\%) & 234/317 (73.8\%) \\
GDSCALE & 429/577 (74.4\%) & 291/291 (100.0\%) & 152/152 (100.0\%) & 289/520 (55.6\%) & 234/317 (73.8\%) \\
Global CDR & 427/577 (74.0\%) & 291/291 (100.0\%) & 152/152 (100.0\%) & 289/520 (55.6\%) & 234/317 (73.8\%) \\
FAQ & 561/577 (97.2\%) & 167/291 (57.4\%) & 151/152 (99.3\%) & 489/520 (94.0\%) & 306/317 (96.5\%) \\
NPI-Q & 563/577 (97.6\%) & 167/291 (57.4\%) & 152/152 (100.0\%) & 491/520 (94.4\%) & 306/317 (96.5\%) \\
\bottomrule
\end{tabularx}
\par\smallskip\footnotesize
\DIFaddFL{Source records contain 1,751 observed and 7,534 missing score cells across 1,857 participants and five scores. All available scores were retained. Historical filling used diagnostic-group rules: MMSE values were drawn from group-specific integer sets; GDSCALE and NPI-Q used group-specific fixed values; FAQ used fixed values or group-specific integer sets/ranges; and CDR used fixed values or a group-specific Bernoulli choice for EMCI/MCI-U. These are simulated inputs, not clinical observations or group-mean imputation. Their use can encode diagnostic-group information. The original random states were not recovered, and no re-imputation was performed. Age, sex and weight have no blank entries in the prepared table, which does not establish the clinical validity of every entry. The clinical source row's Image ID does not identify the actual MRI/PET acquisition used by the model.
}\end{table}

%DIF >  END supplementary/tables/clinical_missingness


%DIF >  BEGIN supplementary/tables/reproducibility
\begin{table}[htbp]
\caption{{\color{blue}\textbf{Reproducibility settings and scope of the revised analyses.}}}
\label{tab:reproducibility}
\centering\small
\begin{tabularx}{\textwidth}{p{0.25\textwidth}X}
\toprule
Component & Recorded settings and limitations \\
\midrule
Reconstructed generation & 1,857 subjects; fixed two-fold partitions of 929/928, split seed 42. Full, TracerOnly, TabularMLP and Single did not use a fixed training seed; deterministic repeated-training behavior is not claimed. Epoch 165 is the principal comparison endpoint; MRI-only Swin-UNETR uses its prespecified epoch 200. \\
Classification & Task-specific subject-level stratified three-fold splits (seed 42); 15\% stratified internal validation within each outer training fold. Fit seed $42+100\times\mathrm{task\ index}+\mathrm{fold}$, with task index AD/CN=0, MCI/CN=1, EMCI/LMCI=2 and folds 1--3. Training lasts 200 epochs with validation every two epochs; checkpoint selected by validation ACC; outer test evaluated once after reloading. \\
Classification cohorts & AD/CN: 470 (165 AD, 305 CN); MCI/CN: 741 (436 EMCI+LMCI, 305 CN); EMCI/LMCI: 436 (287 EMCI, 149 LMCI). MCI-U excluded. The task-specific partitions are shared across input conditions. \\
Source identity & The historical classification Syn-PET inputs are distinct from the reconstructed generation comparisons; a direct checkpoint binding between these image families was not established. Existing classification inputs were not regenerated. \\
Clinical source & Clinical rows selected by first CSV record per participant; five score fields contain diagnosis-conditioned fills. Original filling random states and complete actual acquisition mappings were not recovered. \\
\bottomrule
\end{tabularx}
\end{table}

%DIF >  END supplementary/tables/reproducibility


%DIF >  BEGIN supplementary/tables/conditioning
\subsection*{\DIFadd{Supplementary Methods: conditioning prompts and adapter analysis}}
\DIFadd{The following is the historical clinical prompt template, with bracketed field placeholders instead of participant identifiers or invented measurements. Values use the historical formatting (rounded age and scores, one decimal for weight and CDR). Its use does not validate the filled values as clinical observations.
}\begin{quote}\small
\DIFadd{Subject is a }[\DIFadd{age}]\DIFadd{-year-old }[\DIFadd{sex}] \DIFadd{with a weight of }[\DIFadd{weight}] \DIFadd{kg. The global Clinical Dementia Rating (CDR) score, which assesses dementia severity (0: no dementia to 3: severe dementia), is }[\DIFadd{global\_cdr}]\DIFadd{. The Mini-Mental State Examination (MMSE) score, assessing cognitive function (0: severe impairment to 30: normal), is }[\DIFadd{mmse\_score}]\DIFadd{. The Geriatric Depression Scale (GDS) score, screening depression (0: no depression to 15: severe depression), is }[\DIFadd{gds\_score}]\DIFadd{. The Functional Activities Questionnaire (FAQ) score, assessing daily activity impairment (0: no impairment to 30: severe impairment), is }[\DIFadd{faq\_score}]\DIFadd{. The Neuropsychiatric Inventory Questionnaire (NPI-Q) Total Score, assessing neuropsychiatric symptom burden (0: no symptoms to higher scores indicating greater burden), is }[\DIFadd{npi\_q\_score}]\DIFadd{.
}\end{quote}
\noindent\DIFadd{\textbf{FDG tracer prompt.}
}\begin{quote}\small
\DIFadd{FDG PET is a functional brain imaging technique that visualizes the dynamic changes in glucose metabolism, directly linked to neuronal energy demands and synaptic activity. It serves as a tool to assess functional connectivity and energy utilization across brain regions. Areas with decreased metabolic activity, such as those affected by neurodegenerative diseases, should exhibit reduced signal intensity. High-intensity metabolic hotspots in gray matter (e.g., the cerebral cortex and basal ganglia) are key markers of neuronal activity. 
}\end{quote}
\noindent\DIFadd{\textbf{AV45 tracer prompt.}
}\begin{quote}\small
\DIFadd{AV45 PET is a molecular imaging technique that highlights the static distribution of amyloid-beta plaques, a critical pathological marker of Alzheimer's disease. This imaging modality provides a spatial map of amyloid deposition in cortical regions (e.g., the temporal, parietal, and frontal lobes) and can distinguish amyloid-positive areas from amyloid-negative white matter regions. The primary focus is on identifying amyloid deposition patterns to assess disease progression and pathological burden.
}\end{quote}
\DIFadd{The tracer texts are fixed model inputs; their statements about pathology or progression are not evidence that synthetic PET establishes a biological diagnosis or predicts progression.
}

\begin{table}[htbp]
\caption{{\color{blue}\textbf{Adapter-level sensitivity analysis of the cross-tracer soft target.} Descriptive means across two folds, using 1,856 QC-filtered participants. Only training-fold PET features were used for fitting; held-out features were evaluated after optimization. No cross-term means removal of the penalty ($\lambda_{\mathrm{cross}}=0$), not a target of zero. No generator was retrained or evaluated in this analysis.}}
\label{tab:supp_tau_adapter}
\centering\small
\begin{tabular}{lccccc}
\toprule
Setting & \multicolumn{2}{c}{Training alignment} & \multicolumn{2}{c}{Held-out alignment} & Cross-tracer \\
 & FDG & AV45 & FDG & AV45 & similarity \\
\midrule
$\tau=0.70$ & 0.9437 & 0.9410 & 0.9436 & 0.9410 & 0.7032 \\
$\tau=0.82$ & 0.9687 & 0.9661 & 0.9687 & 0.9661 & 0.8214 \\
$\tau=0.90$ & 0.9838 & 0.9813 & 0.9838 & 0.9812 & 0.9018 \\
$\tau=0.98$ & 0.9948 & 0.9923 & 0.9948 & 0.9923 & 0.9807 \\
No cross-term & 0.9952 & 0.9927 & 0.9952 & 0.9926 & 0.9897 \\
\bottomrule
\end{tabular}
\end{table}

%DIF >  END supplementary/tables/conditioning


%DIF >  BEGIN supplementary/tables/regional_agreement
\begin{landscape}
\begin{table}[p]
\caption{{\color{blue}\textbf{Regional normalized uptake agreement and descriptive regression for the 16 representative tracer--ROI combinations in Figure~\ref{fig:4}A,B.} Overall cohorts from the author-designated historical ROI CSVs; N can vary by ROI.}}
\label{tab:regional_agreement}
\centering\scriptsize\setlength{\tabcolsep}{2pt}
\resizebox{\linewidth}{!}{\begin{tabular}{llrccccccccc}
\toprule
Tracer & ROI (ID) & N & Mean diff. [95\% CI] & MAE & CCC & ICC(A,1) & BA bias [95\% LoA] & $p$ & $q$ & OLS slope [95\% CI] & Intercept [95\% CI] \\
\midrule
AV45 & L precuneus (1025) & 1365 & \shortstack{0.00052\\{}[-0.00056, 0.00159]} & 0.0163 & 0.9863 & 0.9863 & \shortstack{0.0005\\{}[-0.0391, 0.0401]} & 0.3445 & 0.9786 & \shortstack{1.0127\\{}[1.0038, 1.0215]} & \shortstack{-0.0064\\{}[-0.0114, -0.0014]} \\
AV45 & R precuneus (2025) & 1365 & \shortstack{0.00021\\{}[-0.00079, 0.00120]} & 0.0150 & 0.9887 & 0.9887 & \shortstack{0.0002\\{}[-0.0366, 0.0370]} & 0.6846 & 0.9786 & \shortstack{1.0008\\{}[0.9928, 1.0088]} & \shortstack{-0.0002\\{}[-0.0047, 0.0043]} \\
AV45 & L frontalpole (1032) & 1365 & \shortstack{0.00033\\{}[-0.00072, 0.00138]} & 0.0157 & 0.9808 & 0.9808 & \shortstack{0.0003\\{}[-0.0385, 0.0392]} & 0.5384 & 0.9786 & \shortstack{1.0037\\{}[0.9932, 1.0143]} & \shortstack{-0.0012\\{}[-0.0055, 0.0032]} \\
AV45 & R frontalpole (2032) & 1364 & \shortstack{0.00039\\{}[-0.00068, 0.00147]} & 0.0160 & 0.9815 & 0.9815 & \shortstack{0.0004\\{}[-0.0392, 0.0400]} & 0.4746 & 0.9786 & \shortstack{1.0007\\{}[0.9904, 1.0110]} & \shortstack{0.0001\\{}[-0.0042, 0.0045]} \\
AV45 & L temporalpole (1033) & 1365 & \shortstack{-0.00009\\{}[-0.00117, 0.00099]} & 0.0162 & 0.9591 & 0.9591 & \shortstack{-0.0001\\{}[-0.0399, 0.0398]} & 0.8769 & 0.9786 & \shortstack{0.9794\\{}[0.9640, 0.9947]} & \shortstack{0.0078\\{}[0.0019, 0.0137]} \\
AV45 & R temporalpole (2033) & 1365 & \shortstack{0.00031\\{}[-0.00076, 0.00138]} & 0.0163 & 0.9644 & 0.9644 & \shortstack{0.0003\\{}[-0.0392, 0.0399]} & 0.5667 & 0.9786 & \shortstack{1.0087\\{}[0.9942, 1.0232]} & \shortstack{-0.0032\\{}[-0.0092, 0.0027]} \\
AV45 & L superiorfrontal (1028) & 1365 & \shortstack{-0.00037\\{}[-0.00146, 0.00073]} & 0.0165 & 0.9806 & 0.9807 & \shortstack{-0.0004\\{}[-0.0407, 0.0400]} & 0.5102 & 0.9786 & \shortstack{1.0002\\{}[0.9896, 1.0107]} & \shortstack{-0.0005\\{}[-0.0058, 0.0049]} \\
AV45 & R superiorfrontal (2028) & 1365 & \shortstack{0.00063\\{}[-0.00041, 0.00168]} & 0.0157 & 0.9822 & 0.9822 & \shortstack{0.0006\\{}[-0.0380, 0.0393]} & 0.2354 & 0.9786 & \shortstack{0.9970\\{}[0.9869, 1.0071]} & \shortstack{0.0021\\{}[-0.0030, 0.0073]} \\
FDG & L precuneus (1025) & 1598 & \shortstack{-0.00003\\{}[-0.00100, 0.00093]} & 0.0156 & 0.9573 & 0.9573 & \shortstack{-0.0000\\{}[-0.0386, 0.0385]} & 0.9433 & 0.9774 & \shortstack{0.9844\\{}[0.9699, 0.9990]} & \shortstack{0.0108\\{}[0.0006, 0.0209]} \\
FDG & R precuneus (2025) & 1598 & \shortstack{0.00010\\{}[-0.00086, 0.00105]} & 0.0156 & 0.9598 & 0.9598 & \shortstack{0.0001\\{}[-0.0380, 0.0382]} & 0.8428 & 0.9774 & \shortstack{0.9983\\{}[0.9841, 1.0125]} & \shortstack{0.0013\\{}[-0.0086, 0.0111]} \\
FDG & L frontalpole (1032) & 1598 & \shortstack{0.00024\\{}[-0.00074, 0.00122]} & 0.0160 & 0.9586 & 0.9586 & \shortstack{0.0002\\{}[-0.0389, 0.0394]} & 0.6273 & 0.9410 & \shortstack{1.0039\\{}[0.9895, 1.0184]} & \shortstack{-0.0019\\{}[-0.0100, 0.0061]} \\
FDG & R frontalpole (2032) & 1597 & \shortstack{-0.00032\\{}[-0.00130, 0.00065]} & 0.0159 & 0.9634 & 0.9634 & \shortstack{-0.0003\\{}[-0.0394, 0.0387]} & 0.5160 & 0.9410 & \shortstack{1.0041\\{}[0.9906, 1.0177]} & \shortstack{-0.0026\\{}[-0.0103, 0.0050]} \\
FDG & L temporalpole (1033) & 1598 & \shortstack{-0.00034\\{}[-0.00132, 0.00063]} & 0.0159 & 0.9430 & 0.9430 & \shortstack{-0.0003\\{}[-0.0392, 0.0385]} & 0.4901 & 0.9410 & \shortstack{0.9934\\{}[0.9764, 1.0105]} & \shortstack{0.0023\\{}[-0.0046, 0.0091]} \\
FDG & R temporalpole (2033) & 1598 & \shortstack{0.00068\\{}[-0.00030, 0.00166]} & 0.0161 & 0.9449 & 0.9449 & \shortstack{0.0007\\{}[-0.0386, 0.0400]} & 0.1754 & 0.9410 & \shortstack{0.9902\\{}[0.9736, 1.0069]} & \shortstack{0.0047\\{}[-0.0022, 0.0116]} \\
FDG & L superiorfrontal (1028) & 1598 & \shortstack{0.00093\\{}[-0.00006, 0.00191]} & 0.0159 & 0.9438 & 0.9438 & \shortstack{0.0009\\{}[-0.0385, 0.0403]} & 0.0657 & 0.9410 & \shortstack{1.0101\\{}[0.9931, 1.0271]} & \shortstack{-0.0052\\{}[-0.0156, 0.0051]} \\
FDG & R superiorfrontal (2028) & 1598 & \shortstack{0.00081\\{}[-0.00017, 0.00178]} & 0.0158 & 0.9465 & 0.9465 & \shortstack{0.0008\\{}[-0.0382, 0.0398]} & 0.1055 & 0.9410 & \shortstack{1.0020\\{}[0.9855, 1.0185]} & \shortstack{-0.0004\\{}[-0.0104, 0.0096]} \\
\bottomrule
\end{tabular}}
\par\smallskip\footnotesize
\DIFaddFL{Differences and Bland--Altman (BA) bias are synthetic minus real uptake (arbitrary units after per-image min--max normalization). Mean-difference confidence intervals and individual BA limits of agreement are distinct. CCC is Lin's concordance correlation and ICC(A,1) is a two-way absolute-agreement single-measure coefficient. Paired two-sided $t$-test $p$ values are raw; BH-FDR $q$ values use all 98 ROIs within each tracer's overall cohort. OLS regresses synthetic on real uptake and gives unadjusted descriptive 95\% CIs. Non-significant differences and slopes close to one do not establish equivalence, clinical SUV/SUVR calibration or interchangeable biological measurements.
}\end{table}
\end{landscape}

%DIF >  END supplementary/tables/regional_agreement


%DIF >  BEGIN supplementary/tables/regional_effect_sizes
\begin{table}[p]
\caption{{\color{blue}\textbf{AD--CN regional effect sizes and BH-FDR-adjusted comparisons for Figure~\ref{fig:4}C--F.} Cohen's $d$ is AD minus CN divided by the pooled within-group sample s.d.}}
\label{tab:regional_effect_sizes}
\centering\scriptsize\setlength{\tabcolsep}{3pt}
\begin{tabularx}{\textwidth}{llXrrrrrr}
\toprule
Panel & Tracer & ROI (ID) & $N_{AD}$ & $N_{CN}$ & Real $d$ & Syn $d$ & Real $q$ & Syn $q$ \\
\midrule
C & AV45 & L precuneus (1025) & 179 & 511 & 1.6783 & 1.6437 & $7.73\times10^{-66}$ & $1.25\times10^{-63}$ \\
C & AV45 & L inferiortemporal (1009) & 179 & 511 & 1.6831 & 1.6337 & $4.68\times10^{-66}$ & $4.74\times10^{-63}$ \\
C & AV45 & L middletemporal (1015) & 179 & 511 & 1.6482 & 1.6236 & $4.67\times10^{-64}$ & $1.72\times10^{-62}$ \\
C & AV45 & L superiortemporal (1030) & 179 & 511 & 1.5780 & 1.5434 & $5.64\times10^{-60}$ & $8.51\times10^{-58}$ \\
C & AV45 & L transversetemporal (1034) & 179 & 511 & 1.5133 & 1.5288 & $2.84\times10^{-56}$ & $5.61\times10^{-57}$ \\
D & FDG & L isthmuscingulate (1010) & 303 & 371 & -1.2531 & -1.1774 & $2.69\times10^{-48}$ & $2.10\times10^{-43}$ \\
D & FDG & R isthmuscingulate (2010) & 303 & 371 & -1.3972 & -1.3201 & $1.18\times10^{-57}$ & $1.94\times10^{-52}$ \\
D & FDG & L supramarginal (1031) & 303 & 371 & -0.4848 & -0.4308 & $2.14\times10^{-9}$ & $1.10\times10^{-7}$ \\
D & FDG & Left Hippocampus (17) & 302 & 371 & -0.6401 & -0.6042 & $4.27\times10^{-15}$ & $1.27\times10^{-13}$ \\
D & FDG & Right Hippocampus (53) & 303 & 371 & -0.7518 & -0.7189 & $4.59\times10^{-20}$ & $1.54\times10^{-18}$ \\
E & AV45 & Left VentralDC (28) & 179 & 511 & -0.1050 & -0.0933 & 0.2368 & 0.2858 \\
E & AV45 & Right VentralDC (60) & 179 & 511 & -0.1211 & -0.0986 & 0.1729 & 0.2618 \\
E & AV45 & Right Hippocampus (53) & 179 & 511 & -0.0068 & 0.0006 & 0.9373 & 0.9949 \\
E & AV45 & Left Cerebellum White Matter (7) & 179 & 511 & -0.0745 & -0.1010 & 0.3996 & 0.2529 \\
E & AV45 & Left Hippocampus (17) & 179 & 511 & 0.0701 & 0.1232 & 0.4241 & 0.1667 \\
F & FDG & Left Thalamus (10) & 303 & 371 & -0.0097 & 0.0582 & 0.9008 & 0.4667 \\
F & FDG & Right Thalamus (49) & 303 & 371 & -0.0668 & -0.0530 & 0.3964 & 0.4991 \\
F & FDG & L parsorbitalis (1019) & 303 & 371 & -0.1294 & -0.0913 & 0.1035 & 0.2506 \\
F & FDG & L lateraloccipital (1011) & 303 & 371 & -0.1239 & -0.0978 & 0.1171 & 0.2229 \\
F & FDG & R lingual (2013) & 303 & 371 & -0.0493 & -0.0435 & 0.5303 & 0.5744 \\
\bottomrule
\end{tabularx}
\par\smallskip\footnotesize
\DIFaddFL{Separate two-sided Student's $t$-tests compare AD and CN for each tracer and PET type. BH-FDR correction uses all 98 ROIs within each of four tracer-by-PET-type families, rather than just the displayed regions. Positive $d$ indicates higher uptake in AD. All values are derived from min--max-normalized regional intensities, not SUV/SUVR. Preserved significance or similar effect sizes do not establish clinical equivalence.
}\end{table}

%DIF >  END supplementary/tables/regional_effect_sizes


%DIF >  BEGIN supplementary/tables/whole_volume_generation
%DIF >  Source-backed publication fragment; source records retained independently.
%DIF >  Source: G3/distribution_summary.csv
\begin{landscape}
\begingroup
\footnotesize
\setlength{\tabcolsep}{4pt}
\begin{longtable}{lllllll}
\caption{{\color{blue}Whole-volume OOF generation metrics (mean $\pm$ sample s.d.)}}\label{tab:whole_volume_generation}\\
\toprule
Model & Tracer & N & SSIM & PSNR & MSE & MAE \\
\midrule
\endfirsthead
\toprule
Model & Tracer & N & SSIM & PSNR & MSE & MAE \\
\midrule
\endhead
\bottomrule
\endfoot
Full & AV45 & 1365 & 0.884408 \(\pm\) 0.0405641 & 24.9664 \(\pm\) 2.26749 & 0.00372914 \(\pm\) 0.00332507 & 0.0309555 \(\pm\) 0.0110338 \\
Full & FDG & 1598 & 0.921129 \(\pm\) 0.0416829 & 27.223 \(\pm\) 2.14503 & 0.00224384 \(\pm\) 0.00268009 & 0.024129 \(\pm\) 0.00961375 \\
ClinicalUNet & AV45 & 1365 & 0.860412 \(\pm\) 0.0374027 & 24.637 \(\pm\) 2.02806 & 0.00388392 \(\pm\) 0.00241146 & 0.0324286 \(\pm\) 0.00954057 \\
ClinicalUNet & FDG & 1598 & 0.900311 \(\pm\) 0.0356273 & 26.6116 \(\pm\) 1.89074 & 0.00247114 \(\pm\) 0.00196941 & 0.0258502 \(\pm\) 0.00818077 \\
SwinUNETR & AV45 & 1365 & 0.874091 \(\pm\) 0.0369388 & 24.5068 \(\pm\) 2.11433 & 0.00403421 \(\pm\) 0.00249859 & 0.0320631 \(\pm\) 0.00986262 \\
SwinUNETR & FDG & 1598 & 0.918317 \(\pm\) 0.0360297 & 26.8441 \(\pm\) 1.9234 & 0.00234143 \(\pm\) 0.00181013 & 0.0241961 \(\pm\) 0.00799078 \\
TabularMLP & AV45 & 1365 & 0.874855 \(\pm\) 0.0407899 & 25.1009 \(\pm\) 2.28897 & 0.00363223 \(\pm\) 0.00321331 & 0.0308469 \(\pm\) 0.0110175 \\
TabularMLP & FDG & 1598 & 0.912262 \(\pm\) 0.0417162 & 27.2919 \(\pm\) 2.12402 & 0.00220694 \(\pm\) 0.0026887 & 0.0241971 \(\pm\) 0.00947495 \\
TracerOnly & AV45 & 1365 & 0.885748 \(\pm\) 0.0408475 & 24.9695 \(\pm\) 2.27491 & 0.00373732 \(\pm\) 0.00335201 & 0.0308555 \(\pm\) 0.0111714 \\
TracerOnly & FDG & 1598 & 0.926504 \(\pm\) 0.0407948 & 27.2745 \(\pm\) 2.14158 & 0.00221882 \(\pm\) 0.00270009 & 0.0236521 \(\pm\) 0.00949478 \\
Single & AV45 & 1365 & 0.873767 \(\pm\) 0.041007 & 24.9848 \(\pm\) 2.21914 & 0.00369036 \(\pm\) 0.00320827 & 0.031269 \(\pm\) 0.0107387 \\
Single & FDG & 1598 & 0.926933 \(\pm\) 0.0400237 & 27.0223 \(\pm\) 2.11834 & 0.00234305 \(\pm\) 0.00287526 & 0.0239414 \(\pm\) 0.00968223 \\
\end{longtable}
\noindent\parbox{\linewidth}{\footnotesize \DIFadd{Mean plus/minus sample SD. Whole-volume frozen OOF; FDG1598, AV451365. Full165, ClinicalUNet165, Swin200 MRI-only, Single165, TabularMLP165, TracerOnly165. SSIM unitless; PSNR dB; MSE squared and MAE unsquared normalized-intensity arbitrary units. Independent training realizations; reconstructed runs, not historical Figure2.}}
\endgroup
\end{landscape}

%DIF >  END supplementary/tables/whole_volume_generation


%DIF >  BEGIN supplementary/tables/paired_generation
%DIF >  Source-backed publication fragment; source records retained independently.
%DIF >  Source: G3/paired_comparisons.csv
\begin{landscape}
\begingroup
\footnotesize
\setlength{\tabcolsep}{4pt}
\begin{longtable}{lllllllllll}
\caption{{\color{blue}Paired whole-volume generation comparisons}}\label{tab:paired_generation}\\
\toprule
model & tracer & metric & n & mean\_delta & ci\_low & ci\_high & cohen\_dz & Full better pct & p\_raw & p\_holm \\
\midrule
\endfirsthead
\toprule
model & tracer & metric & n & mean\_delta & ci\_low & ci\_high & cohen\_dz & Full better pct & p\_raw & p\_holm \\
\midrule
\endhead
\bottomrule
\endfoot
ClinicalUNet & fdg & ssim & 1598 & 0.0208181 & 0.0197791 & 0.0218175 & 0.990561 & 89.8623 & 1.7603e-239 & 1.40824e-238 \\
ClinicalUNet & fdg & psnr & 1598 & 0.611381 & 0.5664 & 0.656514 & 0.648354 & 76.5957 & 6.43853e-124 & 3.86312e-123 \\
ClinicalUNet & fdg & mse & 1598 & -0.0002273 & -0.000275411 & -0.000158772 & -0.189496 & 76.5957 & 6.04086e-14 & 1.31586e-13 \\
ClinicalUNet & fdg & mae & 1598 & -0.00172122 & -0.0019156 & -0.00151429 & -0.417989 & 75.4068 & 6.9045e-58 & 3.45225e-57 \\
ClinicalUNet & av45 & ssim & 1365 & 0.0239956 & 0.0227111 & 0.025253 & 0.978795 & 87.1062 & 2.29545e-201 & 1.60681e-200 \\
ClinicalUNet & av45 & psnr & 1365 & 0.329409 & 0.249281 & 0.40985 & 0.220001 & 62.0513 & 9.68848e-16 & 3.87539e-15 \\
ClinicalUNet & av45 & mse & 1365 & -0.000154787 & -0.000257116 & -3.72771e-05 & -0.0752304 & 62.0513 & 0.00551979 & 0.00551979 \\
ClinicalUNet & av45 & mae & 1365 & -0.00147312 & -0.00185782 & -0.00109689 & -0.206517 & 63.7363 & 4.3862e-14 & 1.31586e-13 \\
SwinUNETR & fdg & ssim & 1598 & 0.00281191 & 0.00204043 & 0.00354844 & 0.175086 & 69.2741 & 3.77993e-12 & 1.88996e-11 \\
SwinUNETR & fdg & psnr & 1598 & 0.378929 & 0.334097 & 0.423834 & 0.40333 & 69.0238 & 2.67057e-54 & 1.8694e-53 \\
SwinUNETR & fdg & mse & 1598 & -9.75975e-05 & -0.000146119 & -3.08259e-05 & -0.0822636 & 69.0238 & 0.00102917 & 0.00205835 \\
SwinUNETR & fdg & mae & 1598 & -6.70702e-05 & -0.000253042 & 0.000134417 & -0.0166045 & 55.8198 & 0.506936 & 0.506936 \\
SwinUNETR & av45 & ssim & 1365 & 0.0103163 & 0.00965209 & 0.0109524 & 0.822615 & 83.4432 & 2.31981e-155 & 1.85585e-154 \\
SwinUNETR & av45 & psnr & 1365 & 0.459673 & 0.393978 & 0.526625 & 0.375353 & 66.8132 & 5.23792e-41 & 3.14275e-40 \\
SwinUNETR & av45 & mse & 1365 & -0.000305079 & -0.000395415 & -0.000193256 & -0.160805 & 66.8132 & 3.58919e-09 & 1.07676e-08 \\
SwinUNETR & av45 & mae & 1365 & -0.00110762 & -0.0014185 & -0.000784845 & -0.18539 & 61.4652 & 1.11886e-11 & 4.47543e-11 \\
TabularMLP & fdg & ssim & 1598 & 0.00886656 & 0.00812186 & 0.00959853 & 0.578166 & 75.8448 & 3.44323e-102 & 2.75458e-101 \\
TabularMLP & fdg & psnr & 1598 & -0.0689397 & -0.108799 & -0.029079 & -0.083771 & 47.3091 & 0.000830676 & 0.00425048 \\
TabularMLP & fdg & mse & 1598 & 3.68949e-05 & 1.37346e-05 & 5.99572e-05 & 0.0775476 & 47.3091 & 0.00196924 & 0.00787697 \\
TabularMLP & fdg & mae & 1598 & -6.81443e-05 & -0.000213963 & 8.04426e-05 & -0.0227372 & 53.3166 & 0.363529 & 0.727058 \\
TabularMLP & av45 & ssim & 1365 & 0.00955263 & 0.00863504 & 0.0104854 & 0.54034 & 74.2125 & 5.41378e-78 & 3.78964e-77 \\
TabularMLP & av45 & psnr & 1365 & -0.134409 & -0.210528 & -0.0561215 & -0.0918656 & 46.8132 & 0.000708413 & 0.00425048 \\
TabularMLP & av45 & mse & 1365 & 9.69028e-05 & 1.94112e-05 & 0.000173411 & 0.0677817 & 46.8132 & 0.0123868 & 0.0371605 \\
TabularMLP & av45 & mae & 1365 & 0.000108651 & -0.000232867 & 0.000435387 & 0.0173217 & 51.1355 & 0.522301 & 0.727058 \\
TracerOnly & fdg & ssim & 1598 & -0.00537535 & -0.00590994 & -0.00484902 & -0.47433 & 26.9086 & 1.79244e-72 & 1.43395e-71 \\
TracerOnly & fdg & psnr & 1598 & -0.0514788 & -0.084006 & -0.0198543 & -0.0783114 & 48.8736 & 0.0017767 & 0.00888351 \\
TracerOnly & fdg & mse & 1598 & 2.5016e-05 & 5.74477e-06 & 4.40603e-05 & 0.0629846 & 48.8736 & 0.0119058 & 0.0476234 \\
TracerOnly & fdg & mae & 1598 & 0.000476913 & 0.000360989 & 0.000595372 & 0.197824 & 42.3029 & 4.84122e-15 & 3.38885e-14 \\
TracerOnly & av45 & ssim & 1365 & -0.0013401 & -0.00186153 & -0.000811475 & -0.132833 & 44.1026 & 1.03289e-06 & 6.19734e-06 \\
TracerOnly & av45 & psnr & 1365 & -0.00307791 & -0.0547682 & 0.0476535 & -0.0032182 & 47.9853 & 0.905373 & 1 \\
TracerOnly & av45 & mse & 1365 & -8.18237e-06 & -5.70224e-05 & 4.02036e-05 & -0.00891382 & 47.9853 & 0.741957 & 1 \\
TracerOnly & av45 & mae & 1365 & 0.000100047 & -0.000115014 & 0.000313809 & 0.024782 & 46.8132 & 0.360042 & 1 \\
Single & fdg & ssim & 1598 & -0.00580378 & -0.00642246 & -0.00517901 & -0.434763 & 34.2303 & 4.22015e-62 & 2.95411e-61 \\
Single & fdg & psnr & 1598 & 0.200636 & 0.159293 & 0.242418 & 0.237103 & 60.5131 & 8.98722e-21 & 5.39233e-20 \\
Single & fdg & mse & 1598 & -9.92098e-05 & -0.000126504 & -7.21876e-05 & -0.177651 & 60.5131 & 1.84961e-12 & 9.24804e-12 \\
Single & fdg & mae & 1598 & 0.000187601 & 3.76336e-05 & 0.00034206 & 0.0596374 & 47.3717 & 0.0172416 & 0.0527503 \\
Single & av45 & ssim & 1365 & 0.0106401 & 0.00985376 & 0.0114433 & 0.707251 & 78.1685 & 2.31929e-122 & 1.85543e-121 \\
Single & av45 & psnr & 1365 & -0.0183712 & -0.0770958 & 0.0388809 & -0.0169039 & 49.8168 & 0.532384 & 0.532384 \\
Single & av45 & mse & 1365 & 3.87787e-05 & -1.71053e-05 & 9.51783e-05 & 0.0366412 & 49.8168 & 0.176043 & 0.352086 \\
Single & av45 & mae & 1365 & -0.000313478 & -0.000560844 & -6.84978e-05 & -0.0671769 & 54.8718 & 0.0131876 & 0.0527503 \\
\end{longtable}
\noindent\parbox{\linewidth}{\footnotesize \DIFadd{Delta=Full minus comparator; positive favors Full for SSIM/PSNR and negative for MSE/MAE. Cohen dz signed.10000 paired resamples within frozen folds, seed42; percentile95\% CIs unadjusted. Two-sided paired t; separate Holm8-test family per comparator (2tracers x4metrics). ClinicalUNet statistics reused; TracerOnly t/dz reused and missing bootstrap supplemented. CIs do not estimate training/seed variability. MRI-only Swin inputs and endpoint differ.}}
\endgroup
\end{landscape}

%DIF >  END supplementary/tables/paired_generation


%DIF >  BEGIN supplementary/tables/masked_summary
\begin{table*}[t]
\centering
\caption{{\color{blue}Mask-restricted quality of frozen two-fold OOF PET predictions.}}
\label{tab:r3_11_masked_summary}
\resizebox{\textwidth}{!}{%
\begin{tabular}{lllrrrr}
\hline
Model & Tracer & $N$ & SSIM $\uparrow$ & PSNR (dB) $\uparrow$ & MSE $\downarrow$ & MAE $\downarrow$ \\
\hline
DIReCT++ Full (e165) & FDG & 1597 & 0.91844 $\pm$ 0.04026 & 23.954 $\pm$ 2.271 & 0.004763 $\pm$ 0.004028 & 0.05257 $\pm$ 0.01921 \\
DIReCT++ Full (e165) & AV45 & 1364 & 0.86999 $\pm$ 0.06048 & 21.803 $\pm$ 2.602 & 0.007959 $\pm$ 0.005469 & 0.06845 $\pm$ 0.02394 \\
Clinical Regression 3D U-Net (e165) & FDG & 1597 & 0.90208 $\pm$ 0.04256 & 23.373 $\pm$ 2.019 & 0.005262 $\pm$ 0.003889 & 0.05509 $\pm$ 0.01795 \\
Clinical Regression 3D U-Net (e165) & AV45 & 1364 & 0.86098 $\pm$ 0.05868 & 21.529 $\pm$ 2.363 & 0.008237 $\pm$ 0.005320 & 0.06976 $\pm$ 0.02276 \\
Swin-UNETR (e200; MRI only) & FDG & 1597 & 0.89237 $\pm$ 0.03938 & 23.565 $\pm$ 2.013 & 0.004997 $\pm$ 0.003300 & 0.05404 $\pm$ 0.01672 \\
Swin-UNETR (e200; MRI only) & AV45 & 1364 & 0.83615 $\pm$ 0.05878 & 21.273 $\pm$ 2.429 & 0.008791 $\pm$ 0.005652 & 0.07228 $\pm$ 0.02363 \\
\hline
\end{tabular}}
\par\smallskip
\begin{minipage}{\textwidth}\footnotesize
\DIFaddFL{Values are mean $\pm$ sample SD across subjects. SSIM is unitless; PSNR is in dB; MSE and MAE use squared and unsquared }[\DIFaddFL{0,1}] \DIFaddFL{normalized-uptake units, respectively. All models use byte-identical frozen target PET arrays and identical subject-specific masks. The mask is the frozen SynthSeg nonzero-label union, including cortex, white matter, subcortical structures, cerebellum, brainstem, ventricles and CSF; it is not a pure parenchymal mask. SSIM uses the existing E1 3D Gaussian implementation (kernel 7, sigma 1.5, data range 1) and only centers whose entire 7$\times$7$\times$7 window lies inside the mask. The frozen mask QC exclusion removes one subject from each tracer; no new exclusions or spatial/intensity transformations were applied. Full and Clinical U-Net use epoch 165; MRI-only Swin-UNETR uses the prespecified final epoch 200. All predictions are frozen two-fold OOF outputs.
}\end{minipage}
\end{table*}

%DIF >  END supplementary/tables/masked_summary


%DIF >  BEGIN supplementary/tables/masked_paired
\begin{table*}[t]
\centering
\caption{{\color{blue}Paired mask-restricted PET image-quality comparisons.}}
\label{tab:r3_11_masked_paired}
\resizebox{\textwidth}{!}{%
\begin{tabular}{lllrrrrrrrr}
\hline
Comparison & Tracer & Metric & $N$ & Mean $\Delta$ & 95\% bootstrap CI & $t$ & Cohen $d_z$ & Full better (\%) & Raw $p$ & Holm $p$ \\
\hline
A: Full $-$ Clinical U-Net & FDG & SSIM & 1597 & 0.01636 & [0.01577, 0.01699] & 51.477 & 1.288 & 95.7 & $<10^{-300}$ & $<10^{-300}$ \\
A: Full $-$ Clinical U-Net & FDG & PSNR & 1597 & 0.581 & [0.529, 0.634] & 21.474 & 0.537 & 72.3 & 4.55e-90 & 3.18e-89 \\
A: Full $-$ Clinical U-Net & FDG & MSE & 1597 & -0.000499 & [-0.000572, -0.000425] & -13.116 & -0.328 & 72.3 & 2.13e-37 & 1.28e-36 \\
A: Full $-$ Clinical U-Net & FDG & MAE & 1597 & -0.00253 & [-0.00295, -0.00210] & -11.670 & -0.292 & 66.9 & 2.96e-30 & 1.18e-29 \\
A: Full $-$ Clinical U-Net & AV45 & SSIM & 1364 & 0.00901 & [0.00754, 0.01051] & 11.881 & 0.322 & 69.0 & 4.79e-31 & 2.39e-30 \\
A: Full $-$ Clinical U-Net & AV45 & PSNR & 1364 & 0.273 & [0.178, 0.370] & 5.585 & 0.151 & 57.8 & 2.82e-08 & 8.45e-08 \\
A: Full $-$ Clinical U-Net & AV45 & MSE & 1364 & -0.000278 & [-0.000488, -0.000066] & -2.609 & -0.071 & 57.8 & 0.00917 & 0.00962 \\
A: Full $-$ Clinical U-Net & AV45 & MAE & 1364 & -0.00131 & [-0.00223, -0.00040] & -2.824 & -0.076 & 56.4 & 0.00481 & 0.00962 \\
B: Full $-$ Swin-UNETR & FDG & SSIM & 1597 & 0.02607 & [0.02553, 0.02663] & 92.200 & 2.307 & 99.0 & $<10^{-300}$ & $<10^{-300}$ \\
B: Full $-$ Swin-UNETR & FDG & PSNR & 1597 & 0.389 & [0.337, 0.443] & 14.239 & 0.356 & 66.8 & 2.11e-43 & 1.26e-42 \\
B: Full $-$ Swin-UNETR & FDG & MSE & 1597 & -0.000233 & [-0.000311, -0.000151] & -5.733 & -0.143 & 66.8 & 1.18e-08 & 1.18e-08 \\
B: Full $-$ Swin-UNETR & FDG & MAE & 1597 & -0.00148 & [-0.00191, -0.00105] & -6.681 & -0.167 & 63.3 & 3.26e-11 & 6.52e-11 \\
B: Full $-$ Swin-UNETR & AV45 & SSIM & 1364 & 0.03383 & [0.03274, 0.03493] & 59.128 & 1.601 & 95.2 & $<10^{-300}$ & $<10^{-300}$ \\
B: Full $-$ Swin-UNETR & AV45 & PSNR & 1364 & 0.530 & [0.451, 0.609] & 13.216 & 0.358 & 65.1 & 1.3e-37 & 6.52e-37 \\
B: Full $-$ Swin-UNETR & AV45 & MSE & 1364 & -0.000832 & [-0.001003, -0.000658] & -9.563 & -0.259 & 65.1 & 5.15e-21 & 1.55e-20 \\
B: Full $-$ Swin-UNETR & AV45 & MAE & 1364 & -0.00383 & [-0.00459, -0.00309] & -10.003 & -0.271 & 64.1 & 8.81e-23 & 3.52e-22 \\
\hline
\end{tabular}}
\par\smallskip
\begin{minipage}{\textwidth}\footnotesize
\DIFaddFL{$\Delta$ = DIReCT++ Full minus comparator on exactly matched subject, tracer and held-out fold. Positive differences favor Full for SSIM/PSNR; negative differences favor Full for MSE/MAE. Cohen's dz uses this signed difference divided by its sample SD. Full-better proportions use strict metric-specific improvement. Two-sided paired t-tests; 95\% percentile bootstrap CIs use 10,000 paired resamples, seed 42, stratified by held-out fold with original fold sizes retained. Holm correction is applied separately to comparison A and comparison B, with eight tests (2 tracers $\times$ 4 metrics) in each family; CIs are unadjusted. Uncertainty concerns the frozen OOF subjects, not retraining/seed variability. Statistical significance is not clinical equivalence. The two comparisons differ in information and endpoint: A uses matched MRI/clinical/tracer information at e165; B compares Full e165 with MRI-only Swin-UNETR e200. A displayed p $<10^{-300}$ denotes numerical underflow of the computed p value, not an exactly zero probability.
}\end{minipage}
\end{table*}

%DIF >  END supplementary/tables/masked_paired


%DIF >  BEGIN supplementary/tables/generation_cases
\begin{figure}[p]
\centering
\includegraphics[width=0.97\textwidth]{revision_cases_fdg}
\caption{{\color{blue}\textbf{FDG synthesis in median, 10th-percentile and worst QC-passed cases.} Rows follow the frozen whole-volume SSIM selection; columns show MRI, real PET, synthetic PET and absolute-error maps. PET intensities and errors use the displayed [0,1] normalized-intensity scales. Selection is within the existing QC-passed set, without new exclusions or visual replacement; rankings use whole-volume SSIM rather than masked SSIM.}}
\label{fig:cases_fdg}
\end{figure}
\begin{figure}[p]
\centering
\includegraphics[width=0.97\textwidth]{revision_cases_av45}
\caption{{\color{blue}\textbf{AV45 synthesis in median, 10th-percentile and worst QC-passed cases.} Rows follow the frozen whole-volume SSIM selection; columns show MRI, real PET, synthetic PET and absolute-error maps. PET intensities and errors use the displayed [0,1] normalized-intensity scales. Selection is within the existing QC-passed set, without new exclusions or visual replacement; rankings use whole-volume SSIM rather than masked SSIM.}}
\label{fig:cases_av45}
\end{figure}

%DIF >  END supplementary/tables/generation_cases


%DIF >  BEGIN supplementary/tables/classification_summary
\begin{landscape}
\begingroup\small\setlength{\tabcolsep}{5pt}
\begin{longtable}{lccccc}
\caption{{\color{blue}Three-fold classification results: mean $\pm$ sample s.d.; [conditional 95\% CI].}}\label{tab:classification_summary}\\
\toprule
Input & ACC & AUC & SEN & SPE & Balanced ACC \\
\midrule\endfirsthead
\toprule Input & ACC & AUC & SEN & SPE & Balanced ACC \\\midrule\endhead
\multicolumn{6}{l}{\textbf{AD/CN}} \\
MRI & \shortstack{79.56 $\pm$ 6.20\\{}[75.95, 82.97]} & \shortstack{87.02 $\pm$ 5.44\\{}[83.51, 90.15]} & \shortstack{75.15 $\pm$ 10.01\\{}[68.48, 81.21]} & \shortstack{81.96 $\pm$ 6.90\\{}[77.70, 85.91]} & \shortstack{78.56 $\pm$ 6.59\\{}[74.63, 82.24]} \\
MRI+Real FDG & \shortstack{89.36 $\pm$ 0.72\\{}[86.39, 91.92]} & \shortstack{94.28 $\pm$ 0.55\\{}[91.77, 96.51]} & \shortstack{81.21 $\pm$ 6.39\\{}[75.15, 87.27]} & \shortstack{93.78 $\pm$ 4.41\\{}[90.85, 96.40]} & \shortstack{87.50 $\pm$ 1.06\\{}[84.07, 90.69]} \\
MRI+Syn FDG & \shortstack{88.72 $\pm$ 3.31\\{}[85.93, 91.48]} & \shortstack{94.40 $\pm$ 0.89\\{}[91.96, 96.58]} & \shortstack{80.00 $\pm$ 12.73\\{}[73.94, 86.06]} & \shortstack{93.45 $\pm$ 3.95\\{}[90.51, 96.07]} & \shortstack{86.73 $\pm$ 5.21\\{}[83.35, 89.97]} \\
MRI+Real AV45 & \shortstack{85.53 $\pm$ 2.67\\{}[82.34, 88.51]} & \shortstack{92.54 $\pm$ 2.53\\{}[89.91, 94.92]} & \shortstack{74.55 $\pm$ 1.82\\{}[67.88, 81.21]} & \shortstack{91.47 $\pm$ 4.86\\{}[88.19, 94.42]} & \shortstack{83.01 $\pm$ 1.77\\{}[79.27, 86.58]} \\
MRI+Syn AV45 & \shortstack{84.04 $\pm$ 2.80\\{}[80.84, 87.43]} & \shortstack{91.81 $\pm$ 3.54\\{}[89.02, 94.23]} & \shortstack{81.21 $\pm$ 2.10\\{}[75.15, 86.68]} & \shortstack{85.56 $\pm$ 4.90\\{}[81.62, 89.49]} & \shortstack{83.39 $\pm$ 2.00\\{}[79.87, 86.91]} \\
MRI+Real both & \shortstack{90.64 $\pm$ 2.59\\{}[87.87, 93.18]} & \shortstack{94.38 $\pm$ 0.90\\{}[91.84, 96.61]} & \shortstack{84.24 $\pm$ 2.78\\{}[78.18, 89.70]} & \shortstack{94.09 $\pm$ 2.62\\{}[91.43, 96.40]} & \shortstack{89.17 $\pm$ 2.60\\{}[85.92, 92.14]} \\
MRI+Syn both & \shortstack{88.08 $\pm$ 2.06\\{}[85.09, 90.85]} & \shortstack{93.82 $\pm$ 3.21\\{}[91.18, 96.07]} & \shortstack{86.06 $\pm$ 5.84\\{}[80.61, 90.91]} & \shortstack{89.18 $\pm$ 0.93\\{}[85.57, 92.46]} & \shortstack{87.62 $\pm$ 2.90\\{}[84.29, 90.67]} \\
Clinical-only & \shortstack{100.00 $\pm$ 0.00\\{}[100.00, 100.00]} & \shortstack{100.00 $\pm$ 0.00\\{}[100.00, 100.00]} & \shortstack{100.00 $\pm$ 0.00\\{}[100.00, 100.00]} & \shortstack{100.00 $\pm$ 0.00\\{}[100.00, 100.00]} & \shortstack{100.00 $\pm$ 0.00\\{}[100.00, 100.00]} \\
MRI+Clinical & \shortstack{94.89 $\pm$ 3.54\\{}[92.77, 96.81]} & \shortstack{99.17 $\pm$ 0.96\\{}[98.50, 99.66]} & \shortstack{92.12 $\pm$ 4.20\\{}[87.88, 95.76]} & \shortstack{96.39 $\pm$ 3.44\\{}[94.10, 98.36]} & \shortstack{94.26 $\pm$ 3.65\\{}[91.90, 96.43]} \\
\midrule
\multicolumn{6}{l}{\textbf{EMCI+LMCI/CN}} \\
MRI & \shortstack{67.48 $\pm$ 0.23\\{}[64.37, 70.58]} & \shortstack{71.59 $\pm$ 1.80\\{}[67.83, 75.25]} & \shortstack{79.57 $\pm$ 11.66\\{}[75.89, 83.23]} & \shortstack{50.12 $\pm$ 17.10\\{}[44.56, 55.38]} & \shortstack{64.85 $\pm$ 2.73\\{}[61.57, 68.12]} \\
MRI+Real FDG & \shortstack{61.27 $\pm$ 2.70\\{}[57.76, 64.78]} & \shortstack{64.82 $\pm$ 2.79\\{}[60.74, 68.77]} & \shortstack{63.30 $\pm$ 3.89\\{}[58.69, 67.69]} & \shortstack{58.35 $\pm$ 1.73\\{}[52.79, 63.90]} & \shortstack{60.83 $\pm$ 2.48\\{}[57.14, 64.49]} \\
MRI+Syn FDG & \shortstack{62.48 $\pm$ 3.25\\{}[59.11, 65.86]} & \shortstack{65.77 $\pm$ 5.56\\{}[61.81, 69.81]} & \shortstack{67.00 $\pm$ 12.26\\{}[62.64, 71.13]} & \shortstack{56.10 $\pm$ 14.64\\{}[50.82, 61.35]} & \shortstack{61.55 $\pm$ 3.28\\{}[58.05, 65.11]} \\
MRI+Real AV45 & \shortstack{66.80 $\pm$ 2.10\\{}[63.56, 69.91]} & \shortstack{72.00 $\pm$ 0.79\\{}[68.23, 75.75]} & \shortstack{75.45 $\pm$ 8.66\\{}[71.54, 79.35]} & \shortstack{54.43 $\pm$ 11.27\\{}[48.87, 59.69]} & \shortstack{64.94 $\pm$ 2.42\\{}[61.54, 68.28]} \\
MRI+Syn AV45 & \shortstack{65.99 $\pm$ 1.21\\{}[62.75, 69.23]} & \shortstack{71.52 $\pm$ 1.85\\{}[67.73, 75.10]} & \shortstack{73.17 $\pm$ 6.51\\{}[69.04, 77.30]} & \shortstack{55.76 $\pm$ 11.84\\{}[50.49, 61.32]} & \shortstack{64.47 $\pm$ 2.73\\{}[61.05, 67.90]} \\
MRI+Real both & \shortstack{63.70 $\pm$ 3.54\\{}[60.19, 67.21]} & \shortstack{69.15 $\pm$ 2.42\\{}[65.13, 72.90]} & \shortstack{65.60 $\pm$ 9.91\\{}[61.24, 69.96]} & \shortstack{60.97 $\pm$ 8.42\\{}[55.41, 66.21]} & \shortstack{63.29 $\pm$ 2.81\\{}[59.76, 66.79]} \\
MRI+Syn both & \shortstack{65.45 $\pm$ 1.64\\{}[62.48, 68.56]} & \shortstack{69.46 $\pm$ 1.34\\{}[65.60, 73.18]} & \shortstack{79.33 $\pm$ 12.80\\{}[75.65, 83.00]} & \shortstack{45.53 $\pm$ 14.60\\{}[39.95, 51.08]} & \shortstack{62.43 $\pm$ 0.94\\{}[59.14, 65.71]} \\
Clinical-only & \shortstack{94.87 $\pm$ 2.60\\{}[93.25, 96.36]} & \shortstack{98.93 $\pm$ 0.70\\{}[98.14, 99.56]} & \shortstack{93.59 $\pm$ 4.17\\{}[91.29, 95.87]} & \shortstack{96.72 $\pm$ 1.14\\{}[94.43, 98.68]} & \shortstack{95.15 $\pm$ 2.27\\{}[93.62, 96.59]} \\
MRI+Clinical & \shortstack{85.16 $\pm$ 8.58\\{}[82.73, 87.45]} & \shortstack{93.07 $\pm$ 6.10\\{}[91.14, 94.77]} & \shortstack{88.51 $\pm$ 12.11\\{}[85.52, 91.26]} & \shortstack{80.35 $\pm$ 18.35\\{}[76.08, 84.28]} & \shortstack{84.43 $\pm$ 9.22\\{}[81.82, 86.84]} \\
\midrule
\multicolumn{6}{l}{\textbf{EMCI/LMCI}} \\
MRI & \shortstack{64.90 $\pm$ 3.09\\{}[61.92, 67.89]} & \shortstack{62.81 $\pm$ 2.50\\{}[57.39, 68.38]} & \shortstack{17.40 $\pm$ 12.64\\{}[12.03, 23.43]} & \shortstack{89.53 $\pm$ 6.57\\{}[85.70, 93.01]} & \shortstack{53.46 $\pm$ 4.63\\{}[50.12, 56.96]} \\
MRI+Real FDG & \shortstack{64.21 $\pm$ 4.93\\{}[60.30, 68.11]} & \shortstack{61.66 $\pm$ 3.35\\{}[56.04, 67.09]} & \shortstack{36.86 $\pm$ 16.63\\{}[29.48, 44.26]} & \shortstack{78.34 $\pm$ 15.69\\{}[73.80, 82.87]} & \shortstack{57.60 $\pm$ 1.88\\{}[53.23, 61.87]} \\
MRI+Syn FDG & \shortstack{65.36 $\pm$ 3.46\\{}[61.47, 69.26]} & \shortstack{63.02 $\pm$ 1.14\\{}[57.49, 68.72]} & \shortstack{32.84 $\pm$ 5.47\\{}[25.46, 40.88]} & \shortstack{82.20 $\pm$ 6.97\\{}[77.67, 86.39]} & \shortstack{57.52 $\pm$ 2.15\\{}[53.15, 62.08]} \\
MRI+Real AV45 & \shortstack{65.82 $\pm$ 3.65\\{}[62.14, 69.49]} & \shortstack{60.62 $\pm$ 2.37\\{}[55.03, 66.28]} & \shortstack{29.50 $\pm$ 9.10\\{}[22.16, 36.87]} & \shortstack{84.66 $\pm$ 1.30\\{}[80.15, 88.83]} & \shortstack{57.08 $\pm$ 4.92\\{}[52.84, 61.28]} \\
MRI+Syn AV45 & \shortstack{65.37 $\pm$ 5.56\\{}[61.47, 69.27]} & \shortstack{62.12 $\pm$ 6.88\\{}[56.59, 67.70]} & \shortstack{36.19 $\pm$ 6.89\\{}[28.76, 43.61]} & \shortstack{80.46 $\pm$ 11.73\\{}[75.92, 84.97]} & \shortstack{58.32 $\pm$ 2.46\\{}[53.92, 62.83]} \\
MRI+Real both & \shortstack{64.90 $\pm$ 4.88\\{}[60.78, 68.81]} & \shortstack{63.04 $\pm$ 1.25\\{}[57.58, 68.56]} & \shortstack{41.05 $\pm$ 19.06\\{}[33.66, 48.48]} & \shortstack{77.31 $\pm$ 13.00\\{}[72.42, 81.85]} & \shortstack{59.18 $\pm$ 5.52\\{}[54.89, 63.59]} \\
MRI+Syn both & \shortstack{65.82 $\pm$ 2.06\\{}[61.69, 70.17]} & \shortstack{66.33 $\pm$ 0.54\\{}[61.14, 71.76]} & \shortstack{43.03 $\pm$ 10.87\\{}[34.99, 51.10]} & \shortstack{77.71 $\pm$ 8.82\\{}[72.83, 82.25]} & \shortstack{60.37 $\pm$ 1.12\\{}[55.87, 65.19]} \\
Clinical-only & \shortstack{99.08 $\pm$ 0.40\\{}[98.16, 99.77]} & \shortstack{99.94 $\pm$ 0.08\\{}[99.79, 100.00]} & \shortstack{98.67 $\pm$ 2.31\\{}[96.67, 100.00]} & \shortstack{99.31 $\pm$ 0.60\\{}[98.26, 100.00]} & \shortstack{98.99 $\pm$ 0.85\\{}[97.81, 99.83]} \\
MRI+Clinical & \shortstack{85.32 $\pm$ 13.80\\{}[82.34, 88.53]} & \shortstack{89.60 $\pm$ 9.08\\{}[86.11, 93.02]} & \shortstack{85.90 $\pm$ 12.00\\{}[80.50, 91.28]} & \shortstack{85.07 $\pm$ 18.35\\{}[81.25, 88.89]} & \shortstack{85.49 $\pm$ 12.24\\{}[82.16, 88.89]} \\
\midrule
\bottomrule
\end{longtable}
\noindent\footnotesize \DIFadd{All metrics are percentages; s.d. is in percentage points. CIs use 5,000 paired fold-by-class bootstrap resamples of frozen predictions and three-fold equal weighting. Intervals are marginal, unadjusted and conditional on the fixed models and splits. Clinical-only AD/CN's zero empirical errors yield a degenerate }[\DIFadd{100,100}] \DIFadd{bootstrap interval, not a population guarantee. Clinical inputs include diagnosis-conditioned fills. Sample sizes are AD/CN=470, EMCI+LMCI/CN=741 and EMCI/LMCI=436; MCI-U is excluded.
}\endgroup
\end{landscape}

%DIF >  END supplementary/tables/classification_summary


%DIF >  BEGIN supplementary/tables/classification_imaging_comparisons
%DIF >  Draft only. Requires longtable, booktabs; wide tables also pdflscape.
%DIF >  Source: G5/existing_imaging_paired_comparisons.csv
\begin{landscape}
\begingroup
\footnotesize
\setlength{\tabcolsep}{4pt}
\begin{longtable}{lllllll}
\caption{{\color{blue}Paired imaging classification differences}}\label{tab:classification_comparisons}\\
\toprule
task & target & reference & metric & delta\_pp & ci\_low\_pp & ci\_high\_pp \\
\midrule
\endfirsthead
\toprule
task & target & reference & metric & delta\_pp & ci\_low\_pp & ci\_high\_pp \\
\midrule
\endhead
\bottomrule
\endfoot
ad\_cn & mri\_syn\_both & mri & ACC & 8.51843 & 4.47901 & 12.5606 \\
ad\_cn & mri\_syn\_both & mri & AUC & 6.80203 & 3.38872 & 10.2184 \\
ad\_cn & mri\_syn\_both & mri & SEN & 10.9091 & 3.63636 & 18.7879 \\
ad\_cn & mri\_syn\_both & mri & SPE & 7.21866 & 2.31023 & 12.1369 \\
ad\_cn & mri\_syn\_both & mri & balanced\_ACC & 9.06388 & 4.68182 & 13.457 \\
ad\_cn & mri\_syn\_both & mri\_real\_both & ACC & -2.55185 & -5.31194 & -0.00272198 \\
ad\_cn & mri\_syn\_both & mri\_real\_both & AUC & -0.55129 & -2.11379 & 0.976434 \\
ad\_cn & mri\_syn\_both & mri\_real\_both & SEN & 1.81818 & -3.0303 & 6.66667 \\
ad\_cn & mri\_syn\_both & mri\_real\_both & SPE & -4.90843 & -8.18935 & -1.95431 \\
ad\_cn & mri\_syn\_both & mri\_real\_both & balanced\_ACC & -1.54513 & -4.34698 & 1.11062 \\
mci\_cn & mri\_syn\_both & mri & ACC & -2.02429 & -5.66802 & 1.61943 \\
mci\_cn & mri\_syn\_both & mri & AUC & -2.12981 & -6.29816 & 2.09665 \\
mci\_cn & mri\_syn\_both & mri & SEN & -0.239332 & -4.60404 & 4.1301 \\
mci\_cn & mri\_syn\_both & mri & SPE & -4.59458 & -11.1499 & 1.93174 \\
mci\_cn & mri\_syn\_both & mri & balanced\_ACC & -2.41695 & -6.28258 & 1.42342 \\
mci\_cn & mri\_syn\_both & mri\_real\_both & ACC & 1.75439 & -1.88934 & 5.39811 \\
mci\_cn & mri\_syn\_both & mri\_real\_both & AUC & 0.308348 & -3.14901 & 3.81391 \\
mci\_cn & mri\_syn\_both & mri\_real\_both & SEN & 13.727 & 9.14344 & 18.3059 \\
mci\_cn & mri\_syn\_both & mri\_real\_both & SPE & -15.4436 & -21.3454 & -9.55154 \\
mci\_cn & mri\_syn\_both & mri\_real\_both & balanced\_ACC & -0.858316 & -4.5701 & 3.0065 \\
emci\_lmci & mri\_syn\_both & mri & ACC & 0.91954 & -3.44355 & 5.50626 \\
emci\_lmci & mri\_syn\_both & mri & AUC & 3.52145 & -2.42365 & 9.50467 \\
emci\_lmci & mri\_syn\_both & mri & SEN & 25.6327 & 16.9248 & 33.7415 \\
emci\_lmci & mri\_syn\_both & mri & SPE & -11.8165 & -17.0651 & -6.5826 \\
emci\_lmci & mri\_syn\_both & mri & balanced\_ACC & 6.90807 & 2.01983 & 11.9828 \\
emci\_lmci & mri\_syn\_both & mri\_real\_both & ACC & 0.91954 & -3.66714 & 5.50937 \\
emci\_lmci & mri\_syn\_both & mri\_real\_both & AUC & 3.29387 & -2.11555 & 8.53926 \\
emci\_lmci & mri\_syn\_both & mri\_real\_both & SEN & 1.98639 & -6.13639 & 10.0548 \\
emci\_lmci & mri\_syn\_both & mri\_real\_both & SPE & 0.398392 & -4.83187 & 5.965 \\
emci\_lmci & mri\_syn\_both & mri\_real\_both & balanced\_ACC & 1.19239 & -3.73404 & 6.05574 \\
\end{longtable}
\noindent\parbox{\linewidth}{\footnotesize \DIFadd{Original6 prespecified comparisons reused; Syn dual PET minus MRI orReal dual PET. Marginal95\%5000 fold/class bootstrap CI, uncorrected for multiplicity. No overall DeLong claim and no equivalence inference.}}
\endgroup
\end{landscape}

%DIF >  END supplementary/tables/classification_imaging_comparisons


%DIF >  BEGIN supplementary/tables/classification_clinical_comparisons
%DIF >  Draft only. Requires longtable, booktabs; wide tables also pdflscape.
%DIF >  Source: G5/clinical_paired_comparisons.csv
\begin{landscape}
\begingroup
\footnotesize
\setlength{\tabcolsep}{4pt}
\begin{longtable}{lllllll}
\caption{{\color{blue}Paired comparisons with clinical baselines}}\label{tab:classification_clinical_comparisons}\\
\toprule
task & target & reference & metric & delta\_pp & ci\_low\_pp & ci\_high\_pp \\
\midrule
\endfirsthead
\toprule
task & target & reference & metric & delta\_pp & ci\_low\_pp & ci\_high\_pp \\
\midrule
\endhead
\bottomrule
\endfoot
ad\_cn & mri\_syn\_both & clinical & ACC & -11.9168 & -14.907 & -9.14853 \\
ad\_cn & mri\_syn\_both & clinical & AUC & -6.17603 & -8.82061 & -3.9256 \\
ad\_cn & mri\_syn\_both & clinical & SEN & -13.9394 & -19.3939 & -9.09091 \\
ad\_cn & mri\_syn\_both & clinical & SPE & -10.8167 & -14.4309 & -7.53899 \\
ad\_cn & mri\_syn\_both & clinical & balanced\_ACC & -12.378 & -15.7128 & -9.32558 \\
ad\_cn & mri\_syn\_both & mri\_clinical & ACC & -6.81175 & -10.2129 & -3.62159 \\
ad\_cn & mri\_syn\_both & mri\_clinical & AUC & -5.34224 & -7.87614 & -3.17673 \\
ad\_cn & mri\_syn\_both & mri\_clinical & SEN & -6.06061 & -11.5152 & -0.606061 \\
ad\_cn & mri\_syn\_both & mri\_clinical & SPE & -7.20896 & -11.1467 & -3.59469 \\
ad\_cn & mri\_syn\_both & mri\_clinical & balanced\_ACC & -6.63478 & -10.198 & -3.20788 \\
ad\_cn & mri\_clinical & mri & ACC & 15.3302 & 12.1373 & 18.9395 \\
ad\_cn & mri\_clinical & mri & AUC & 12.1443 & 9.19898 & 15.4179 \\
ad\_cn & mri\_clinical & mri & SEN & 16.9697 & 10.9091 & 23.0303 \\
ad\_cn & mri\_clinical & mri & SPE & 14.4276 & 10.4995 & 18.6825 \\
ad\_cn & mri\_clinical & mri & balanced\_ACC & 15.6987 & 12.1825 & 19.4301 \\
ad\_cn & clinical & mri & ACC & 20.4352 & 17.0314 & 24.046 \\
ad\_cn & clinical & mri & AUC & 12.9781 & 9.84708 & 16.491 \\
ad\_cn & clinical & mri & SEN & 24.8485 & 18.7879 & 31.5152 \\
ad\_cn & clinical & mri & SPE & 18.0353 & 14.0911 & 22.3031 \\
ad\_cn & clinical & mri & balanced\_ACC & 21.4419 & 17.7596 & 25.3654 \\
mci\_cn & mri\_syn\_both & clinical & ACC & -29.4197 & -32.7935 & -26.0459 \\
mci\_cn & mri\_syn\_both & clinical & AUC & -29.4749 & -33.432 & -25.6737 \\
mci\_cn & mri\_syn\_both & clinical & SEN & -14.2592 & -18.3908 & -10.1259 \\
mci\_cn & mri\_syn\_both & clinical & SPE & -51.1907 & -57.1087 & -45.2433 \\
mci\_cn & mri\_syn\_both & clinical & balanced\_ACC & -32.7249 & -36.388 & -29.1387 \\
mci\_cn & mri\_syn\_both & mri\_clinical & ACC & -19.7031 & -23.2119 & -16.3293 \\
mci\_cn & mri\_syn\_both & mri\_clinical & AUC & -23.6109 & -27.4845 & -19.805 \\
mci\_cn & mri\_syn\_both & mri\_clinical & SEN & -9.17651 & -13.0977 & -5.26531 \\
mci\_cn & mri\_syn\_both & mri\_clinical & SPE & -34.8185 & -41.3675 & -28.6029 \\
mci\_cn & mri\_syn\_both & mri\_clinical & balanced\_ACC & -21.9975 & -25.6565 & -18.3658 \\
mci\_cn & mri\_clinical & mri & ACC & 17.6788 & 14.4399 & 20.9177 \\
mci\_cn & mri\_clinical & mri & AUC & 21.4811 & 18.0947 & 24.9968 \\
mci\_cn & mri\_clinical & mri & SEN & 8.93718 & 5.25744 & 12.6059 \\
mci\_cn & mri\_clinical & mri & SPE & 30.2239 & 24.0114 & 36.1681 \\
mci\_cn & mri\_clinical & mri & balanced\_ACC & 19.5805 & 15.9627 & 23.0734 \\
mci\_cn & clinical & mri & ACC & 27.3954 & 24.0216 & 30.7692 \\
mci\_cn & clinical & mri & AUC & 27.3451 & 23.6158 & 31.1969 \\
mci\_cn & clinical & mri & SEN & 14.0198 & 9.88663 & 18.3782 \\
mci\_cn & clinical & mri & SPE & 46.5961 & 40.3837 & 52.5011 \\
mci\_cn & clinical & mri & balanced\_ACC & 30.308 & 26.6813 & 33.8772 \\
emci\_lmci & mri\_syn\_both & clinical & ACC & -33.2593 & -37.3973 & -28.8994 \\
emci\_lmci & mri\_syn\_both & clinical & AUC & -33.6068 & -38.8337 & -28.1938 \\
emci\_lmci & mri\_syn\_both & clinical & SEN & -55.6327 & -63.701 & -46.9656 \\
emci\_lmci & mri\_syn\_both & clinical & SPE & -21.5936 & -26.7982 & -17.0577 \\
emci\_lmci & mri\_syn\_both & clinical & balanced\_ACC & -38.6131 & -43.3044 & -33.6618 \\
emci\_lmci & mri\_syn\_both & mri\_clinical & ACC & -19.4993 & -24.0844 & -14.9047 \\
emci\_lmci & mri\_syn\_both & mri\_clinical & AUC & -23.2724 & -28.5706 & -17.9808 \\
emci\_lmci & mri\_syn\_both & mri\_clinical & SEN & -42.8707 & -51.0061 & -34.2177 \\
emci\_lmci & mri\_syn\_both & mri\_clinical & SPE & -7.35746 & -12.6171 & -2.13816 \\
emci\_lmci & mri\_syn\_both & mri\_clinical & balanced\_ACC & -25.1141 & -30.0744 & -20.0502 \\
emci\_lmci & mri\_clinical & mri & ACC & 20.4188 & 16.517 & 24.5458 \\
emci\_lmci & mri\_clinical & mri & AUC & 26.7938 & 21.7408 & 31.6749 \\
emci\_lmci & mri\_clinical & mri & SEN & 68.5034 & 60.4762 & 75.8912 \\
emci\_lmci & mri\_clinical & mri & SPE & -4.45906 & -8.98026 & 0.0840643 \\
emci\_lmci & mri\_clinical & mri & balanced\_ACC & 32.0222 & 27.4872 & 36.569 \\
emci\_lmci & clinical & mri & ACC & 34.1789 & 31.1872 & 37.1706 \\
emci\_lmci & clinical & mri & AUC & 37.1282 & 31.5136 & 42.5461 \\
emci\_lmci & clinical & mri & SEN & 81.2653 & 75.2245 & 86.6395 \\
emci\_lmci & clinical & mri & SPE & 9.77705 & 6.28655 & 13.6111 \\
emci\_lmci & clinical & mri & balanced\_ACC & 45.5212 & 41.9676 & 48.8719 \\
\end{longtable}
\noindent\parbox{\linewidth}{\footnotesize \DIFadd{12prespecified comparisons x5metrics. Delta target-reference; equal-weight3fold means.5000 paired fold/class bootstrap95\% CI, seed20261102. Marginal intervals without simultaneous coverage. Does not estimate repeated-training variation.}}
\endgroup
\end{landscape}

%DIF >  END supplementary/tables/classification_clinical_comparisons


%DIF >  BEGIN supplementary/tables/classification_delong
%DIF >  Draft only. Requires longtable, booktabs; wide tables also pdflscape.
%DIF >  Source: G5/existing_imaging_delong_holm.csv;G5/clinical_delong_by_fold.csv
\begin{landscape}
\begingroup
\footnotesize
\setlength{\tabcolsep}{4pt}
\begin{longtable}{llllllll}
\caption{{\color{blue}Paired DeLong comparisons by outer fold}}\label{tab:classification_delong}\\
\toprule
task & fold & target & reference & delta\_AUC\_pp & p\_raw & p\_Holm & family \\
\midrule
\endfirsthead
\toprule
task & fold & target & reference & delta\_AUC\_pp & p\_raw & p\_Holm & family \\
\midrule
\endhead
\bottomrule
\endfoot
ad\_cn & 1 & mri\_syn\_both & mri & 8.05704 & 0.00609343 & 0.109682 & existing18 \\
ad\_cn & 1 & mri\_syn\_both & mri\_real\_both & 0.730838 & 0.461309 & 1 & existing18 \\
ad\_cn & 2 & mri\_syn\_both & mri & 3.65419 & 0.113616 & 1 & existing18 \\
ad\_cn & 2 & mri\_syn\_both & mri\_real\_both & 0.855615 & 0.40185 & 1 & existing18 \\
ad\_cn & 3 & mri\_syn\_both & mri & 8.69487 & 0.0202858 & 0.344858 & existing18 \\
ad\_cn & 3 & mri\_syn\_both & mri\_real\_both & -3.24032 & 0.0689018 & 1 & existing18 \\
mci\_cn & 1 & mri\_syn\_both & mri & -5 & 0.205433 & 1 & existing18 \\
mci\_cn & 1 & mri\_syn\_both & mri\_real\_both & -3.23191 & 0.299826 & 1 & existing18 \\
mci\_cn & 2 & mri\_syn\_both & mri & -2.52874 & 0.425017 & 1 & existing18 \\
mci\_cn & 2 & mri\_syn\_both & mri\_real\_both & 0.0338066 & 0.990617 & 1 & existing18 \\
mci\_cn & 3 & mri\_syn\_both & mri & 1.13929 & 0.767688 & 1 & existing18 \\
mci\_cn & 3 & mri\_syn\_both & mri\_real\_both & 4.12315 & 0.203511 & 1 & existing18 \\
emci\_lmci & 1 & mri\_syn\_both & mri & 2 & 0.73837 & 1 & existing18 \\
emci\_lmci & 1 & mri\_syn\_both & mri\_real\_both & 4.8125 & 0.336152 & 1 & existing18 \\
emci\_lmci & 2 & mri\_syn\_both & mri & 6.10119 & 0.144769 & 1 & existing18 \\
emci\_lmci & 2 & mri\_syn\_both & mri\_real\_both & 1.70068 & 0.736135 & 1 & existing18 \\
emci\_lmci & 3 & mri\_syn\_both & mri & 2.46316 & 0.633676 & 1 & existing18 \\
emci\_lmci & 3 & mri\_syn\_both & mri\_real\_both & 3.36842 & 0.383536 & 1 & existing18 \\
ad\_cn & 1 & mri\_syn\_both & clinical & -4.63458 & 0.00165732 & 0.00864088 & new36 \\
ad\_cn & 1 & mri\_syn\_both & mri\_clinical & -2.7451 & 0.037946 & 0.112295 & new36 \\
ad\_cn & 1 & mri\_clinical & mri & 10.8021 & 3.40179e-05 & 0.000408215 & new36 \\
ad\_cn & 1 & clinical & mri & 12.6916 & 3.45292e-05 & 0.000408215 & new36 \\
ad\_cn & 2 & mri\_syn\_both & clinical & -4.02852 & 0.037572 & 0.112295 & new36 \\
ad\_cn & 2 & mri\_syn\_both & mri\_clinical & -4.0107 & 0.0374316 & 0.112295 & new36 \\
ad\_cn & 2 & mri\_clinical & mri & 7.66488 & 0.00123441 & 0.00864088 & new36 \\
ad\_cn & 2 & clinical & mri & 7.68271 & 0.00126115 & 0.00864088 & new36 \\
ad\_cn & 3 & mri\_syn\_both & clinical & -9.86499 & 0.000460643 & 0.00414579 & new36 \\
ad\_cn & 3 & mri\_syn\_both & mri\_clinical & -9.27093 & 0.000566367 & 0.00453094 & new36 \\
ad\_cn & 3 & mri\_clinical & mri & 17.9658 & 5.14555e-08 & 8.74744e-07 & new36 \\
ad\_cn & 3 & clinical & mri & 18.5599 & 1.01156e-07 & 1.61849e-06 & new36 \\
mci\_cn & 1 & mri\_syn\_both & clinical & -30.4801 & 6.99058e-18 & 2.4467e-16 & new36 \\
mci\_cn & 1 & mri\_syn\_both & mri\_clinical & -18.1136 & 1.54431e-07 & 2.31646e-06 & new36 \\
mci\_cn & 1 & mri\_clinical & mri & 13.1136 & 1.69917e-07 & 2.37884e-06 & new36 \\
mci\_cn & 1 & clinical & mri & 25.4801 & 3.22727e-15 & 9.68181e-14 & new36 \\
mci\_cn & 2 & mri\_syn\_both & clinical & -30.6694 & 4.38095e-19 & 1.57714e-17 & new36 \\
mci\_cn & 2 & mri\_syn\_both & mri\_clinical & -25.0575 & 2.29151e-14 & 6.18707e-13 & new36 \\
mci\_cn & 2 & mri\_clinical & mri & 22.5287 & 5.89174e-13 & 1.41402e-11 & new36 \\
mci\_cn & 2 & clinical & mri & 28.1406 & 1.51617e-16 & 4.85174e-15 & new36 \\
mci\_cn & 3 & mri\_syn\_both & clinical & -27.2752 & 2.00445e-14 & 5.61247e-13 & new36 \\
mci\_cn & 3 & mri\_syn\_both & mri\_clinical & -27.6617 & 3.99106e-15 & 1.15741e-13 & new36 \\
mci\_cn & 3 & mri\_clinical & mri & 28.801 & 1.27578e-16 & 4.33764e-15 & new36 \\
mci\_cn & 3 & clinical & mri & 28.4145 & 1.7613e-15 & 5.46004e-14 & new36 \\
emci\_lmci & 1 & mri\_syn\_both & clinical & -33.0625 & 1.99145e-12 & 3.78376e-11 & new36 \\
emci\_lmci & 1 & mri\_syn\_both & mri\_clinical & -15.7917 & 0.00198083 & 0.00864088 & new36 \\
emci\_lmci & 1 & mri\_clinical & mri & 17.7917 & 0.000226449 & 0.00226449 & new36 \\
emci\_lmci & 1 & clinical & mri & 35.0625 & 1.37494e-12 & 3.02487e-11 & new36 \\
emci\_lmci & 2 & mri\_syn\_both & clinical & -33.6947 & 8.53209e-13 & 1.96238e-11 & new36 \\
emci\_lmci & 2 & mri\_syn\_both & mri\_clinical & -20.0255 & 4.42156e-07 & 5.74803e-06 & new36 \\
emci\_lmci & 2 & mri\_clinical & mri & 26.1267 & 3.80927e-12 & 6.85669e-11 & new36 \\
emci\_lmci & 2 & clinical & mri & 39.7959 & 1.36072e-16 & 4.49038e-15 & new36 \\
emci\_lmci & 3 & mri\_syn\_both & clinical & -34.0632 & 1.42116e-12 & 3.02487e-11 & new36 \\
emci\_lmci & 3 & mri\_syn\_both & mri\_clinical & -34 & 1.62294e-12 & 3.24588e-11 & new36 \\
emci\_lmci & 3 & mri\_clinical & mri & 36.4632 & 4.30757e-14 & 1.11997e-12 & new36 \\
emci\_lmci & 3 & clinical & mri & 36.5263 & 4.98084e-14 & 1.24521e-12 & new36 \\
\end{longtable}
\noindent\parbox{\linewidth}{\footnotesize \DIFadd{Two-sided paired DeLong only within each outer test fold. Holm separately for original18 and new36 test families. No pooled/combined/mean p-value or three-fold overall significance is asserted. Undefined variance would remain NA and count conservatively in family size.}}
\endgroup
\end{landscape}

%DIF >  END supplementary/tables/classification_delong


%DIF >  BEGIN supplementary/tables/classification_roc
\begin{figure}[p]\centering
\includegraphics[width=0.32\textwidth]{revision_roc_ad_cn}
\includegraphics[width=0.32\textwidth]{revision_roc_mci_cn}
\includegraphics[width=0.32\textwidth]{revision_roc_emci_lmci}
\caption{{\color{blue}\textbf{Cross-sectional classification ROC curves for AD/CN, EMCI+LMCI/CN and EMCI/LMCI (left to right).} Curves average three fold-specific TPRs at a fixed FPR grid; shaded regions are pointwise conditional 95\% bootstrap bands, not simultaneous confidence bands. Legend AUC values are equally weighted empirical fold AUC means. Comparisons include MRI, MRI+Real both, MRI+Syn both, Clinical-only and MRI+Clinical. The clinical fields include diagnosis-related scores and historical group-conditioned fills; high clinical AUC is not evidence of independent synthetic-PET benefit.}}
\label{fig:classification_roc}
\end{figure}

%DIF >  END supplementary/tables/classification_roc


%DIF >  BEGIN supplementary/tables/e1_score_removal
\begin{table}[p]\centering\small
\caption{{\color{blue}\textbf{E1 fixed-weight score-removal sensitivity.} Differences are No-Scores minus Full, except the nonnegative output difference between predictions. This is a distinct E1 Full checkpoint family, not the reconstructed Full used in the generator-comparison tables.}}
\label{tab:e1_score_removal}
\begin{tabularx}{\textwidth}{llrXrrr}
\toprule
Tracer & Metric & N & Mean $\Delta$ & CI lower & CI upper & Holm $p$ \\
\midrule
FDG & ROI MAE & 1597 & 3.63611e-05 & 1.18272e-05 & 6.01943e-05 & 0.0035996 \\
FDG & Masked MAE & 1597 & -1.39166e-05 & -3.50725e-05 & 7.21746e-06 & 0.22038 \\
FDG & PSNR & 1597 & -0.00417362 & -0.00673816 & -0.00160625 & -- \\
FDG & SSIM (7) & 1597 & 0.000220733 & 0.000182378 & 0.000260921 & -- \\
FDG & SSIM (11) & 1597 & 0.000215856 & 0.000178967 & 0.000253801 & -- \\
FDG & Output difference & 1597 & 0.00101927 & 0.00100572 & 0.00103318 & -- \\
AV45 & ROI MAE & 1364 & 7.87648e-05 & 4.56388e-06 & 0.000152402 & 0.039996 \\
AV45 & Masked MAE & 1364 & 0.0001445 & 9.57843e-05 & 0.000193458 & 0.00039996 \\
AV45 & PSNR & 1364 & -0.0201903 & -0.0247279 & -0.0157227 & -- \\
AV45 & SSIM (7) & 1364 & 0.000684918 & 0.000619745 & 0.000750227 & -- \\
AV45 & SSIM (11) & 1364 & 0.000554805 & 0.000491236 & 0.000619548 & -- \\
AV45 & Output difference & 1364 & 0.00110469 & 0.0010868 & 0.00112321 & -- \\
\bottomrule
\end{tabularx}
\par\smallskip\footnotesize
\DIFaddFL{The MRI, tracer context and model weights are fixed; the eight-field patient token is replaced with a frozen Age+Sex+Weight token, not an empty token. No model is retrained. ROI/masked MAE use normalized-intensity units, PSNR is dB, and SSIM is unitless. CIs use 10,000 within-fold paired bootstrap resamples (seed 20261017). ROI and masked MAE are the primary endpoints: paired sign-permutation tests use 10,000 permutations and Holm correction across the two tracers separately within each endpoint. Other intervals are descriptive. The full-reference measurements were reused; whole-cohort Full predictions were not replayed anew. Nonzero output sensitivity and small error changes do not prove independent clinical value.
}\end{table}

%DIF >  END supplementary/tables/e1_score_removal


%DIF >  BEGIN supplementary/tables/demographic_stratification
%DIF >  R2-16 subgroup tables; generated from summary_metrics.csv.
%DIF >  Requires graphicx (resizebox). No manuscript file is modified.

\begin{table}[htbp]
\centering
\scriptsize
\caption{{\color{blue}Age- and sex-stratified classification performance: AD vs CN.}}
\label{tab:r2-16-ad-cn}
\resizebox{\textwidth}{!}{%
\begin{tabular}{llrrrcc}
\hline
Subgroup & Input & $N$ & Negative & Positive & AUC (\%) [95\% CI] & Balanced ACC (\%) [95\% CI] \\
\hline
Age $<75$ & MRI & 256 & 177 & 79 & 88.17 [83.22, 92.56] & 80.27 [74.84, 85.42] \\
 & MRI+Syn dual PET & 256 & 177 & 79 & 93.02 [88.60, 96.90] & 87.71 [83.01, 91.96] \\
 & MRI+Real dual PET & 256 & 177 & 79 & 93.69 [88.98, 97.46] & 90.10 [85.63, 93.98] \\
\hline
Age $\geq75$ & MRI & 214 & 128 & 86 & 84.99 [79.58, 89.90] & 76.42 [70.59, 82.07] \\
 & MRI+Syn dual PET & 214 & 128 & 86 & 93.83 [90.34, 96.73] & 86.33 [81.41, 90.75] \\
 & MRI+Real dual PET & 214 & 128 & 86 & 94.55 [91.36, 97.17] & 87.74 [83.01, 92.14] \\
\hline
Female & MRI & 226 & 155 & 71 & 84.85 [78.40, 90.29] & 77.05 [70.86, 82.84] \\
 & MRI+Syn dual PET & 226 & 155 & 71 & 94.55 [90.43, 97.73] & 88.02 [83.17, 92.52] \\
 & MRI+Real dual PET & 226 & 155 & 71 & 94.70 [90.12, 97.88] & 88.45 [83.25, 93.04] \\
\hline
Male & MRI & 244 & 150 & 94 & 87.54 [82.46, 91.80] & 78.59 [73.01, 83.57] \\
 & MRI+Syn dual PET & 244 & 150 & 94 & 93.07 [89.24, 96.32] & 86.89 [82.16, 91.03] \\
 & MRI+Real dual PET & 244 & 150 & 94 & 93.15 [89.20, 96.41] & 88.74 [84.24, 92.83] \\
\hline
\end{tabular}}
\par\smallskip
\begin{minipage}{\textwidth}\scriptsize
\DIFaddFL{Metrics are equal means of three within-fold subgroup estimates. Intervals are marginal 95\% percentile intervals from 5,000 paired resamples within fold, subgroup and true class, conditional on fixed models and folds; no multiplicity correction was applied. Age and sex are separate subgroup analyses. Positive classes are AD, EMCI+LMCI and LMCI for the respective tasks. Recorded age could not be aligned to the MRI visit because the recovered export lacks visit dates.
}\end{minipage}
\end{table}

\begin{table}[htbp]
\centering
\scriptsize
\caption{{\color{blue}Age- and sex-stratified classification performance: EMCI+LMCI vs CN.}}
\label{tab:r2-16-mci-cn}
\resizebox{\textwidth}{!}{%
\begin{tabular}{llrrrcc}
\hline
Subgroup & Input & $N$ & Negative & Positive & AUC (\%) [95\% CI] & Balanced ACC (\%) [95\% CI] \\
\hline
Age $<75$ & MRI & 456 & 177 & 279 & 71.56 [66.47, 76.49] & 66.28 [62.07, 70.67] \\
 & MRI+Syn dual PET & 456 & 177 & 279 & 67.96 [62.99, 72.82] & 60.84 [56.65, 65.07] \\
 & MRI+Real dual PET & 456 & 177 & 279 & 69.28 [64.47, 74.07] & 63.35 [58.88, 67.95] \\
\hline
Age $\geq75$ & MRI & 285 & 128 & 157 & 71.44 [65.36, 77.37] & 62.91 [57.85, 68.12] \\
 & MRI+Syn dual PET & 285 & 128 & 157 & 71.86 [65.66, 77.68] & 64.94 [60.10, 69.99] \\
 & MRI+Real dual PET & 285 & 128 & 157 & 70.82 [64.85, 76.86] & 64.11 [58.89, 69.50] \\
\hline
Female & MRI & 351 & 155 & 196 & 72.83 [67.25, 78.06] & 64.58 [59.83, 69.35] \\
 & MRI+Syn dual PET & 351 & 155 & 196 & 68.19 [62.27, 73.84] & 62.11 [57.49, 66.79] \\
 & MRI+Real dual PET & 351 & 155 & 196 & 69.47 [63.90, 74.92] & 61.80 [56.79, 66.74] \\
\hline
Male & MRI & 390 & 150 & 240 & 70.45 [65.00, 75.74] & 65.26 [60.83, 69.83] \\
 & MRI+Syn dual PET & 390 & 150 & 240 & 70.79 [65.50, 75.90] & 62.89 [58.40, 67.35] \\
 & MRI+Real dual PET & 390 & 150 & 240 & 68.85 [63.24, 74.17] & 64.64 [59.72, 69.43] \\
\hline
\end{tabular}}
\par\smallskip
\begin{minipage}{\textwidth}\scriptsize
\DIFaddFL{Metrics are equal means of three within-fold subgroup estimates. Intervals are marginal 95\% percentile intervals from 5,000 paired resamples within fold, subgroup and true class, conditional on fixed models and folds; no multiplicity correction was applied. Age and sex are separate subgroup analyses. Positive classes are AD, EMCI+LMCI and LMCI for the respective tasks. Recorded age could not be aligned to the MRI visit because the recovered export lacks visit dates.
}\end{minipage}
\end{table}

\begin{table}[htbp]
\centering
\scriptsize
\caption{{\color{blue}Age- and sex-stratified classification performance: EMCI vs LMCI.}}
\label{tab:r2-16-emci-lmci}
\resizebox{\textwidth}{!}{%
\begin{tabular}{llrrrcc}
\hline
Subgroup & Input & $N$ & Negative & Positive & AUC (\%) [95\% CI] & Balanced ACC (\%) [95\% CI] \\
\hline
Age $<75$ & MRI & 279 & 192 & 87 & 63.05 [56.22, 70.16] & 54.83 [50.58, 59.32] \\
 & MRI+Syn dual PET & 279 & 192 & 87 & 67.22 [60.31, 73.81] & 59.69 [54.02, 65.38] \\
 & MRI+Real dual PET & 279 & 192 & 87 & 60.90 [53.61, 67.96] & 58.22 [53.01, 63.18] \\
\hline
Age $\geq75$ & MRI & 157 & 95 & 62 & 60.34 [50.90, 69.53] & 50.57 [44.73, 56.32] \\
 & MRI+Syn dual PET & 157 & 95 & 62 & 63.68 [54.37, 72.48] & 60.21 [52.28, 67.63] \\
 & MRI+Real dual PET & 157 & 95 & 62 & 63.13 [54.27, 71.58] & 59.35 [51.65, 67.06] \\
\hline
Female & MRI & 196 & 124 & 72 & 62.55 [54.22, 70.81] & 55.06 [49.87, 60.73] \\
 & MRI+Syn dual PET & 196 & 124 & 72 & 67.23 [59.36, 74.89] & 61.64 [54.97, 68.46] \\
 & MRI+Real dual PET & 196 & 124 & 72 & 63.21 [55.17, 71.01] & 60.65 [54.32, 67.15] \\
\hline
Male & MRI & 240 & 163 & 77 & 64.67 [56.97, 72.05] & 52.86 [48.50, 57.37] \\
 & MRI+Syn dual PET & 240 & 163 & 77 & 65.84 [58.33, 73.24] & 58.65 [52.77, 64.92] \\
 & MRI+Real dual PET & 240 & 163 & 77 & 63.29 [55.47, 71.07] & 57.54 [51.54, 63.57] \\
\hline
\end{tabular}}
\par\smallskip
\begin{minipage}{\textwidth}\scriptsize
\DIFaddFL{Metrics are equal means of three within-fold subgroup estimates. Intervals are marginal 95\% percentile intervals from 5,000 paired resamples within fold, subgroup and true class, conditional on fixed models and folds; no multiplicity correction was applied. Age and sex are separate subgroup analyses. Positive classes are AD, EMCI+LMCI and LMCI for the respective tasks. Recorded age could not be aligned to the MRI visit because the recovered export lacks visit dates.
}\end{minipage}
\end{table}

%DIF >  END supplementary/tables/demographic_stratification
\DIFaddend 


% END supplementary/supplementary

\clearpage
\endgroup

\printbibliography

\end{document}
